\documentclass{article}
\usepackage{graphicx} % Required for inserting images

\usepackage{amsmath}
\usepackage{amssymb}
\usepackage{amsfonts}
\usepackage{algorithm}
\usepackage{algpseudocode}
\usepackage{dsfont} % Add this to your preamble
\usepackage{hyperref} % for hyperlinks

\usepackage[showdeletions]{color-edits}
%\usepackage[suppress]{color-edits}
\addauthor{il}{blue}
\addauthor{kc}{red}

\title{optimization section v3}
\author{kumarc}
\date{July 2026}

\begin{document}

\maketitle

\section{Main Idea}
We have a few interesting and novel parts of our paper. The majority of our methods are very comparable to existing methods in the political redistricting literature. The key difference between the school zoning and the political redistricting problem is that \begin{enumerate}
    \item the school district constraints are far more restrictive and even finding feasible solutions can be difficult, whereas for political redistricting it is more about proving optimality / generating a large number of solutions
    \item we are able to fix centroids without excessive precomputation since we care less about optimality and schools serve as logical centroid locations.
\end{enumerate} 
Multi-level, MCMC, and MIP strategies have already been applied in the political redistricting literature. We offer a few key insights:
\begin{enumerate}
    \item We apply political redistricting methods to school rezoning.
    \item We concretely (albeit almost exclusively empirically) show that CP is better than MIP at this problem, with strong implications for the political redistricting literature.
    \item We show that ReCom and relaxed ReCom can not effectively solve the school zoning problem (though short bursts and adaptive short bursts find significant success).
    \item We provide a concrete multi-level strategy that is highly effective under tight constraints and allows us to have improved results; other multi-level strategies in the literature are likely to result in infeasible solutions.
    \item We introduce a way to optimize for submodular choice set utility functions via logic-based Benders decomposition and show how it can be used as a heuristic for improving school choice outcomes.
    \item We formulate market-clearing equilibrium (MID) and Sample Average Approximation (SAA) cutting-plane models to directly optimize utilitarian welfare over student budget sets.
    \item We derive a congestion-priced upper bound on stable-matching welfare that is a function of the zoning rather than a constant, is valid for every non-negative price vector and every lottery seed, and --- because it separates across applicants --- admits an exact finite closed-form outer approximation requiring no recourse LP. This yields a cutting-plane method with finite convergence to a certified optimum of the bound, and prices the zone-confinement cost that the a priori welfare constants leave unpriced. Replicated over four seeds and two centroid configurations, it improves budget welfare over SAA by about $1\%$ while stopping on its own separation certificate in roughly a tenth of the wall time.
\end{enumerate}

Many of these ideas can (and arguably should be) used in the political redistricting literature.

\section{Base Solvers}
We represent a school district as an undirected, connected graph $G = (V, E)$, where each vertex $v \in V$ corresponds to a distinct geographic unit (e.g., a census block or block group) and each edge $e = \{u, v\} \in E$ indicates that two units are spatially adjacent. We seek to partition the graph into exactly $k$ zones, denoted  $$\mathcal{P} = \{Z_1, Z_2, \dots, Z_k\}$$ which must satisfy:
\begin{enumerate}
    \item $\bigcup_{i=1}^k Z_i = V$
    \item $Z_i \cap Z_j = \emptyset \qquad \forall i \neq j$
    \item $Z_i \neq \emptyset \qquad \forall i \in [k]$
\end{enumerate}

We define the structural requirements of a valid school zoning partition $\mathcal{P}$ based on $G$, centroids $C$, number of zones $k$, and balance parameters $\delta$.

\subsection{Boundary Compactness (Objective Function)}
To ensure that zones are geographically compact and pass the visual ``eye test'' for stakeholders, we evaluate the quality of a partition by its boundary size. The boundary of a partition $\mathcal{P}$, denoted $\partial \mathcal{P}$, is the set of all edges cut by the partition:
$$\partial \mathcal{P} = \left\{\{u,v\} \in E \mid u \in Z_i, v \in Z_j, i \neq j\right\}$$
Our fundamental objective is to find a partition that minimizes the total number of cut edges, $|\partial \mathcal{P}|$ (or, more generally, the total boundary length $\sum_{\{u,v\} \in \partial \mathcal{P}} w_{u,v}$ where $w_{u,v} \ge 0$ is the shared geographic boundary length).

\subsection{Constraints}
\subsubsection{Centroid and Anchor Alignment}
To reduce symmetry in our later formulations, each zone $Z_i$ is anchored by a designated centroid vertex $c_i \in V$, where $C = \{c_1, c_2, \ldots, c_k \}$. Schools serve as the natural geographic anchors for these centroids. A valid partition must structurally align each zone with exactly one unique centroid:
$$c_i \in Z_i \qquad \forall i \in [k]$$

\subsubsection{School Balance}
Let $S$ denote the total set of schools in the district. Also let $s(v) \in \mathbb{Z}_{\geq 0}$ denote the number of schools in a vertex $v\in V$. Note that this can potentially be greater than 1 for coarse graphs. Then we have $S(Z_i) = \sum_{v \in Z_i} s(v)$ for some zone $i \in [k]$. To prevent resource and choice imbalances, the total number of schools allocated to any given zone must be distributed as evenly as possible across the district:
$$\left\lfloor \frac{|S|}{k} \right\rfloor \leq S(Z_i) \leq \left\lceil \frac{|S|}{k} \right\rceil \qquad \forall Z_i \in \mathcal{P}$$

\subsubsection{Zone Contiguity}
A zone cannot be split into geographically isolated pockets. Topologically, this requires that the induced subgraph of each zone must form a single connected component:
$$\text{The induced subgraph } G[Z_i] \text{ is connected } \forall Z_i \in \mathcal{P}$$

\subsubsection{Capacity and Enrollment Tolerances}
Each vertex $v \in V$ carries an associated estimated student population $n_v \in \mathbb{R}_{\geq 0}$ and a school seat capacity $c_v \in \mathbb{Z}_{\geq 0}$ (where $c_v > 0$ only if $v \in S$). For any zone $Z_i$, we define its total student enrollment as $N(Z_i) = \sum_{v \in Z_i} n_v$ and its total operational capacity as $C(Z_i) = \sum_{v \in Z_i} c_v$.

To protect schools from severe over-enrollment or under-utilization, the ratio of total seats to total students within any zone must fall within specified lower shortage tolerances ($\delta^{-}$) and upper overage tolerances ($\delta^{+}$):
$$(1-\delta^{-}) \leq \frac{C(Z_i)}{N(Z_i)} \leq (1+\delta^{+}) \qquad \forall Z_i \in \mathcal{P}$$

\subsubsection{Demographic Diversity and Socioeconomic Balance}
Let $f_v \in \mathbb{R}_{\geq 0}$ denote the estimated count of students qualifying for Free and Reduced Lunch (FRL) programs at vertex $v$. The district-wide FRL proportion is a fixed constant calculated as:
$$F = \frac{\sum_{v \in V} f_v}{\sum_{v \in V} n_v}$$
To counteract socioeconomic isolation, the localized FRL proportion of any individual zone $Z_i$ must deviate from the global target $F$ by no more than an allowable variance threshold $\delta_F \in \mathbb{R}_{>0}$:
$$\left| \frac{\sum_{v \in Z_i} f_v}{N(Z_i)} - F \right| \leq \delta_F \qquad \forall Z_i \in \mathcal{P}$$

Let $\mathcal{I}$ represent the set of all valid input configurations for a problem with $(G, C, k, \delta)$ and $\mathcal{P}$ be the space of all possible partitions. Define the feasible set $\Omega(G, C, k, \delta) \subseteq \mathcal{P}$ as the subset of partitions that satisfy all constraints given the input $(G, C, k, \delta)$.

We define the function $B: \mathcal{I} \to \mathcal{P} \cup \{\perp\}$ to map input configurations to a partition or the bottom element (representing infeasibility) as follows:

\begin{equation}
B(G, C, k, \delta) = 
\begin{cases} 
\arg\min_{\mathcal{P} \in \Omega(G, C, k, \delta)} |\partial \mathcal{P}| & \text{if } \Omega(G, C, k, \delta) \neq \emptyset \\
\perp & \text{if } \Omega(G, C, k, \delta) = \emptyset 
\end{cases}
\end{equation}
Where $\perp$ denotes an infeasible input. Ties between solutions with identical cut edges are broken arbitrarily. All solvers will define various implementations of this base solver structure. 

\subsection{Math Programming}

To solve this problem we will employ constraint programming (CP), a technique related to Mixed Integer Programming (MIP). CP differs from MIP in that it relies on boolean logic, domain reduction, and clause learning to efficiently prune search spaces, excelling in problems where combinatorial constraints impose severe barriers to finding feasible solutions. As we demonstrate empirically, CP-SAT significantly outperforms MIP backends on this problem, largely due to the combinatorial difficulty of satisfying the contiguity and balance constraints simultaneously. Notably, CP-SAT requires integer coefficients, so we employ integer scaling factors to encode constraints that involve fractional values.

We formulate the base model using binary assignment variables $x_{v,z} \in \{0, 1\}$, indicating whether geographic vertex $v \in V$ is assigned to zone $z \in [k]$, and cut boundary variables $b_{i,j} \in \{0, 1\}$ for each edge $(i,j) \in E$:

\begin{align}
\min_{x, b} \quad & \sum_{(i,j)\in E} b_{i,j} \label{cp:obj} \\
\textrm{s.t.} \quad & \sum_{z \in [k]} x_{v,z} = 1 && \forall v \in V \label{cp:assignment} \\
& x_{c(z), z} = 1 && \forall z \in [k] \label{cp:centroid} \\
& x_{v,z} = 1 \implies \bigvee_{u \in \text{cn}(v, c(z))} x_{u,z} && \forall z \in [k], v \in V \setminus \{c(z)\} \label{cp:contiguity} \\
& \left\lfloor \frac{|S|}{k} \right\rfloor \leq \sum_{v \in V} s(v) x_{v,z} \leq \left\lceil \frac{|S|}{k} \right\rceil && \forall z \in [k] \label{cp:schools} \\
& (1-\delta^{-}) N(Z_z) \leq C(Z_z) \leq (1+\delta^{+}) N(Z_z) && \forall z \in [k] \label{cp:capacity} \\
& (F-\delta_F) N(Z_z) \leq \sum_{v \in V} f_v x_{v,z} \leq (F+\delta_F) N(Z_z) && \forall z \in [k] \label{cp:frl} \\
& (x_{i,z} \land \neg x_{j,z}) \lor (x_{j,z} \land \neg x_{i,z}) \implies b_{i,j} = 1 && \forall (i,j) \in E, z \in [k] \label{cp:cut} \\
& x_{v,z} \in \{0, 1\} && \forall v \in V, z \in [k] \label{cp:x_domain} \\
& b_{i,j} \in \{0, 1\} && \forall (i,j) \in E \label{cp:b_domain}
\end{align}

Each inequality directly realizes one of the structural requirements introduced in Section 2.2:
\begin{itemize}
    \item \textbf{Objective \eqref{cp:obj}:} Minimizes the total number of cut edges across all zone boundaries.
    \item \textbf{Unique Assignment \eqref{cp:assignment}:} Enforces that every geographic unit $v \in V$ is assigned to exactly one zone $z \in [k]$.
    \item \textbf{Centroid Anchoring \eqref{cp:centroid}:} Fixes each designated zone centroid $c(z)$ to zone $z$, eliminating zone permutation symmetries.
    \item \textbf{Zone Contiguity \eqref{cp:contiguity}:} Enforces graph connectedness through a rooted closer-neighbor condition based on the formulation of Shirabe (2005) and Validi et al. (2022). For each vertex $v \neq c(z)$ and zone $z$, let $\text{cn}(v, c(z))$ denote the set of adjacent neighbors that are strictly closer to centroid $c(z)$ in spatial geometry:
    $$\text{cn}(v, c(z)) = \{u \in N(v) \mid \text{dist}(u, c(z)) < \text{dist}(v, c(z))\}$$
    Constraint \eqref{cp:contiguity} stipulates that if vertex $v$ joins zone $z$, at least one neighbor $u \in \text{cn}(v, c(z))$ must also join zone $z$. Because distances to $c(z)$ strictly decrease along this chain, this constraint induces a cycle-free directed path from every assigned vertex back to the centroid $c(z)$, ensuring that the induced subgraph $G[Z_z]$ is connected.
    \item \textbf{School Count Balance \eqref{cp:schools}:} Distributes the $|S|$ schools evenly among all $k$ zones, restricting zone $z$'s school count between $\lfloor |S|/k \rfloor$ and $\lceil |S|/k \rceil$.
    \item \textbf{Capacity Tolerances \eqref{cp:capacity}:} Ensures that the aggregate seat capacity within zone $z$ remains within lower shortage tolerance $\delta^-$ and upper overage tolerance $\delta^+$ of the resident student population.
    \item \textbf{Socioeconomic Balance \eqref{cp:frl}:} Bounds zone $z$'s Free and Reduced Lunch (FRL) proportion within a maximum deviation $\delta_F$ from the district-wide mean $F$. (The corresponding multi-group racial and ethnic diversity balance constraints follow an analogous form and are detailed in the Appendix).
    \item \textbf{Cut Edge Implication \eqref{cp:cut}:} Enforces that $b_{i,j} = 1$ via Boolean half-reification whenever endpoints $i$ and $j$ of edge $(i,j) \in E$ are assigned to different zones (i.e., if vertex $i$ joins zone $z$ but $j$ does not, or vice versa).
\end{itemize}

Equivalently, CP-SAT receives these constraints in zero-right-hand-side form. Because CP-SAT requires integer coefficients, each inequality is scaled by a large integer multiplier $M$ (typically $M = 10^4$) and its coefficients are rounded conservatively before being added to the model. For example, the lower FRL constraint is written as:
\begin{equation*}
\sum_{v\in V}
\lfloor M \cdot \left(f_v-(F-\delta_F)n_v\right) \rfloor \cdot x_{v,z}
\geq 0 \qquad \forall z \in [k]
\end{equation*}
where the floor function preserves conservative bounding without weakening the feasibility guarantees.

\subsection{MCMC Models}
In the political redistricting literature, Markov Chain Monte Carlo (MCMC) methods—most prominently Recombination (ReCom)—are widely used to sample vast ensembles of contiguous districts. In political redistricting, the primary constraint is typically equal population balance (a single scalar target). In contrast, school rezoning introduces a substantially more restrictive constraint manifold defined by capacity bounds, tight demographic diversity bounds, and exact discrete school count allocations.

\subsubsection{ReCom Proposal Kernel}
All ReCom-family variants share a common combinatorial transition kernel operating on partition $\mathcal{P}_t = \{Z_1, \dots, Z_k\}$ of graph $G = (V, E)$: \kccomment{i dont like this section}
\begin{enumerate}
    \item \textbf{Adjacent Zone Pair Selection:} Let $\mathcal{A}(\mathcal{P}_t)$ denote the set of pairs of zones sharing at least one boundary edge:
    \begin{equation}
    \mathcal{A}(\mathcal{P}_t) = \left\{(a, b) \mid 1 \le a < b \le k, \; \exists (u, v) \in E \text{ with } u \in Z_a, v \in Z_b\right\}
    \end{equation}
    A pair $(a, b)$ is drawn uniformly at random from $\mathcal{A}(\mathcal{P}_t)$. Adaptive Short Bursts replaces this uniform draw with the violation-weighted draw of Section~\ref{subsubsec:adaptive-short-bursts}; every other variant uses it as stated.
    \item \textbf{Uniform Spanning Tree Sampling:} Form the merged region $U = Z_a \cup Z_b$. Because $Z_a$ and $Z_b$ are individually contiguous and share boundary edges, the induced subgraph $G[U]$ is connected. A spanning tree $T = (U, E(T))$ is sampled uniformly from the set $\mathcal{T}(G[U])$ of all spanning trees of $G[U]$ via Wilson's loop-erased random walk algorithm (Wilson, 1996), ensuring that each tree is drawn with exact probability $1 / |\mathcal{T}(G[U])|$.
    \item \textbf{Spanning Tree Cuts and Zone Contiguity:} Tree $T$ is rooted at an arbitrary root vertex $r \in U$. Each non-root vertex $u \in U \setminus \{r\}$ and its directed edge to its parent $(u, \text{parent}(u)) \in E(T)$ uniquely defines a tree bisection into subtree $T_u$ and complement $U \setminus T_u$. Because $T$ spans $G[U]$, the induced subgraphs $G[T_u]$ and $G[U \setminus T_u]$ are connected by construction, guaranteeing topological contiguity without post-hoc connectivity checks.
    \item \textbf{Candidate Reassignments:} Each subtree cut yields two candidate assignments:
    \begin{align}
    \text{Orientation 1:} \quad Z_a' &= T_u, & Z_b' &= U \setminus T_u \\
    \text{Orientation 2:} \quad Z_a' &= U \setminus T_u, & Z_b' &= T_u
    \end{align}
    Orientations that violate designated centroid anchoring ($c_a \notin Z_a'$ or $c_b \notin Z_b'$) or explicit vertex candidate restrictions are pruned, producing a legal candidate cut pool $\mathcal{C}$.
    \item \textbf{Fast Additive Evaluation:} Vertex attributes (students $n_v$, capacity $c_v$, schools $s(v)$, and demographic counts $f_v$) are additive set functions. Subtree statistics for each candidate cut $c \in \mathcal{C}$ are evaluated in a single post-order traversal over $T$, e.g., $N(T_u) = n_u + \sum_{w \in \text{children}(u)} N(T_w)$ and $N(U \setminus T_u) = N(U) - N(T_u)$. The updated boundary cost is computed from non-tree edge crossings via lowest common ancestor prefix sums:
    \begin{equation}
    |\partial \mathcal{P}_c| = |\partial \mathcal{P}_t| - |\partial(Z_a, Z_b)| + |\partial(Z_a^{(c)}, Z_b^{(c)})|
    \end{equation}
\end{enumerate}

\subsubsection{Standard ReCom}
Standard ReCom (DeFord et al., 2021) evaluates candidate proposals under a hard feasibility filter:
\begin{enumerate}
    \item \textbf{Initial Seeding:} Because no valid baseline map exists for school rezoning, the chain is initialized via a Voronoi diagram anchored at designated school centroids $C = \{c_1, \dots, c_k\}$, where each vertex $v$ is assigned to its geographically nearest centroid: $z^{(0)}(v) = \arg\min_{z \in [k]} \text{dist}(v, c_z)$. Disconnected components are reallocated to adjacent zones maintaining centroid paths, producing an initial contiguous partition $\mathcal{P}_0$. In general, $\mathcal{P}_0 \notin \Omega(G, C, k, \delta)$.
    \item \textbf{Proposal and Rejection:} Let $\mathcal{C}_{\text{feas}} = \{c \in \mathcal{C} \mid \mathcal{P}_c \in \Omega(G, C, k, \delta)\}$ be the set of candidate cuts where both prospective zones satisfy all capacity, school count, and demographic constraints. If $\mathcal{C}_{\text{feas}} \neq \emptyset$, a candidate cut is drawn uniformly: $c \sim \text{Uniform}(\mathcal{C}_{\text{feas}})$; otherwise, $c \sim \text{Uniform}(\mathcal{C})$. The transition is accepted if and only if the prospective partition is globally feasible:
    \begin{equation}
    \mathcal{P}_{t+1} = \begin{cases}
    \mathcal{P}_c & \text{if } \mathcal{P}_c \in \Omega(G, C, k, \delta) \\
    \mathcal{P}_t & \text{if } \mathcal{P}_c \notin \Omega(G, C, k, \delta)
    \end{cases}
    \end{equation}
\end{enumerate}
\textbf{Theoretical Bottleneck on School Rezoning:} In political redistricting, baseline legislative plans start within the feasible manifold $\Omega$, and the only primary constraint is equal population. In school rezoning, $\Omega$ is the simultaneous intersection of $m \ge 6k$ discrete constraints (exact school counts $\lfloor |S|/k \rfloor \le S(Z) \le \lceil |S|/k \rceil$, tight capacity bounds, and FRL balance). Because constraints are enforced post-cut, the probability that a uniform spanning tree cut simultaneously satisfies all bounds is vanishingly small: $P(c \in \mathcal{C}_{\text{feas}}) \approx 0$. As a result, standard ReCom rejects virtually all moves and remains trapped at $\mathcal{P}_0$.

\subsubsection{Relaxed ReCom}
To navigate through infeasible space and bypass zero acceptance, Relaxed ReCom eliminates the hard rejection step ($P_{\text{accept}} = 1$) and samples candidate cuts from a biased categorical distribution that steers the chain toward feasibility and geometric compactness.

For each candidate cut $c \in \mathcal{C}$ producing prospective zones $Z_a^{(c)}$ and $Z_b^{(c)}$, we compute the unnormalized scoring function $\mu(c)$:
\begin{equation}
\ln \mu(c) = \theta_{\text{trees}} \left( \text{cr}(Z_a^{(c)}) + \text{cr}(Z_b^{(c)}) \right) + \sum_{\kappa \in \mathcal{K}} \theta_\kappa \left( \ln \kappa(Z_a^{(c)}) + \ln \kappa(Z_b^{(c)}) \right)
\end{equation}
where:
\begin{itemize}
    \item $\text{cr}(Z) = \max(0, |E[Z]| - |V(Z)| + 1)$ is the cycle rank (cyclomatic number) of the induced subgraph $G[Z]$. By Kirchhoff's Matrix Tree Theorem, the number of spanning trees $|\mathcal{T}(G[Z])|$ scales exponentially with the cycle rank: $|\mathcal{T}(G[Z])| \approx \exp(\text{cr}(Z))$. Weighting by $\theta_{\text{trees}} = 1$ penalizes stringy, bottlenecked boundaries and favors internally dense, compact subgraphs.
    \item $\mathcal{K}$ is the set of balance metrics evaluated on each prospective zone:
    \begin{enumerate}
        \item Node count: $\kappa_{\text{nodes}}(Z) = |V(Z)|$, with weight $\theta_{\text{nodes}} = 1$
        \item Student enrollment: $\kappa_{\text{students}}(Z) = N(Z)$, with weight $\theta_{\text{students}} = 1$
        \item Seat capacity: $\kappa_{\text{seats}}(Z) = C(Z)$, with weight $\theta_{\text{seats}} = 1$
        \item Socioeconomic need: $\kappa_{\text{frl}}(Z) = F(Z)$, with weight $\theta_{\text{frl}} = 3$
        \item School count: $\kappa_{\text{sch}}(Z) = S(Z)$, with large weight $\theta_{\text{sch}} = 45$ to heavily penalize school count disparities
        \item Relative capacity mismatch: $\kappa_{\text{shortage\%}}(Z) = \max\left( \frac{|N(Z) - C(Z)|}{\max(N(Z), \epsilon)}, \epsilon \right)$ with weight $\theta_{\text{shortage\%}} = 10$ and $\epsilon = 10^{-12}$.
    \end{enumerate}
\end{itemize}
Candidate cut $c$ is drawn from the categorical distribution:
\begin{equation}
P(c) = \frac{\mu(c)}{\sum_{c' \in \mathcal{C}} \mu(c')}
\end{equation}
The selected move is accepted unconditionally ($\mathcal{P}_{t+1} \leftarrow \mathcal{P}_c$). While Relaxed ReCom traverses infeasible states freely, unconstrained random walks with static weights drift over long trajectories and frequently fail to settle into the narrow intersection of all feasibility constraints. These weights and features were selected via manual tuning. 

\subsubsection{Short Bursts}
Building upon the Short Bursts framework (Cannon et al., 2020), we replace the unbounded random walk with repeated localized search bursts of fixed length $L$ (typically $L = 25$). Starting from an incumbent partition $\mathcal{P}^{(b)}$ at burst iteration $b$, the chain executes $L$ unconstrained transitions using the ReCom or Relaxed ReCom proposal kernel, generating a trajectory of visited states $\mathcal{S}^{(b)} = \{\mathcal{P}_1, \dots, \mathcal{P}_L\}$. If ReCom is used, infeasible states should not be rejected, similar to Relaxed ReCom.

To compare intermediate partitions across heterogeneous constraints with disparate numerical scales (e.g., student counts $\sim 10^3$, school counts $\sim 10^0$, demographic proportions $\sim 10^{-1}$), we track the running maximum observed violation for each constraint $j \in [m]$ across all visited states up to step $\tau$:
\begin{equation}
M_j^{(\tau)} = \max_{0 \le t \le \tau} v_j(\mathcal{P}_t)
\end{equation}
where $v_j(\mathcal{P}) = \sum_{z=1}^k v_j(Z_z) \ge 0$ is the total linear violation slack for constraint $j$. The dynamically normalized violation penalty is defined as:
\begin{equation}
\Phi(\mathcal{P}) = \sum_{j: M_j^{(\tau)} > 0} \frac{v_j(\mathcal{P})}{M_j^{(\tau)}}
\end{equation}
Partitions are evaluated under a strict lexicographic Pareto ordering $\prec$ that prioritizes feasibility before optimizing boundary compactness:
\begin{equation}
\mathcal{P} \prec \mathcal{Q} \iff \begin{cases}
\mathcal{P} \in \Omega \quad \text{and} \quad \mathcal{Q} \notin \Omega \\
\mathcal{P}, \mathcal{Q} \in \Omega \quad \text{and} \quad |\partial \mathcal{P}| < |\partial \mathcal{Q}| \\
\mathcal{P}, \mathcal{Q} \notin \Omega \quad \text{and} \quad \Phi(\mathcal{P}) < \Phi(\mathcal{Q}) - \epsilon_{\text{tol}}
\end{cases}
\end{equation}
with tolerance $\epsilon_{\text{tol}} = 10^{-6}$. At the conclusion of burst $b$, the best visited state is identified:
\begin{equation}
\mathcal{P}^* = \arg\min_{\mathcal{P} \in \mathcal{S}^{(b)}} \mathcal{P} \quad \text{under } \prec
\end{equation}
If $\mathcal{P}^* \prec \mathcal{P}^{(b)}$, the improvement is accepted as the new incumbent: $\mathcal{P}^{(b+1)} \leftarrow \mathcal{P}^*$; otherwise, the chain resets to the incumbent: $\mathcal{P}^{(b+1)} \leftarrow \mathcal{P}^{(b)}$. This ratcheting mechanism allows the chain to explore infeasible configurations locally while guaranteeing that the search never drifts permanently into worse regions.

\subsubsection{Adaptive Short Bursts}\label{subsubsec:adaptive-short-bursts}
Although Short Bursts guides the search restarts, proposal generation within each burst remains unguided. When certain constraints (such as school count balance) are far more rigid than others, random tree bisections rarely decrease multiple violations simultaneously. We address this by formulating Adaptive Short Bursts, which dynamically couples an augmented Lagrangian proposal distribution with Adam dual multiplier ascent. The same Lagrangian energy guides both halves of the proposal: which pair of zones is merged, and which cut of the merged region is taken.

\paragraph{Percentage-Normalized Constraint Violations:} Constraint residuals are converted into fractional percentage deviations:
\begin{equation}
\widetilde{v}_{\text{cap}}(Z) = \frac{v_{\text{cap}}(Z)}{N(Z)} \times 100\%, \quad \widetilde{v}_{\text{sch}}(Z) = \frac{v_{\text{sch}}(Z)}{\bar{S}} \times 100\%
\end{equation}
where $\bar{S} = |S|/k$ is the target average school count per zone. This is done so that residuals are usually \kccomment{How do i say this better} greater than 1 unless satisfied so that squaring still has an effect on the district. \kccomment{this also relies on saying that we can decompose constraints with linear penalties }


\paragraph{Augmented Lagrangian Energy Formulation:}\kccomment{most ai naming of all time}

Let $\boldsymbol{\eta} = (\eta_1, \dots, \eta_m) \in \mathbb{R}_+^m$ denote dual penalty multipliers, initialized to $\eta_j^{(0)} = 1.0$. For each candidate cut $c \in \mathcal{C}$ yielding prospective partition $\mathcal{P}_c$, we define the candidate energy function:
\begin{equation}
E(c; \mathbf{w}) = |\partial \mathcal{P}_c| + \sum_{j=1}^m w_j \cdot \widetilde{v}_j(\mathcal{P}_c)^2
\end{equation}
where penalty weights $w_j$ are dynamically coupled to the incumbent partition's cut boundary cost:
\begin{equation}
w_j = \max\left(1.0, \; |\partial \mathcal{P}^{(b)}| + \eta_j \cdot \widetilde{v}_j(\mathcal{P}^{(b)})^2\right)
\end{equation}

\paragraph{Violation-Weighted Zone Pair Selection:} The energy $E(c; \mathbf{w})$ scores candidate cuts only after a zone pair has been merged, while the pair itself is drawn uniformly (Step 1 of the shared kernel). A burst therefore spends most of its proposals rebisecting zones that are already balanced, and the rigid constraints it needs to repair are precisely the ones concentrated in a few zones. We localize the same energy to a single zone. Writing $\partial(Z_a) = \bigcup_{b \neq a} \partial(Z_a, Z_b)$ for the cut edges incident to $Z_a$, define
\begin{equation}
E_{\mathrm{zone}}(Z_a; \mathbf{w}) = |\partial(Z_a)| + \sum_{j=1}^m w_j \cdot \widetilde{v}_j(Z_a)^2 ,
\end{equation}
which restricts both terms of $E(\cdot; \mathbf{w})$ to $Z_a$: the share of the objective that zone pays in boundary edges, and the violations for which it is individually responsible. The first zone is drawn from the Boltzmann distribution over zones with at least one neighbor,
\begin{equation}
P(a) = \frac{\exp\left( + \frac{E_{\mathrm{zone}}(Z_a; \mathbf{w})}{T} \right)}{\sum_{a'} \exp\left( + \frac{E_{\mathrm{zone}}(Z_{a'}; \mathbf{w})}{T} \right)} ,
\end{equation}
and its partner uniformly from its neighbors, $P(b \mid a) = 1/|\mathcal{N}(a)|$ with $\mathcal{N}(a) = \{b : (a, b) \in \mathcal{A}(\mathcal{P}_t)\}$. The resulting pair $(a, b)$ enters Step 2 of the kernel unchanged, and the same dual-coupled weights $\mathbf{w}$ and temperature $T$ serve both draws.

The sign of the exponent is reversed relative to cut selection below: cut selection suppresses high-energy candidates, whereas zone selection seeks them, so each burst is spent on the zone currently furthest from feasibility. Summing $E_{\mathrm{zone}}$ over zones double counts every cut pair, but a common scale factor is absorbed by $T$ and does not affect the draw. Greedy selection does not fixate, because per-zone violations are recomputed after every accepted transition: repairing the selected zone promotes a different zone to the top of the distribution. At $T = 1.0$ the weights $w_j$ grow large enough that this draw is effectively $\arg\max$, and on block\_2 at $k = 6$ the selected zone still changed on $58\%$ of consecutive proposals, with all six zones drawn. Biasing the pair draw costs no sampling guarantee, since Relaxed ReCom has already set $P_{\text{accept}} = 1$ and abandoned detailed balance; these chains are optimizers, not samplers.

\kccomment{Empirics for a results table, if we want them: at $120$s with \texttt{stop\_on\_feasible}, violation-weighted pair selection reaches feasibility $4$--$12\times$ faster than the uniform draw at $k = 6, 8$, and $4/5$ vs $1/5$ of seeds at $k = 10$ on both block\_2 and Block\_1. Standardizing $E_{\mathrm{zone}}$ across zones before the softmax makes $T$ scale-free but did not change feasibility rates ($32/60$ either way), so it is not used.}

\paragraph{Softmax Proposal Distribution:} Within each burst, candidate cuts are sampled from a Boltzmann distribution parameterized by temperature $T > 0$ (default $T = 1.0$):
\begin{equation}
P(c) = \frac{\exp\left( - \frac{E(c; \mathbf{w})}{T} \right)}{\sum_{c' \in \mathcal{C}} \exp\left( - \frac{E(c'; \mathbf{w})}{T} \right)}
\end{equation}
This exponentially suppresses candidate cuts that worsen heavily penalized violations, actively steering spanning tree bisections toward the feasible boundary.

\paragraph{Adam Dual Ascent Across Bursts:} At the end of burst $b$, after selecting the best partition $\mathcal{P}^*$ via the lexicographic ordering $\prec$, the dual ascent subgradient is computed from the squared violations of $\mathcal{P}^*$:
\begin{equation}
g_j^{(b)} = \widetilde{v}_j(\mathcal{P}^*)^2 \qquad \forall j \in [m]
\end{equation}
The dual multipliers $\boldsymbol{\eta}$ are updated using projected Adam optimization: \kccomment{Do i need to explain what adam is?}
\begin{align}
\mathbf{m}_b &= \beta_1 \mathbf{m}_{b-1} + (1 - \beta_1) \mathbf{g}^{(b)} \\
\mathbf{u}_b &= \beta_2 \mathbf{u}_{b-1} + (1 - \beta_2) (\mathbf{g}^{(b)} \odot \mathbf{g}^{(b)}) \\
\widehat{\mathbf{m}}_b &= \frac{\mathbf{m}_b}{1 - \beta_1^b}, \qquad \widehat{\mathbf{u}}_b = \frac{\mathbf{u}_b}{1 - \beta_2^b} \\
\eta_j^{(b+1)} &= \max\left(0.0, \; \eta_j^{(b)} + \gamma \frac{\widehat{m}_{j, b}}{\sqrt{\widehat{u}_{j, b}} + \epsilon_{\text{Adam}}}\right)
\end{align}
with learning rate $\gamma = 0.1$, moment decay rates $\beta_1 = 0.9, \beta_2 = 0.999$, and $\epsilon_{\text{Adam}} = 10^{-8}$. The updated multipliers $\boldsymbol{\eta}^{(b+1)}$ and incumbent $\mathcal{P}^{(b+1)}$ then determine the proposal weights $\mathbf{w}$ for the subsequent burst. \kccomment{Do OR people need any justification for why im using adam? Because neither I, nor anyone else on earth, knows why adam works as well as it does}

\paragraph{Closed-Loop Feedback Dynamics:} If a constraint $j$ remains persistently violated across bursts ($\widetilde{v}_j > 0$), its gradient $g_j > 0$ inflates $\eta_j$ and $w_j$, exponentially suppressing candidate cuts that fail to satisfy constraint $j$. When constraint $j$ is satisfied ($\widetilde{v}_j = 0$), $g_j = 0$, allowing the second moment to decay and stabilizing $\eta_j$. Once all constraints reach feasibility ($\widetilde{\mathbf{v}} = \mathbf{0}$), the energy collapses to $E(c; \mathbf{w}) = |\partial \mathcal{P}_c|$, and the Boltzmann distribution purely drives cut boundary minimization. Zone selection collapses likewise to $E_{\mathrm{zone}}(Z_a; \mathbf{w}) = |\partial(Z_a)|$, which targets the zone holding the most boundary edges, and hence the one with the most to gain from rebisection.

\section{Hierarchical Graph Contraction}

The computational complexity of the zoning problem scales rapidly with the number of vertices $|V|$ and zones $k$. Informally, the discrete search space contains $k^{|V|}$ possible assignments. Solving this directly at the census block level (thousands of nodes) is often computationally intractable. To resolve this, we develop a principled graph contraction hierarchy that creates arbitrary-fidelity coarsened graphs.

Let our base graph be $G_0 = (V_0, E_0)$. We construct a contracted graph $G_1 = (V_1, E_1)$ with $|V_1| = k_1 < |V_0|$ vertices that preserves the topological and spatial structure of $G_0$. Formally, $G_1$ is a valid contraction of $G_0$ under a surjective mapping $\phi_1: V_0 \to V_1$ if:
\begin{description}
    \item[Connectivity:] For every vertex $u \in V_1$, the induced preimage subgraph $G_0[\phi_1^{-1}(u)]$ is connected.
    \item[Adjacency:] For any two distinct vertices $u, v \in V_1$, an edge exists in $G_1$ if and only if their preimages are adjacent in $G_0$:
    \[
    (u, v) \in E_1 \iff \exists x \in \phi_1^{-1}(u), y \in \phi_1^{-1}(v) \text{ such that } (x, y) \in E_0
    \]
\end{description}

We compute this partitioning using KaHIP (Sanders and Schulz, 2013). Crucially, we do not contract school vertices. Let $S \subset V_0$ denote the set of vertices containing schools. When constructing any coarse graph $G_{\ell+1}$ from $G_\ell$, every school node $s \in S$ is held out as an indivisible singleton vertex ($V_{\ell+1}^{\text{school}} = S$). KaHIP partitions only the remaining non-school geographic units $V_\ell \setminus S$ into $k_{\ell+1} - |S|$ contiguous blocks. Coarse edges between aggregated residential blocks and school singletons are derived from crossing base edges: $(u, s) \in E_{\ell+1} \iff \exists x \in \phi^{-1}(u) \text{ such that } (x, s) \in E_\ell$.

Preserving school nodes as singletons is essential for three reasons:
\begin{enumerate}
    \item \textbf{School count exactness:} It ensures $s(v) = 1$ for each school node and $s(v) = 0$ for all non-school aggregates, enabling exact enforcement of the school balance constraint $\lfloor |S|/k \rfloor \le S(Z_i) \le \lceil |S|/k \rceil$ at any graph hierarchy level. If schools were merged into aggregate blocks, multiple schools could be bundled into a single coarse vertex, potentially rendering the school balance constraint infeasible.
    \item \textbf{Centroid anchor preservation:} Distinct school singletons ensure that designated zone centroids $c_i \in S$ remain individually assignable and are never co-contracted into the same coarse vertex.
    \item \textbf{Capacity localization:} School seat capacities $c_s$ remain precisely associated with their geographic points, preserving spatial accuracy in capacity-to-student balance.
\end{enumerate}

\section{Strategies}
\subsection{Multilevel Strategy}
While coarsened graphs can be solved independently, their primary utility lies in a multi-level optimization strategy. Specifically, we solve the zoning problem hierarchically across a sequence of levels $G_n, G_{n-1}, \dots, G_0$, from coarsest to finest. 

At each transition from level $i$ to level $i-1$, the coarser partition $\mathcal{P}_i$ is projected onto the finer graph $G_{i-1}$ using the contraction mapping $\phi_i$. Vertices situated deeply within the interior of a projected zone—whose graph distance to the zone boundary exceeds a specified radius $r$—have their assignments fixed to that zone. Only vertices near the boundary (within distance $r$) are treated as candidate decision variables in the finer optimization solve. The projected partition also provides a high-quality warm-start hint. This restricts the decision space at fine levels exclusively to boundary refinement, yielding substantial computational speedups.

\begin{algorithm}
\caption{Multilevel Strategy}
\begin{algorithmic}[1]
\Require Sequence of graphs $G_n, G_{n-1}, \dots, G_0$ from coarsest to finest, base solver $B$, parameters $\delta$, boundary relaxation radius $r$.
\Procedure{Multilevel}{$\{G_i\}_{i=n}^0, B, \delta, r$}
    \State $\mathcal{P}_n \gets B(G_n, C, k, \delta)$ \Comment{Solve coarsest level globally}
    \For{$i \leftarrow n-1$ down to $0$}
        \State $\tilde{\mathcal{P}}_i \gets \text{Project}(\mathcal{P}_{i+1}, G_i)$ \Comment{Project partition via contraction map}
        \State $V^{\text{bound}} \gets \{u \in V_i \mid \exists v \in N(u) \text{ s.t. } \tilde{\mathcal{P}}_i(u) \neq \tilde{\mathcal{P}}_i(v)\}$
        \State $\mathcal{B}_r \gets \{v \in V_i \mid \text{dist}_{G_i}(v, V^{\text{bound}}) \le r\}$ \Comment{Boundary neighborhood}
        \State Set candidate zones $\mathcal{C}(v) \gets \{\tilde{\mathcal{P}}_i(v)\}$ for all $v \in V_i \setminus \mathcal{B}_r$ \Comment{Fix interior}
        \State Set candidate zones $\mathcal{C}(v) \gets [k]$ for all $v \in \mathcal{B}_r$ \Comment{Relax boundary}
        \State $\mathcal{P}_i \gets B(G_i, C, k, \delta \mid \text{hint}=\tilde{\mathcal{P}}_i, \text{candidates}=\mathcal{C})$
    \EndFor
    \State \Return $\mathcal{P}_0$
\EndProcedure
\end{algorithmic}
\end{algorithm}

This approach yields two major practical advantages:
\begin{enumerate}
    \item We find higher-quality, more compact solutions within the same time limit compared to solving fine graphs from scratch.
    \item We reliably discover feasible solutions at the census block level in parameter regimes where monolithic fine-level solvers fail completely.
\end{enumerate}

\section{Choice Set Optimization}

\kccomment{This should maybe go to the appendix}
Instead of solely evaluating zoning by geometric compactness, we now consider the academic welfare that a zone configuration delivers to families. Under a zoned choice system, a family residing in vertex $v$ gains access to the set of schools assigned to that zone, $S' \subseteq S$. We wish to optimize the expected utility that families derive from their available choice set.

Under a standard Multinomial Logit (MNL) discrete choice model (McFadden, 1974), each student derives utility $U_{i,s} = V_s + \varepsilon_{i,s}$ from school $s$, where $V_s$ is the deterministic mean utility and $\varepsilon_{i,s} \stackrel{\text{i.i.d.}}{\sim} \text{Gumbel}(0, 1)$ represents an idiosyncratic taste shock following a standard Type I Extreme Value distribution. The expected maximum utility of choice set $S'$ is given by the log-sum-exp formula:
\begin{equation}
U(S') = \mathbb{E} \left[ \max_{s \in S'} (V_s + \varepsilon_{i,s}) \right] = \ln \left( \sum_{s \in S'} e^{V_s} \right) + \gamma
\end{equation}
where $\gamma \approx 0.5772$ is the Euler-Mascheroni constant.

\subsection{Gain Bound}
Because the log-sum-exp function is submodular, the marginal gain from adding an additional school $s \in S \setminus S'$ exhibits diminishing marginal returns: for any $S' \subset T \subseteq S$,
\begin{equation}\label{eq:upper-bound}
\Delta^+_{s, S'} = U(S' \cup \{s\}) - U(S') \geq U(T \cup \{s\}) - U(T)
\end{equation}

We sum these bounds to establish an upper bound on the gain from adding an arbitrary subset of schools $G = \{s_1, \dots, s_{n_{\text{gain}}}\} \subseteq S \setminus S'$. Construct the intermediate sequence $S_0 = S'$ and $S_j = S_{j-1} \cup \{s_j\}$ for $j = 1, \dots, n_{\text{gain}}$, terminating at $S_{n_{\text{gain}}} = S' \cup G$. Using a telescoping sum and submodularity ($S' \subseteq S_{j-1}$), we have:
\begin{equation} \label{eq:total-gain-bound}
\begin{aligned}
U(S' \cup G) - U(S') &= \sum_{j=1}^{n_{\text{gain}}} \left( U(S_j) - U(S_{j-1}) \right) \\
&\leq \sum_{j=1}^{n_{\text{gain}}} \left( U(S' \cup \{s_j\}) - U(S') \right) = \sum_{s \in G} \Delta^+_{s, S'}
\end{aligned}
\end{equation}

\subsection{Loss Bound}
Similarly, we establish a lower bound on the utility lost from removing a school $s \in S'$. Submodularity implies that the loss incurred by removing $s$ from the complete district set $S$ is smaller than the loss from removing it from any smaller subset $S' \subseteq S$:
\begin{equation}\label{eq:loss-bound}
\Delta_s^- = U(S) - U(S \setminus \{s\}) \leq U(S') - U(S' \setminus \{s\})
\end{equation}
Crucially, $\Delta_s^-$ depends only on the global set $S$ and provides a uniform lower bound for any choice set $S'$.

For an arbitrary subset of lost schools $L = \{s_1, \dots, s_{n_{\text{loss}}}\} \subseteq S'$, construct the sequence $S_0 = S'$ and $S_j = S_{j-1} \setminus \{s_j\}$ ending at $S_{n_{\text{loss}}} = S' \setminus L$. Telescoping gives:
\begin{equation} \label{eq:total-loss-bound}
\begin{aligned}
U(S') - U(S' \setminus L) &= \sum_{j=1}^{n_{\text{loss}}} \left( U(S_{j-1}) - U(S_j) \right) \geq \sum_{s \in L} \Delta_s^- \\
U(S' \setminus L) &\leq U(S') - \sum_{s \in L} \Delta_s^-
\end{aligned}
\end{equation}

\subsection{Combined Upper Bound}
Combining the gain and loss bounds, we obtain an outer linear approximation for the utility of any arbitrary alternative choice set $\tilde{S} \subseteq S$ evaluated from baseline $S'$. Let $G = \tilde{S} \setminus S'$ be the gained schools and $L = S' \setminus \tilde{S}$ be the lost schools:
\begin{equation} \label{eq:total-bound}
U(\tilde{S}) \leq U(S') + \sum_{s \in \tilde{S} \setminus S'} \Delta_{s, S'}^+ - \sum_{s \in S' \setminus \tilde{S}} \Delta_s^-
\end{equation}

\subsection{Logic-Based Benders Decomposition}
To maximize total student utility while maintaining compactness, we reformulate cut edge minimization as a budget constraint:
$$\sum_{(i,j)\in E} b_{i,j} \leq \delta_{\text{cut}} |E|$$
where $\delta_{\text{cut}} \in (0, 1)$ sets the maximum allowable cut edge proportion.

Let $a_{v, l(s)} \in \{0, 1\}$ indicate whether vertex $v$ and school location $l(s)$ are assigned to the same zone:
\begin{equation}\label{eq:access_def}
a_{v, l(s)} \iff \bigvee_{z \in [k]} (x_{v,z} \land x_{l(s),z})
\end{equation}
which is linearized via McCormick envelopes in MIP or boolean operators in CP-SAT.

We introduce auxiliary decision variables $\hat{U}_v \in \mathbb{R}$ representing the total expected choice set utility of residents at vertex $v \in V$, and solve the resulting problem via logic-based Benders decomposition:

\begin{description}
    \item[Master Problem:] Maximize $\sum_{v \in V} \hat{U}_v$ subject to the base structural constraints \eqref{cp:assignment}--\eqref{cp:frl}, the compactness budget, access definitions \eqref{eq:access_def}, and accumulated Benders cuts. In iteration 0, $\hat{U}_v$ is unbounded from above, yielding an initial feasible partition $\mathcal{P}^0 = (x^0, z^0)$.
    \item[Subproblem:] At iteration $t$, given zoning partition $z^t$, evaluate the exact choice set available at each vertex $v$, $S_v^t = \{s \in S \mid z^t(l(s)) = z^t(v)\}$, and compute true utility $U_v(S_v^t)$ alongside marginal gains $\Delta_{t,s,v}^+ = U_v(S_v^t \cup \{s\}) - U_v(S_v^t)$ and global losses $\Delta_{s,v}^- = U_v(S) - U_v(S \setminus \{s\})$.
    \item[Benders Optimality Cuts:] Add unconditional linear outer approximation cuts to the master problem:
    \begin{equation}\label{eq:benders_cut}
    \hat{U}_v \leq U_v(S_v^t) + \sum_{s \in S \setminus S_v^t} \Delta_{t,s,v}^+ a_{v, l(s)} - \sum_{s \in S_v^t} \Delta_{s,v}^- (1 - a_{v, l(s)}) \qquad \forall v \in V
    \end{equation}
    % Alternatively, cuts may be aggregated across all vertices into a single total utility bound:
    % \begin{equation}\label{eq:benders_cut_agg}
    % \sum_{v \in V} \hat{U}_v \leq \sum_{v \in V} \left( U_v(S_v^t) + \sum_{s \in S \setminus S_v^t} \Delta_{t,s,v}^+ a_{v, l(s)} - \sum_{s \in S_v^t} \Delta_{s,v}^- (1 - a_{v, l(s)}) \right)
    % \end{equation}
\end{description}

\noindent\textbf{Theorem 1 (Finite Convergence and Optimality).} \textit{The logic-based Benders decomposition terminates in a finite number of iterations at a globally optimal zone partition $\mathcal{P}^*$.}

\textit{Proof.} Let $\mathcal{X}$ denote the finite set of feasible binary zone assignment vectors $x \in \{0, 1\}^{|V| \times k}$ satisfying the base constraints. 
First, by inequality \eqref{eq:total-bound}, every generated Benders cut \eqref{eq:benders_cut} is a valid global upper bound: for any $x \in \mathcal{X}$, the master variable satisfies $\hat{U}_v(x) \geq U_v(S_v(x))$ for all $v$.
Second, the cut is exact at the evaluated point $x^t$: setting $x = x^t$, we have $a_{v, l(s)} = 0$ for all $s \notin S_v^t$ and $(1 - a_{v, l(s)}) = 0$ for all $s \in S_v^t$, which collapses the right-hand side of \eqref{eq:benders_cut} exactly to $U_v(S_v^t)$. Thus $\hat{U}_v(x^t) = U_v(S_v^t)$.
Finally, because $\mathcal{X}$ is finite, the master problem cannot generate infinitely many distinct solutions. When the master problem produces an incumbent partition $x^{t+1}$ whose objective matches its true subproblem evaluation, we have: \kccomment{dunno if this is sufficient}
$$\sum_{v \in V} U_v(S_v(x^{t+1})) = \sum_{v \in V} \hat{U}_v(x^{t+1}) \geq \sum_{v \in V} \hat{U}_v(x) \geq \sum_{v \in V} U_v(S_v(x)) \qquad \forall x \in \mathcal{X}$$
proving that $x^{t+1}$ is a global maximum. $\blacksquare$

\section{Budget Set Optimization}
Choice set optimization maximizes the expected utility of the options accessible to a family. However, in over-subscribed urban school districts, access does not guarantee enrollment: admission is governed by centralized market clearing under school priority tiers and lottery tiebreakers. We now formulate budget set optimization to maximize total utilitarian welfare under actual admissions.

\subsection{Market-Clearing Cutoff (MID) Optimization}
We represent student demand through compressed applicant types $g \in \mathcal{G}$, where $n_g$ is the number of students of type $g$. Each type $g$ is characterized by a resident node $v(g)$, an ordered preference list of programs $(s(g,1), s(g,2), \dots, s(g,R_g))$, corresponding priority tiers $\rho(g,r) \in \mathbb{Z}_{\ge 0}$, and utilities $u(g,r)$ satisfying $u(g,1) \geq u(g,2) \geq \dots \geq u(g,R_g) > 0$.

Each program $s \in \mathcal{S}$ has capacity $q_s$ and is assigned a continuous market-clearing cutoff $P_s \in \mathbb{R}_{\geq 0}$. Let $\omega \in [0, 1]$ denote an applicant's single uniform lottery number. An applicant in priority tier $\rho$ qualifies for program $s$ if and only if $\rho + \omega \geq P_s$, or equivalently $\omega \geq P_s - \rho$. The rejected lottery mass is therefore:
\begin{equation}
t_{s, \rho} = \min\left(1, \max(P_s - \rho, 0)\right)
\end{equation}

For choice $r$ of type $g$, access depends on the zoning assignment $x$. Let $a_{g, s} \in \{0, 1\}$ indicate whether program $s$ is accessible to type $g$ ($a_{g,s} = 1$ if citywide or located in type $g$'s assigned zone). The effective rejected mass is:
\begin{equation}
e_{g,r} = 
\begin{cases}
t_{s(g,r), \rho(g,r)} & \text{if } a_{g, s(g,r)} = 1 \\
1 & \text{if } a_{g, s(g,r)} = 0
\end{cases}
\end{equation}

An applicant enrolls at choice $r$ if and only if they are rejected from all preferred choices $1, \dots, r-1$ and qualify at choice $r$. Let $R_{g,r}$ denote the remaining lottery mass rejected by choices $1, \dots, r$. Assignment follows the recurrence relation:
\begin{equation}\label{eq:mid_recurrence}
R_{g,0} = 1, \qquad R_{g,r} = \min\left(R_{g,r-1}, e_{g,r}\right) \quad \forall r \in [R_g]
\end{equation}
The assigned student mass at choice $r$ is then:
\begin{equation}
d_{g,r} = R_{g,r-1} - R_{g,r} = \max\left(0, R_{g,r-1} - e_{g,r}\right)
\end{equation}

\noindent\textbf{Proposition 1 (Recurrence Correctness).} \textit{The mass recurrence \eqref{eq:mid_recurrence} exactly matches the lottery admission interval.}

\textit{Proof.} Under uniform lottery tiebreaker $\omega \in [0, 1]$, applicant $g$ is rejected by alternative $k$ if and only if $\omega < e_{g,k}$. Consequently, $g$ is rejected by all top $r$ alternatives if and only if $\omega < \min_{k=1}^r e_{g,k}$. Defining $R_{g,r} = \min_{k=1}^r e_{g,k}$, we have $R_{g,0} = 1$ and:
$$R_{g,r} = \min_{k=1}^r e_{g,k} = \min\left( \min_{k=1}^{r-1} e_{g,k}, e_{g,r} \right) = \min\left(R_{g,r-1}, e_{g,r}\right)$$
The measure of the set of lottery numbers where applicant $g$ is rejected by choices $1, \dots, r-1$ ($\omega < R_{g,r-1}$) but accepted at choice $r$ ($\omega \geq e_{g,r}$) is the interval length $[\min(R_{g,r-1}, e_{g,r}), R_{g,r-1}]$, which is precisely $R_{g,r-1} - \min(R_{g,r-1}, e_{g,r}) = R_{g,r-1} - R_{g,r}$. $\blacksquare$

Total program demand cannot exceed seat capacity:
\begin{equation}
D_s(x, P) = \sum_{g \in \mathcal{G}} \sum_{r : s(g,r) = s} n_g d_{g,r} \leq q_s \qquad \forall s \in \mathcal{S}
\end{equation}
The MID optimization problem maximizes total utilitarian welfare:
\begin{equation}
\max_{x \in \Omega, P \ge 0} \sum_{g \in \mathcal{G}} \sum_{r=1}^{R_g} n_g u(g,r) d_{g,r} \quad \text{s.t.} \quad D_s(x, P) \leq q_s \quad \forall s \in \mathcal{S}
\end{equation}

When no program offers citywide access, both this objective and the sampled stable-matching objective of Section~\ref{sec:saa} admit an exact \emph{whole-zone} Dantzig--Wolfe decomposition: a column is one label's admissible zone together with its exactly evaluated welfare, the master tiles the district with one zone per label, and the zones are generated on demand from the master's shadow prices. Appendix~\ref{app:dw_mid} develops it --- why the base feasible set factors label by label, why zone welfare is additive under each definition, and why a single zone's sampled stable-matching value is one run of deferred acceptance rather than a solve of the access polytope of Appendix~\ref{app:access_polytope}.

\noindent\textbf{MID Continuum Decomposition:} 
\kccomment{this should def go to the appendix}
The full MID problem couples discrete zoning $x \in \Omega$ with the exact recurrence $R_{g,r} = \min(R_{g,r-1}, e_{g,r})$ across all applicant types and preference ranks. To solve this efficiently, we decompose the problem by maintaining an active prefix vector $\mathbf{k} = (k_g)_{g \in \mathcal{G}}$ with $k_g \in \{0, \dots, R_g\}$.

For each type $g$, alternatives are partitioned into an exact active prefix ($r < k_g$) and a relaxed tail ($r \geq k_g$). The tail is replaced by an optimistic transportation flow $z_{g,r} \ge 0$:
\begin{equation}
\sum_{r = k_g}^{R_g} z_{g,r} \leq R_{g, k_g-1}, \qquad 0 \leq z_{g,r} \leq a_{g,r}(x) \qquad \forall g \in \mathcal{G}, r \geq k_g
\end{equation}
The master problem maximizes the combined welfare over zoning $x \in \Omega$, cutoffs $P \ge 0$, prefix recurrence variables, and tail flows:
\begin{equation}\label{eq:mid_master}
\begin{aligned}
\max_{x, P, R, z} \quad & \sum_{g \in \mathcal{G}} n_g \left( \sum_{r < k_g} u(g,r) (R_{g,r-1} - R_{g,r}) + \sum_{r \geq k_g} u(g,r) z_{g,r} \right) \\
\text{s.t.} \quad & x \in \Omega(G, C, k, \delta) \\
& R_{g,r} = \min(R_{g,r-1}, e_{g,r}) && \forall g \in \mathcal{G}, r < k_g \\
& \sum_{r \geq k_g} z_{g,r} \leq R_{g, k_g-1}, \quad 0 \leq z_{g,r} \leq a_{g,r}(x) && \forall g \in \mathcal{G}, r \geq k_g \\
& \sum_{g, r < k_g : s(g,r)=s} n_g d_{g,r} + \sum_{g, r \geq k_g : s(g,r)=s} n_g z_{g,r} \leq q_s && \forall s \in \mathcal{S}
\end{aligned}
\end{equation}

\noindent\textbf{Proposition 3 (Valid Relaxation).} \textit{The master problem \eqref{eq:mid_master} is a valid relaxation of the full MID problem, satisfying $\text{OPT}(\text{Full MID}) \leq \text{OPT}(\text{Master}(\mathbf{k}))$.}

\textit{Proof.} Under any feasible solution to the full MID problem, true tail assignments $d_{g,r} \ge 0$ satisfy $\sum_{r \geq k_g} d_{g,r} = R_{g, k_g-1} - R_{g, R_g} \leq R_{g, k_g-1}$, obey geographic access $d_{g,r} \leq a_{g,r}(x)$, and clear program capacity. Because the master problem relaxes cutoff consistency within tails, the true solution is feasible in the master with identical objective value. $\blacksquare$

\noindent\textbf{Separation and Optimality:}
Given master candidate $(\bar{x}, \bar{P})$, true demand $D_s(\bar{x}, \bar{P})$ and assigned masses $d_{g,r}(\bar{x}, \bar{P})$ are evaluated using the complete recurrence:
\begin{enumerate}
    \item \textbf{Overload Separation:} If $D_s(\bar{x}, \bar{P}) > q_s$ for any program $s$, we identify a minimal subset of contributing types whose prefix exact demand violates capacity, and increment $k_g$ through the offending rank, cutting off the candidate.
    \item \textbf{Utility-Gap Separation:} If capacity is feasible but $\sum_{r \ge k_g} u(g,r) \bar{z}_{g,r} > \sum_{r \ge k_g} u(g,r) d_{g,r}(\bar{x}, \bar{P})$ for some type $g$, we increment $k_g$ through the first tail rank where $\bar{z}_{g,r} \neq d_{g,r}(\bar{x}, \bar{P})$.
\end{enumerate}
When no separation is possible, complete demand is feasible and true utility equals master utility, proving that candidate $\bar{x}$ is a globally optimal zoning.

\subsection{Sample Average Approximation (SAA)}
\label{sec:saa}
Under strict lottery tiebreakers, student assignment is equivalent to stable matching. Under SAA, we draw independent tie-breaking seeds $\psi \in \Psi$, yielding strict priority orderings $\succ_s^\psi$ and student preferences $\succ_i^\psi$.

For a fixed sample $\psi$ the recourse problem is the fractional stable-admissions program \eqref{saa:primal_obj}--\eqref{saa:domain}, in which the zoning enters only through the binary access indicators $a_{i,s}$ and stability is imposed by the comb inequalities \eqref{saa:comb} of Baïou and Balinski (2000). Appendix~\ref{app:access_polytope} develops that polytope in full: the shaft-and-tooth characterization of blocking, the validity and completeness of the comb family once the set of acceptable pairs is itself a decision variable, the aggregated one-row-per-pair surrogate \eqref{saa:rothblum} that the implementation loads in its place, the separation routine that generates combs, and the variant in which zoning confers a priority bonus rather than exclusive access. Three facts established there are used below.
\begin{enumerate}
    \item Every family in play is valid at every zoning and every seed, so the LP value is an upper bound on sample stable-matching welfare and any cut read off its dual is sound.
    \item With the full comb family the recourse value equals the welfare of deferred acceptance exactly, so the outer approximation is tangent at the candidate zoning rather than merely valid there.
    \item With the aggregated surrogate it is not. On $\text{Block}_2$ the surrogate over-reports realized deferred-acceptance welfare by $194.8$ to $251.8$ units across eight saved zonings, mean $225.0$, and the variation across zonings is what can misorder two candidates.
\end{enumerate}
The derivation below uses only the first.

\noindent\textbf{Dual Outer Approximation and Geographic Aggregation:}
By LP duality, the dual solution $(y^{1,\psi, r}, y^{2,\psi, r}, y^{3,\psi, r}, y^{4,\psi, r})$ computed at iteration $r$ under seed $\psi$ gives a valid global upper bound on assignment utility for any choice of student-program access indicators:
\begin{equation}
\eta_\psi \leq \kappa^{(r, \psi)} + \sum_{i \in \mathcal{I}} \sum_{s \in \mathcal{S}} \beta_{i,s}^{(r, \psi)} a_{i,s}
\end{equation}
where $\kappa^{(r, \psi)} = \sum_{i \in \mathcal{I}} y_i^{1,\psi, r} + \sum_{s \in \mathcal{S}} q_s y_s^{2,\psi, r}$ is the baseline dual welfare intercept, and $\beta_{i,s}^{(r, \psi)} = y_{i,s}^{3,\psi, r} + q_s y_{i,s}^{4,\psi, r}$ collects the multipliers of the two rows in which $a_{i,s}$ appears, so it reads as the marginal dual value of granting student $i$ access to program $s$; Appendix~\ref{app:ap_consequences} notes that this reading, and the sign that comes with it, are specific to the aggregated relaxation and do not survive the addition of comb rows.

Because school zoning access is determined entirely by geographic residence, an applicant $i$ residing at vertex $v(i) \in V$ has access to program $s$ located at school vertex $l(s) \in V$ if and only if $v(i)$ and $l(s)$ are co-zoned. That is, $a_{i,s} \equiv a_{v(i), l(s)}$ for all non-citywide programs (with $a_{i,s} = 1$ identically for citywide programs or when $v(i) = l(s)$). 

Partitioning students by home vertex $\mathcal{I}_v = \{i \in \mathcal{I} \mid v(i) = v\}$ and programs by school vertex $\mathcal{S}_u = \{s \in \mathcal{S} \mid l(s) = u\}$, the student-level dual terms sum to aggregate vertex-pair coefficients:
\begin{equation}
\Gamma_{v, u}^{(r, \psi)} = \sum_{i \in \mathcal{I}_v} \sum_{\substack{s \in \mathcal{S}_u \\ s \text{ non-citywide}}} \beta_{i,s}^{(r, \psi)} \qquad \forall v \in V, u \in S, v \neq u
\end{equation}
Fixed access terms (citywide programs and students residing at a school vertex) are absorbed directly into the intercept:
\begin{equation}
\tilde{\kappa}^{(r, \psi)} = \kappa^{(r, \psi)} + \sum_{i \in \mathcal{I}} \sum_{\substack{s \in \mathcal{S} \\ s \text{ citywide} \lor v(i) = l(s)}} \beta_{i,s}^{(r, \psi)}
\end{equation}

\noindent\textbf{Multicut Master:}
The $|\Psi|$ cuts produced at iteration $r$ could be averaged into the single hyperplane
\begin{equation}\label{eq:saa_agg_cut}
\eta \leq \bar{\kappa}^{(r)} + \sum_{v \in V} \sum_{u \in S} \bar{\Gamma}_{v, u}^{(r)} a_{v, u},
\qquad
\bar{\kappa}^{(r)} = \frac{1}{|\Psi|} \sum_{\psi \in \Psi} \tilde{\kappa}^{(r, \psi)}, \quad \bar{\Gamma}_{v, u}^{(r)} = \frac{1}{|\Psi|} \sum_{\psi \in \Psi} \Gamma_{v, u}^{(r, \psi)},
\end{equation}
bounding a single expected-welfare variable $\eta \in \mathbb{R}$. That is valid, because the sample objective is itself an average, and it is what this implementation originally did. We instead keep a separate epigraph variable $\eta_\psi$ for every sampled seed and maximize their sum --- the multicut form of Benders decomposition (Birge 1985; Birge and Louveaux 1988). Each seed's cut carries the factor $1/|\Psi|$ so that $\sum_\psi \eta_\psi$ remains on the scale of the sample average:
\begin{align}
\max_{x, a, \eta} \quad & \sum_{\psi \in \Psi} \eta_\psi \label{saa:master_obj} \\
\textrm{s.t.} \quad & x \in \Omega(G, C, k, \delta) \label{saa:master_omega} \\
& \eta_\psi \leq \frac{1}{|\Psi|} \Big( \tilde{\kappa}^{(r, \psi)} + \sum_{v \in V} \sum_{u \in S} \Gamma_{v, u}^{(r, \psi)} a_{v, u} \Big) && \forall \psi \in \Psi, \ r = 1, \dots, R \label{saa:master_cut} \\
& a_{v, u} \iff \bigvee_{z \in [k]} (x_{v, z} \land x_{u, z}) && \forall (v, u) \in \mathcal{A}^* \label{saa:master_access} \\
& 0 \leq \eta_\psi \leq \text{Welfare}_{\max} / |\Psi| && \forall \psi \in \Psi \label{saa:master_bounds} \\
& a_{v, u} \in \{0, 1\} && \forall (v, u) \in \mathcal{A}^* \label{saa:master_domain}
\end{align}
where $\mathcal{A}^* = \{(v, u) \in V \times S \mid \exists r, \psi, \ \Gamma_{v, u}^{(r, \psi)} \neq 0\}$ denotes the sparse set of active vertex-school pairs with non-zero marginal dual value, and $\text{Welfare}_{\max}$ is a global a priori upper bound developed in Appendix~\ref{app:welfare_bounds}. Since each $\eta_\psi$ bounds $\frac{1}{|\Psi|} W_\psi(x)$, the domain in \eqref{saa:master_bounds} splits the same constant $|\Psi|$ ways.

\noindent\textbf{Proposition 12 (Multicut Dominance).} \textit{Fix any pool of oracle evaluations $r = 1, \dots, R$ and let $U^{\mathrm{multi}}_R$ and $U^{\mathrm{agg}}_R$ be the optimal values of \eqref{saa:master_obj}--\eqref{saa:master_domain} and of the same master with \eqref{saa:master_cut}--\eqref{saa:master_bounds} replaced by the single-variable rows \eqref{eq:saa_agg_cut} and $0 \le \eta \le \text{Welfare}_{\max}$. Then}
\begin{equation}\label{eq:saa_multicut_chain}
\max_{x \in \Omega} \ \frac{1}{|\Psi|} \sum_{\psi \in \Psi} W_\psi(x) \;\leq\; U^{\mathrm{multi}}_R \;\leq\; U^{\mathrm{agg}}_R .
\end{equation}

\textit{Proof.} For the right inequality, take any $(x, a, \eta)$ feasible for the multicut master and set $\eta^\dagger = \sum_\psi \eta_\psi$. Summing the $|\Psi|$ rows of \eqref{saa:master_cut} at a fixed $r$ reproduces \eqref{eq:saa_agg_cut} exactly, since $\frac{1}{|\Psi|}\sum_\psi \tilde\kappa^{(r,\psi)} = \bar\kappa^{(r)}$ and likewise for the coefficients; and summing \eqref{saa:master_bounds} gives $0 \le \eta^\dagger \le \text{Welfare}_{\max}$. So $(x, a, \eta^\dagger)$ is feasible for the aggregated master with the same objective value, and the aggregated optimum is at least the multicut optimum. For the left inequality, fix any $x \in \Omega$ with induced access pattern $a(x)$. Each row of \eqref{saa:master_cut} is $1/|\Psi|$ times a valid dual bound on $W_\psi$, which holds at every access pattern, so $\eta_\psi = \frac{1}{|\Psi|} W_\psi(x)$ satisfies all of them and lies in the domain \eqref{saa:master_bounds} by validity of $\text{Welfare}_{\max}$; hence $(x, a(x), \eta)$ is feasible with objective $\frac{1}{|\Psi|}\sum_\psi W_\psi(x)$. $\blacksquare$

The two ends of \eqref{eq:saa_multicut_chain} are what matter operationally. The multicut master never reports a weaker bound than the aggregated one from the same oracle calls, because $|\Psi|$ separate rows describe $\sum_\psi \min_r(\cdot)$ whereas one averaged row describes only $\min_r \sum_\psi (\cdot)$, and the former implies the latter. The cost is $|\Psi|$ rows and $|\Psi| - 1$ extra continuous variables per iteration rather than one row, on a master that is already the expensive part of an iteration; that trade is the one empirical question the formulation raises, and it is why the aggregated form is retained as an ablation rather than deleted.

\noindent\textbf{A Priori Welfare Bound:}
The constant $\text{Welfare}_{\max}$ in \eqref{saa:master_bounds} is nominally just the domain of the epigraph variables, but it also caps the master's reported dual bound and hence the reported gap. Appendix~\ref{app:welfare_bounds} develops four candidates. Two are zoning-blind: the first-choice constant $\overline{W}^{\mathrm{FC}}$, which ignores capacity and is therefore loose by the full extent of oversubscription, and the transportation bound $\overline{W}^{\mathrm{TP}}$ of Proposition~4, which respects capacity, is attained integrally, solves in milliseconds and weakly dominates $\overline{W}^{\mathrm{FC}}$ --- on $\text{Block}_2$ by $598$ to $612$ units. Neither prices the zone confinement itself: $\overline{W}^{\mathrm{TP}}$ still hands every applicant a globally best available seat.

The other two do. Appendix~\ref{app:zoned_transport} writes the zoning model itself as a linear program, relaxing each membership variable $x_{v,z}$ to $[0,1]$ so a geographic unit may sit fractionally in several zones, reifies the co-zoning indicators on top of it, and then maximizes the zone-restricted transportation value over that relaxed polytope. The two variants, $\widehat{W}^{\mathrm{ZT}}_{\mathrm{cn}}$ and $\widehat{W}^{\mathrm{ZT}}_{\mathrm{flow}}$, differ only in whether connectedness is written with the closer-neighbor rows \eqref{cp:contiguity} or with a rooted single-commodity flow; Proposition~13 shows both are valid and both lie below $\overline{W}^{\mathrm{TP}}$. On $\text{Block}_2$ with $\texttt{max\_distance} = 3.1$ they cut a further $1{,}134$ to $1{,}241$ units off $\overline{W}^{\mathrm{TP}}$ --- roughly a third of the standing gap --- at a cost of one cached LP solve of two to three seconds. We take $\text{Welfare}_{\max} = \widehat{W}^{\mathrm{ZT}}_{\mathrm{cn}}$, which measured both the tighter and the cheaper of the two on every instance we ran and is the variant that relaxes the rows the master itself carries. What none of the four prices is stability, which is where the remaining distance to the priced bound of Section~\ref{sec:priced_access} lies.

\noindent\textbf{Computational Advantages of Geographic Aggregation:}
Aggregating the student-level duals to vertex pairs --- as opposed to aggregating across seeds, which \eqref{saa:master_cut} no longer does --- confers three advantages:
\begin{enumerate}
    \item \textbf{Dimensionality Reduction:} Rather than tracking $|\mathcal{I}| \times |\mathcal{S}| \approx 4.5 \times 10^5$ student-level access indicators, the master model instantiates binary access variables $a_{v, u}$ only on-demand for the $|\mathcal{A}^*| \approx 1.7 \times 10^3$ pairs $(v, u)$ carrying a non-zero marginal dual value. The access variables are shared across the $|\Psi|$ epigraphs, so the multicut form adds rows and continuous variables but no new integer structure.
    \item \textbf{Exact Equivalence:} Because geographic access $a_{i,s} \equiv a_{v(i), l(s)}$ is strictly identical for all students at vertex $v$, the algebraic aggregation of dual coefficients preserves exact LP outer-approximation guarantees with zero approximation error.
    \item \textbf{Unified Base Solver Architecture:} Because the cuts \eqref{saa:master_cut} share the identical access variable structure $a_{v, u}$ as the choice set objective \kccomment{hallucinated reference, fix }\eqref{eq:benders_cut_agg}, SAA directly reuses the core CP-SAT and MIP master solvers without requiring custom master models; a seed-tagged cut selects which epigraph variable it bounds and changes nothing else.
\end{enumerate}

\subsection{Congestion-Priced Access}
\label{sec:priced_access}
The residual gap left by $\overline{W}^{\mathrm{TP}}$ is the unpriced cost of confining each applicant to the programs co-zoned with their residence. Appendix~\ref{app:priced_access} closes it by replacing the \emph{constant} $\text{Welfare}_{\max}$ with a \emph{function} of the zoning that is still a valid upper bound on stable-matching welfare. Fix a non-negative price $\pi_s$ per seat at program $s$, charge it against utility to obtain the priced value $w_{i,s} = u_{i,s} - \pi_s$, and give each applicant the best priced value they can reach under the zoning, with the outside option available at zero. Summing those per-applicant terms and adding back $\sum_s q_s \pi_s$ gives the bound $\overline{W}^{\mathrm{PA}}(x; \pi)$. It is valid for \emph{every} $\pi \geq 0$ and every seed, so the prices are a free modelling choice rather than a quantity that must be computed correctly; it separates across applicants, so it admits an exact finite closed-form outer description in the master's access variables with no recourse LP at all; and the resulting cutting plane converges finitely to a certified optimum of the bound. On $\text{Block}_2$ it overstates sample-average stable-matching welfare by $249$ to $354$ units against the $3{,}877$ to $4{,}850$ of $\overline{W}^{\mathrm{TP}}$, and by nearly the same amount at zonings produced by five unrelated methods --- the empirical regularity under which maximizing the bound approximates maximizing welfare. Setting $\pi = 0$ isolates what the pricing contributes: the unpriced surrogate converges faster but is indifferent between a program with spare seats and an oversubscribed one, and gave up $1{,}357$ units of realized welfare on one instance at a matched iteration allowance. Proposition~5, Proposition~6 and Theorem~2, the vertex aggregation, the separation rule and the bias decomposition are all in Appendix~\ref{app:priced_access}.

\subsection{Empirical Comparison on Budget Welfare}
The three budget-set formulations optimize three different surrogates, so they can only be compared on a measure external to all of them. We use the welfare of the discrete MID program, written $W^{\mathrm{MID}}$ and called \emph{budget welfare} below, recomputed for every saved zoning by replaying it through the finite-grid MID oracle at base-block resolution. This is the only quantity computed identically for every strategy: each strategy's own reported objective is its own surrogate --- the sample average for SAA, the priced bound $\overline{W}^{\mathrm{PA}}$ for priced access, the MNL logsum for choice-set optimization --- and comparing those to one another compares different functions. We report neither the logsum nor any strategy's own objective here for that reason.

\noindent\textbf{Design:}
All arms solve the same $\text{Block}_2$ instance ($|\mathcal{I}| = 3953$, $|\mathcal{S}| = 123$, $k = 6$) under identical structural constraints, with a $3600\,\mathrm{s}$ solver budget, an allowance of $30$ iterations, and $16$ threads, replicated over seeds $1$--$4$ and both centroid configurations: $40$ runs, no failures. Each run contributes the single zoning, out of every zoning it saved, with the highest budget welfare --- not the one its own surrogate selected. That choice is deliberately generous and is applied uniformly; it removes the confound that the strategies disagree about which of their own iterates to return, and it matters, because $12$ of the $39$ scored runs held an iterate better than the one they handed back. Priced access is run at three price multipliers, $\pi = 0$, $\pi^\star$ and $2\pi^\star$, since Proposition 5 leaves the multiplier free.

Mean $W^{\mathrm{MID}}$ and the seed spread, largest minus smallest, over the four seeds:

\begin{center}
\begin{tabular}{lrrrr}
\hline
& \multicolumn{2}{c}{6-zone-3} & \multicolumn{2}{c}{6-zone-9} \\
arm & mean & spread & mean & spread \\
\hline
Priced access, $2\pi^\star$ & \textbf{14{,}957.9} & 233.7 & 14{,}706.2 & 269.4 \\
Priced access, $\pi^\star$  & 14{,}932.7 & 287.5 & \textbf{14{,}714.9} & 311.8 \\
Priced access, $\pi = 0$    & 14{,}808.4 & 100.6 & 14{,}599.7 & 303.8 \\
SAA                         & 14{,}777.7 & \textbf{68.4} & 14{,}578.7 & \textbf{92.5} \\
Choice set                  & 14{,}052.4 & 1{,}476.4 & 14{,}307.0 & 38.9 \\
\hline
\end{tabular}
\end{center}

\noindent\textbf{Pairing:}
The marginal spreads above are not the right yardstick for the differences, because every arm ran the same four seeds and a seed's difficulty is common to all of them. Differencing within a $(\text{centroid}, \text{seed})$ pair removes that shared component. Over the eight pairs, relative to SAA:

\begin{center}
\begin{tabular}{lrrr}
\hline
arm $-$ SAA & pairs won & mean & range \\
\hline
Priced access, $2\pi^\star$ & \textbf{7 / 8} & \textbf{$+153.9$} & $-78.7$ \ldots $+249.3$ \\
Priced access, $\pi^\star$  & 6 / 8 & $+145.6$ & $-32.0$ \ldots $+334.9$ \\
Priced access, $\pi = 0$    & 5 / 8 & $+25.9$  & $-206.5$ \ldots $+183.9$ \\
Choice set                  & 0 / 7 & $-465.0$ & $-1{,}590.3$ \ldots $-90.3$ \\
\hline
\end{tabular}
\end{center}

Priced access at $\pi^\star$ or $2\pi^\star$ improves on SAA by roughly $150$ units of budget welfare, about $1\%$, winning six or seven of eight pairs. We state that margin rather than a larger one because single-seed runs at a shorter budget suggested $+300$ to $+400$, and the replication does not support it: the effect is real, consistently signed, and smaller than one seed indicated. It is also not uniform --- one or two pairs go the other way.

Three further comparisons are worth separating from that headline. Pairing $2\pi^\star$ against $\pi^\star$ gives four wins of eight and a mean of $+8.3$, so the two multipliers are indistinguishable here and we make no claim between them. Pairing either against $\pi = 0$ gives $+128.0$ (seven of eight) and $+119.7$ (five of eight), so congestion pricing is worth about $130$ units --- real, but an order of magnitude less than the $1{,}357$ it was worth at a $900\,\mathrm{s}$ budget, because with enough iterations the unpriced surrogate largely catches up. And SAA is the \emph{most stable} arm by a wide margin, with a seed spread of $68$ and $93$ against $234$ to $312$ for priced access: it is beaten on the mean while being the most predictable. Choice-set optimization loses every pair, and its $1{,}476$ spread at 6-zone-3 includes one run that returned no feasible zoning at all within its budget, which is why its count there is three rather than four.

\noindent\textbf{Convergence and cost:}
The difference in how the two cutting-plane families terminate is larger than the difference in what they find. Across the $24$ priced-access runs, the relative optimism gap between the master's model value and the true value of $\sum_i g_i$ at the same zoning falls from a median of $29.78\%$ at the first iteration to a median of $0.013\%$, and $23$ of the $24$ runs stop because \emph{no cut separates} --- the exact termination certificate of Theorem 2 --- after a median of $10$ iterations and $466\,\mathrm{s}$, never needing the $30$-iteration allowance or the $3600\,\mathrm{s}$ budget. Across the eight SAA runs the same gap falls from $29.23\%$ only to $7.95\%$, and every run exhausts its iteration allowance at $30$ iterations and $4{,}555\,\mathrm{s}$ of wall time. Priced access therefore reaches a certified optimum of its surrogate in roughly a tenth of the wall time SAA spends failing to close its own gap, and the reason is structural rather than incidental: its cuts are closed form, so an iteration costs one master solve, whereas each SAA iteration also pays for $|\Psi| = 15$ recourse LPs.

Two cautions attach to that. First, the certificate is about $\overline{W}^{\mathrm{PA}}$, not about $W^{\mathrm{MID}}$; the run stops because nothing more can be gained on the surrogate, and the remaining distance to the true optimum is bounded by the surrogate's bias rather than by anything the run measures. Second, choice-set optimization also closes its gap, to a median of $0.01\%$, but its surrogate is the MNL logsum, which is neither an upper nor a lower bound on realized welfare, so its gap certifies nothing about the matching; its poor budget welfare despite an almost exactly converged model is the sharpest illustration in this table that converging on the wrong surrogate is not progress.

Compactness does not have to be traded away for this. At 6-zone-3 the priced-access arms return $1{,}128$ to $1{,}136$ cut edges against $1{,}265$ for SAA, so they are simultaneously better on both panels; at 6-zone-9 they range from $1{,}184$ to $1{,}225$ against SAA's $1{,}179$, which is a wash.

\noindent\textbf{Scope:}
These are four seeds on one instance family at one district size, so the ordering should be read as replicated on this instance rather than established in general, and the $\pm 150$ margin over SAA is small enough that a different instance could plausibly reverse it. What is not marginal, and what we would expect to transfer, is the termination behaviour: a bound that separates across applicants admits an exact finite cut family and therefore a real stopping certificate, whereas one dual hyperplane per seed per iteration over $|\mathcal{I}| \times |\mathcal{S}|$ access coordinates does not, independently of which zoning either one happens to return.

\appendix
\section{Extended Demographic Diversity Constraints}
\label{app:diversity_constraints}
In addition to socioeconomic (FRL) balance, the optimization framework supports multi-group racial and ethnic demographic balance across an arbitrary set of tracked groups $\mathcal{E}$.

Let $r_{v,e} \ge 0$ denote the count of students of ethnicity $e \in \mathcal{E}$ residing in vertex $v \in V$. The district-wide proportion of group $e$ is:
\begin{equation}
R_e = \frac{\sum_{v \in V} r_{v,e}}{\sum_{v \in V} n_v}
\end{equation}
For each group $e \in \mathcal{E}$ and each zone $Z_i \in \mathcal{P}$, the localized proportion must not deviate from $R_e$ by more than tolerance $\delta_R \in \mathbb{R}_{>0}$:
\begin{equation}
\left| \frac{\sum_{v \in Z_i} r_{v,e}}{N(Z_i)} - R_e \right| \leq \delta_R \qquad \forall e \in \mathcal{E}, \forall Z_i \in \mathcal{P}
\end{equation}

In the CP-SAT and MIP formulations, this is linearized into upper and lower integer bounds:
\begin{align}
\sum_{v \in V} \lfloor M \cdot (r_{v,e} - (R_e - \delta_R) n_v) \rfloor \cdot x_{v,z} &\geq 0 \qquad \forall e \in \mathcal{E}, \forall z \in [k] \\
\sum_{v \in V} \lceil M \cdot (r_{v,e} - (R_e + \delta_R) n_v) \rceil \cdot x_{v,z} &\leq 0 \qquad \forall e \in \mathcal{E}, \forall z \in [k]
\end{align}
which are added seamlessly alongside the base FRL constraints without altering the core solver mechanics.

\section{The Stable-Matching Access Polytope}
\label{app:access_polytope}

This appendix develops the sample recourse problem of Section~\ref{sec:saa}. Section~\ref{app:ap_market} fixes the market and the shaft--tooth apparatus, Section~\ref{app:ap_combs} states the comb family and proves it valid and complete at a fixed access pattern, Section~\ref{app:ap_variable} extends both to a market whose set of acceptable pairs is itself a decision variable --- the case the zoning master needs --- and Section~\ref{app:ap_sep} gives the separation routine. Section~\ref{app:ap_bonus} then replaces exclusive access by a zone \emph{priority bonus}, which is the status quo the district actually operates: residence raises a family's priority at the programs in its zone without excluding anyone. The same polytope survives once each applicant is split into a zoned and an unzoned copy. Section~\ref{app:ap_bonus_sep} gives the corresponding separation and Section~\ref{app:ap_consequences} records what changes in the master cuts.

\subsection{The Sample Market}
\label{app:ap_market}

Fix a tie-breaking seed $\psi \in \Psi$. Applicants are $\mathcal{I}$, programs are $\mathcal{S}$, and $\Gamma = \{(i,s) : s \text{ appears on } i\text{'s list}\}$ is the set of acceptable pairs; $|\Gamma| = 12{,}370$ on $\text{Block}_2$ against $|\mathcal{I}| \times |\mathcal{S}| \approx 4.9 \times 10^5$, because ranked lists average $3.13$ entries. Program $s$ has $q_s$ seats and, under seed $\psi$, a strict order $\succ_s^\psi$ over $\{i : (i,s) \in \Gamma\}$ obtained by sorting the priority tier $\rho_{i,s}$ and breaking ties by the lottery draw $\omega_{i,s}^\psi$. Applicant $i$ has a strict order $\succ_i^\psi$ over their list and utilities $u_{i,s}$ consistent with it, so $s \succ_i^\psi s'$ implies $u_{i,s} \ge u_{i,s'}$. We write $\succeq$ for the reflexive closures. Following the proof of Proposition~4, $\mathcal{S}$ contains an outside option $s_0$ with $u_{i,s_0} = 0$, $q_{s_0} = |\mathcal{I}|$, universally acceptable and ranked last by everyone, so $\sum_{s} D_{i,s} = 1$ loses no generality and unassignment is the mass placed on $s_0$.

Zoning enters through the access indicators $a_{i,s} \in \{0,1\}$ of Section~\ref{sec:saa}: $a_{i,s} = 1$ when $s$ is citywide, when $l(s) = v(i)$, or when $v(i)$ and $l(s)$ are co-zoned, and $a_{i,s} = 0$ otherwise. It is convenient to name the market that a given access pattern leaves behind,
\begin{equation}\label{eq:ap_restricted_market}
\Gamma(a) = \{(i,s) \in \Gamma : a_{i,s} = 1\},
\end{equation}
and to write $n_s(a) = |\{i : (i,s) \in \Gamma(a)\}|$ for the number of applicants who both rank $s$ and may attend it. A pair $(i,s) \in \Gamma(a)$ \emph{blocks} an integral matching $D$ if $i$ prefers $s$ to their assignment and $s$ either has a free seat or has admitted someone below $i$ in $\succ_s^\psi$; $D$ is stable if no pair of $\Gamma(a)$ blocks it.

\subsection{Shafts, Teeth and Combs}
\label{app:ap_combs}

The description below is that of Baïou and Balinski (2000), in the notation of Sethuraman, Teo and Qian (2006). For $(i,s) \in \Gamma$ define the \emph{shaft} and the \emph{tooth}
\begin{equation}\label{eq:ap_shaft_tooth}
S^\psi(s,i) = \{(i',s) \in \Gamma : i' \succeq_s^\psi i\},
\qquad
T^\psi(i,s) = \{(i,s') \in \Gamma : s' \succeq_i^\psi s\},
\end{equation}
and for any $\hat\Gamma \subseteq \Gamma$ write $D(\hat\Gamma) = \sum_{(i,s) \in \hat\Gamma} D_{i,s}$. The shaft collects the seats of $s$ that go to applicants $s$ likes at least as much as $i$; the tooth collects the outcomes $i$ likes at least as much as $s$. Both are prefixes of an order and are therefore running sums, which is what makes the whole apparatus cheap: the implementation carries one prefix variable per element of $\Gamma$ on each side, $2|\Gamma|$ in total, rather than re-summing.

\noindent\textbf{Lemma 1 (Blocking in shaft--tooth form).} \textit{Let $D$ be an integral matching and $(i,s) \in \Gamma(a)$. Then $(i,s)$ blocks $D$ if and only if}
\begin{equation}\label{eq:ap_blocking}
D\big(S^\psi(s,i)\big) \le q_s - 1
\qquad\text{and}\qquad
D\big(T^\psi(i,s)\big) = 0 .
\end{equation}

\textit{Proof.} The second condition says $i$ receives nothing weakly preferred to $s$, which is exactly the applicant half of blocking. For the program half, $D(S^\psi(s,i))$ counts the applicants weakly above $i$ admitted to $s$. If that count is at most $q_s - 1$ then the remaining $q_s - D(S^\psi(s,i)) \ge 1$ seats are either empty or held by applicants strictly below $i$, so $s$ can accept $i$, dropping its worst admit if necessary. Conversely if the count is $q_s$ then $s$ is full with applicants it weakly prefers to $i$ and gains nothing from $i$. Note that $i$ itself is in $S^\psi(s,i)$, so under $D(T^\psi(i,s)) = 0$ the count is unchanged by $i$'s own row. $\blacksquare$

Lemma~1 is what the comb construction linearizes. A single applicant witnesses only one seat, so certifying that all $q_s$ seats are unavailable to $i$ requires $q_s$ witnesses. Fix $(i,s) \in \Gamma$ and a collection $K$ of $q_s - 1$ \emph{teeth bases}, distinct except that the vacant slot may repeat, each drawn from
\begin{equation}\label{eq:ap_tooth_bases}
\mathcal{K}^\psi(s,i) = \{k \in \mathcal{I} : (k,s) \in \Gamma, \ k \succ_s^\psi i\} \cup \{\varnothing\},
\end{equation}
where $\varnothing$ is a \emph{vacant slot}, explained below. The \emph{comb} based at $(i,s)$ with teeth $K$ is
\begin{equation}\label{eq:ap_comb_set}
C^\psi(s,i,K) = S^\psi(s,i) \ \cup \ \big(T^\psi(i,s) \setminus \{(i,s)\}\big) \ \cup \ \bigcup_{k \in K} \big(T^\psi(k,s) \setminus \{(k,s)\}\big),
\end{equation}
a shaft with $q_s$ teeth hanging off it, one for the base and one for each element of $K$, with the convention that a vacant slot contributes no pairs to the union and enters only through its credit. Each tooth is stripped of its own $(\cdot,s)$ entry because that entry already lies in the shaft, so nothing is double counted. The comb inequality asserts that a comb carries at least $q_s$ units of mass. Written out, and with access and vacancies priced as derived in Section~\ref{app:ap_variable},
\begin{align}
\max_{D} \quad & \sum_{i \in \mathcal{I}} \sum_{s \in \mathcal{S}} u_{i,s} D_{i,s} \label{saa:primal_obj} \\
\textrm{s.t.} \quad & \sum_{s \in \mathcal{S}} D_{i,s} = 1 && \forall i \in \mathcal{I} \label{saa:assignment} \\
& \sum_{i \in \mathcal{I}} D_{i,s} \leq q_s && \forall s \in \mathcal{S} \label{saa:capacity} \\
& D_{i,s} \leq a_{i,s} && \forall (i,s) \in \Gamma \label{saa:access} \\
& \underbrace{\sum_{i' \succeq_s^\psi i} D_{i', s}}_{\text{shaft}} + \underbrace{\sum_{s' \succ_i^\psi s} D_{i, s'}}_{\text{base tooth}} + \sum_{k \in K} \underbrace{\Big[ \sum_{s' \succ_k^\psi s} D_{k, s'} + \chi_{k,s} \Big]}_{\text{tooth } k} \geq q_s\, a_{i,s} \label{saa:comb} \\
& \qquad \forall (i,s) \in \Gamma \text{ and all admissible } K \nonumber \\
& D_{i,s} \geq 0 && \forall (i,s) \in \Gamma \label{saa:domain}
\end{align}
where the \emph{tooth credit} is
\begin{equation}\label{eq:ap_credit}
\chi_{k,s} = 1 - a_{k,s} \ \ \text{for } k \in \mathcal{I},
\qquad
\chi_{\varnothing,s} = 1 .
\end{equation}
The credit is the only departure from the fixed-market description, and Section~\ref{app:ap_variable} is where it is earned. \kccomment{ai speak} For the moment read \eqref{saa:comb} at $a \equiv 1$ and with no vacant slots, where every credit is zero and the inequality is literally $D(C^\psi(s,i,K)) \ge q_s$.

\noindent\textbf{Proposition 2 (Comb Inequality Stability).} \textit{Fix $\psi$ and a binary access pattern $a$. Then the integral points of \eqref{saa:assignment}--\eqref{saa:domain} are exactly the stable matchings of the market $\Gamma(a)$, and the system has no fractional extreme points: its solution set is $\mathrm{conv}\{D : D \text{ stable in } \Gamma(a)\}$.} 

\textit{Proof.} Validity and integral exactness are Proposition~10 below, which proves them for all $a$ at once. For the absence of fractional extreme points, take $a \equiv 1$ first and assume $n_s \ge q_s$ for every $s$. Then all credits vanish, \eqref{saa:comb} is the comb inequality $D(C) \ge q_s$, and the claim is the theorem of Baïou and Balinski (2000): the system \eqref{saa:assignment}--\eqref{saa:domain} is the stable admissions polytope. Sethuraman, Teo and Qian (2006) give a short proof by exhibiting, for each fractional point, a decomposition of its mass into $q_s$ unit bins per program such that the level sets $\{$best $q_s$ applicants at height $\alpha\}$ are stable matchings for every $\alpha \in (0,1]$; the point is then the convex combination of those level sets induced by $\alpha$ uniform. The general case is the second half of Proposition~10. $\blacksquare$ \kccomment{i dont get this part}

The completeness claim has a consequence worth isolating, because it says what the recourse oracle would return if the comb family were loaded in full.

\noindent\textbf{Corollary 1 (The comb-complete recourse value is deferred-acceptance welfare).} \textit{For every zoning $x$ and seed $\psi$, the optimal value of \eqref{saa:primal_obj}--\eqref{saa:domain} equals the utilitarian welfare of the applicant-proposing deferred-acceptance outcome on $\Gamma(a(x))$.}

\textit{Proof.} By Proposition~2 the feasible set is the convex hull of the stable matchings of $\Gamma(a)$, so the optimum is attained at one of them and equals the largest welfare of a stable matching. Deferred acceptance returns the applicant-optimal stable matching: every applicant weakly prefers it to every other stable matching (Gale and Shapley, 1962; Roth and Sotomayor, 1990). Since $u_{i,\cdot}$ is consistent with $\succ_i^\psi$, each applicant's utility is weakly largest there, hence so is the sum. $\blacksquare$

Corollary~1 is the reason the comb family is worth stating even though the implementation does not load it. An oracle carrying it would be tangent to realized welfare at the candidate zoning, so each seed's Benders cut \eqref{saa:master_cut} would be exact at $x^r$ rather than merely valid, and the incumbent selection problem noted in Section~\ref{sec:saa} would not arise.

\subsection{Access as a Decision Variable}
\label{app:ap_variable}

The zoning master does not solve one market. It searches over $x$, and each $x$ induces a different set $\Gamma(a(x))$ of acceptable pairs. A description of the stable matchings of a \emph{fixed} market is therefore not enough: the master needs one family of inequalities, written in the variables $(D,a)$, that is valid simultaneously at every binary $a$ and that specialises to the complete description of $\Gamma(a)$ at each one. Two things obstruct the naive attempt of multiplying the right-hand side by $a_{i,s}$ and leaving the comb otherwise alone.

First, a tooth base without access is not a competitor. Let $q_s = 2$ with applicants $1 \succ_s^\psi 2$, both ranking only $s$, and let $a_{1,s} = 0$, $a_{2,s} = 1$. The stable matching of $\Gamma(a)$ assigns $2$ to $s$ and leaves $1$ unassigned, but the uncredited comb based at $(2,s)$ with tooth $1$ reads $D_{1,s} + D_{2,s} + 0 + 0 = 1 < 2 = q_s a_{2,s}$. The row is violated by a genuine stable matching, so the family without credits is not valid.

Second, a program may have fewer than $q_s$ interested applicants with access, in which case there are not $q_s - 1$ genuine competitors above $i$ to serve as teeth and, absent the vacant slot, the base $(i,s)$ carries no usable row --- yet stability has content there, since a program that cannot fill must admit everyone who wants it. Both obstructions are the same phenomenon and both are repaired by \eqref{eq:ap_credit}: an applicant who cannot compete for $s$, whether because they lack access or because they do not exist, is credited a full unit, because for the purpose of $s$'s market they are indistinguishable from an applicant already served better than $s$.

\noindent\textbf{Proposition 10 (Zoning-robust comb inequalities).} \textit{Let $\psi$ be fixed.}
\begin{enumerate}
    \item[(i)] \textit{\emph{(Uniform validity.)} Every inequality \eqref{saa:comb} holds at every pair $(D,a)$ with $a \in \{0,1\}^{\Gamma}$ and $D$ an integral stable matching of $\Gamma(a)$.}
    \item[(ii)] \textit{\emph{(Integral exactness.)} Conversely, if $a$ is binary and $D$ is integral and satisfies \eqref{saa:assignment}--\eqref{saa:domain}, then $D$ is a stable matching of $\Gamma(a)$.}
    \item[(iii)] \textit{\emph{(Specialisation.)} At each fixed binary $a$, the rows of \eqref{saa:comb} imply Baïou and Balinski's complete description of the market $\Gamma(a)$, padded as required when $n_s(a) < q_s$. Hence the solution set is exactly $\mathrm{conv}\{D : D \text{ stable in } \Gamma(a)\}$, completing Proposition~2.}
\end{enumerate}

\textit{Proof.} (i) If $a_{i,s} = 0$ the right-hand side is zero and the left-hand side is a sum of non-negative terms. So assume $a_{i,s} = 1$, fix teeth $K$, and write $P = \{i\} \cup K$, a collection of $q_s$ elements, and
$$b_i = \sum_{s' \succ_i^\psi s} D_{i,s'},
\qquad
b_k = \sum_{s' \succ_k^\psi s} D_{k,s'} + \chi_{k,s} \ \ (k \in K),$$
so that the left-hand side of \eqref{saa:comb} is $D(S^\psi(s,i)) + \sum_{p \in P} b_p$. Suppose it is at most $q_s - 1$. Partition $P$ into
$$G = \{p \in P : b_p \ge 1\},
\qquad
F = \{p \in P \setminus G : D_{p,s} = 1\},
\qquad
H = P \setminus (F \cup G).$$
Every $p \in F$ has $b_p = 0$, hence $\chi_{p,s} = 0$, hence $a_{p,s} = 1$ and $p \succeq_s^\psi i$ --- vacant slots and applicants without access carry a credit and lie in $G$ --- so $(p,s) \in S^\psi(s,i)$ and $D(S^\psi(s,i)) \ge |F|$. Therefore the left-hand side is at least $|F| + |G|$, and since $|F| + |G| + |H| = q_s$ the assumed bound forces $H \neq \emptyset$. Pick $k \in H$. As above $\chi_{k,s} = 0$, so $(k,s) \in \Gamma(a)$ and $k \succeq_s^\psi i$; and $b_k = 0$ with $D_{k,s} = 0$ gives $D(T^\psi(k,s)) = 0$. Finally $k \succeq_s^\psi i$ gives $S^\psi(s,k) \subseteq S^\psi(s,i)$, so
$$D\big(S^\psi(s,k)\big) \le D\big(S^\psi(s,i)\big) \le q_s - 1 .$$
By Lemma~1 the pair $(k,s) \in \Gamma(a)$ blocks $D$, contradicting stability. \kccomment{i dont understand this section}

(ii) Let $a$ be binary and $D$ integral and feasible. Constraints \eqref{saa:assignment}, \eqref{saa:capacity} and \eqref{saa:access} make $D$ a matching of $\Gamma(a)$. Suppose $(i,s) \in \Gamma(a)$ blocked it, so by Lemma~1 $D(T^\psi(i,s)) = 0$ and $D(S^\psi(s,i)) \le q_s - 1$. Let $A = \{k : (k,s) \in \Gamma, \ k \succ_s^\psi i\}$ and, with $b_k$ as in part~(i), let $m_0 = |\{k \in A : b_k = 0\}|$. Choose the teeth greedily, taking the $q_s - 1$ smallest brackets of $\mathcal{K}^\psi(s,i)$; since vacant slots have bracket exactly $1$ and are available without limit,
$$\sigma := \sum_{k \in K} b_k \le \max\{0, \ q_s - 1 - m_0\} .$$
The base tooth vanishes because $D(T^\psi(i,s)) = 0$, so the left-hand side of \eqref{saa:comb} is $D(S^\psi(s,i)) + \sigma$. Now $D_{i,s} = 0$, so $D(S^\psi(s,i))$ counts the members of $A$ admitted to $s$; each of those has access and nothing strictly preferred to $s$, hence bracket $0$, hence $D(S^\psi(s,i)) \le m_0$. If $m_0 \ge q_s - 1$ then $\sigma = 0$ and the left-hand side is $D(S^\psi(s,i)) \le q_s - 1$; otherwise it is at most $m_0 + (q_s - 1 - m_0) = q_s - 1$. Either way it falls below $q_s = q_s a_{i,s}$ and this row of \eqref{saa:comb} is violated, contradicting feasibility. So no pair of $\Gamma(a)$ blocks $D$.

(iii) Fix binary $a$. A program with $n_s(a) \ge q_s$ needs no padding, so consider one with $n_s(a) < q_s$ and pad $\Gamma(a)$ with $\varphi_s = q_s - n_s(a)$ phantom applicants whose only acceptable program is $s$ and who occupy the top $\varphi_s$ positions of $\succ_s^\psi$, each with utility zero. In every stable matching of the padded market each phantom is admitted to $s$: were one rejected, $s$ would have to be full, but the admitted set draws from $n_s(a)$ real applicants and $\varphi_s - 1$ other phantoms, at most $q_s - 1$ in all. Real applicants therefore face effective capacity $n_s(a)$, which does not bind, so the stable matchings of the padded market restrict exactly onto those of $\Gamma(a)$ and carry the same welfare. Baïou and Balinski's description applies to the padded market, in which $n_s \ge q_s$ for every program. Eliminate the phantom variables, all of which are fixed at one: a padded comb based at $(i,s)$ using $j$ phantom teeth has left-hand side equal to the real shaft plus $\varphi_s$ --- every phantom lies in the shaft, whether or not it was chosen as a tooth --- plus the base tooth plus the $q_s - 1 - j$ real tooth terms, since a phantom's own tooth is empty. The row of \eqref{saa:comb} with the same real teeth and $j$ vacant slots has the same right-hand side and a left-hand side smaller by $\varphi_s - j \ge 0$, hence implies it. Applicants without access are handled identically, being vacant slots of $\Gamma(a)$ under another name: their credit is $1$ and their shaft contribution is $0$ by \eqref{saa:access}. So every inequality of the complete description is implied, giving solution set $\subseteq \mathrm{conv}$; part~(i) gives $\supseteq$. $\blacksquare$

\kccomment{everything here is confusing}

\noindent\textbf{What the implementation loads.}
The comb family has up to $\binom{n_s - 1}{q_s - 1}$ rows per base pair and is separated, not enumerated. What \texttt{SaaOracle} loads up front is one aggregated stability row per acceptable pair,
\begin{equation}\label{saa:rothblum}
\sum_{i' \succ_s^\psi i} D_{i',s} \;+\; q_s \sum_{s' \succeq_i^\psi s} D_{i,s'} \;\ge\; q_s\, a_{i,s}
\qquad \forall (i,s) \in \Gamma ,
\end{equation}
the many-to-one aggregation of Rothblum's (1992) inequality. It is valid: if $a_{i,s} = 1$ and $i$ receives nothing weakly preferred to $s$ then the second sum is zero and stability forces $q_s$ seats to applicants strictly above $i$, while if $i$ does receive something weakly preferred the second sum is $1$ and alone supplies $q_s$. Being valid, it is implied by the complete description of Proposition~2 and may be added freely. It is also cheap: with the prefix variables of \eqref{eq:ap_shaft_tooth} it is $O(|\Gamma|)$ rows, $12{,}370$ on $\text{Block}_2$, and the $y^{4}$ block of the dual in Section~\ref{sec:saa} is precisely its multiplier vector. Use the \emph{strict} shaft as written; the weak variant is valid too but measurably looser, because it lets $i$'s own mass at $s$ be counted twice.

\kccomment{ai speak}What \eqref{saa:rothblum} is not is complete, and the reason is visible in the derivation. Summing the per-seat inequalities of a seat-cloned program and bounding $\sum_{j} (q_s - j + 1) D_{i, s^{(j)}}$ by $q_s D_{i,s}$, where $s^{(1)}, \dots, s^{(q_s)}$ are the clones, discards the distinction between one applicant holding a whole seat and several holding fractions of it; the multiplier $q_s$ on the own-preference term is what lets a fractional $D_{i,s}$ pay for a deficit in the shaft. Concretely, let $q_s = 2$ and let applicants $1 \succ_s^\psi 2 \succ_s^\psi 3$ rank only $s$, with $u_{1,s} = u_{2,s} = 1$ and $u_{3,s} = 10$, all access granted. The unique stable matching admits $1$ and $2$, for welfare $2$. The point $D_{1,s} = 1$, $D_{2,s} = D_{3,s} = \tfrac12$ satisfies \eqref{saa:assignment}--\eqref{saa:access}, \eqref{saa:domain} and every row of \eqref{saa:rothblum} --- the three rows read $0 + 2 = 2$, $1 + 1 = 2$ and $\tfrac32 + 1 = \tfrac52$ against $q_s = 2$ --- and has objective $6.5$, more than three times the truth. The comb based at $(2,s)$ with the single tooth $1$ is what removes it: its left-hand side is $D_{1,s} + D_{2,s} = \tfrac32 < 2$.

The gap is not merely theoretical. Because Corollary~1 identifies the comb-complete value with deferred-acceptance welfare, the discrepancy between the surrogate LP value and a deferred-acceptance run measures exactly the slack of \eqref{saa:rothblum}. Over eight saved zonings on $\text{Block}_2$ with $k = 6$ the surrogate over-reports realized deferred-acceptance welfare by $194.8$ to $251.8$ units, mean $225.0$, against a total welfare near $14{,}500$; measured instead against the MID continuum value on the same zonings the over-report is a very stable $194.8$ to $195.0$. The variation of roughly $57$ units \emph{across} zonings is the operationally relevant part, since it is what can reorder two candidates: the cuts stay valid, but a strategy that selects its incumbent by surrogate value can select wrongly by that much. This is why Section~\ref{sec:saa} reports budget welfare from an independent replay rather than from the oracle.

\subsection{Comb Separation Under Access}
\label{app:ap_sep}

Separation is exact and requires no optimization beyond a sort. Fix a candidate zoning $x^r$, hence a binary $a$, and a fractional $\bar D$ solving the current restricted LP. For $(k,s) \in \Gamma$ let
\begin{equation}\label{eq:ap_bracket}
\bar b_{k,s} = \sum_{s' \succ_k^\psi s} \bar D_{k,s'} + (1 - a_{k,s}),
\end{equation}
the bracket of \eqref{saa:comb}; all $|\Gamma|$ of these are available as the tooth prefix sums already in the model. Because the teeth enter \eqref{saa:comb} additively and their admissible set $\mathcal{K}^\psi(s,i)$ depends on the base only through $k \succ_s^\psi i$, the most violated comb based at $(i,s)$ is obtained by choosing the $q_s - 1$ smallest brackets above $i$:
\begin{equation}\label{eq:ap_most_violated}
\mathrm{viol}(i,s) = q_s a_{i,s} - \Big[ \underbrace{\textstyle\sum_{i' \succeq_s^\psi i} \bar D_{i',s}}_{\text{shaft prefix}} + \sum_{s' \succ_i^\psi s} \bar D_{i,s'} + \underbrace{\min_{\substack{K \subseteq \mathcal{K}^\psi(s,i) \\ |K| = q_s - 1}} \sum_{k \in K} \bar b_{k,s}}_{\sigma(i,s)} \Big],
\end{equation}
and vacant slots, whose bracket is $1$, enter this minimum on the same footing as everything else --- they are selected exactly when fewer than $q_s - 1$ applicants above $i$ have a bracket below $1$, which is the case Section~\ref{app:ap_variable} introduced them for. One descent of $\succ_s^\psi$ computes \eqref{eq:ap_most_violated} at every base simultaneously. Carry two accumulators: the shaft, a running sum, and $\sigma$, the sum of a max-heap of size $q_s - 1$ holding the smallest brackets seen so far, initialized with $q_s - 1$ vacant slots. On reaching applicant $i$, first add $\bar D_{i,s}$ to the shaft --- the shaft of \eqref{eq:ap_shaft_tooth} is weak and includes $i$'s own row --- then evaluate \eqref{eq:ap_most_violated}, then push $\bar b_{i,s}$ and pop the maximum, so that on reaching the next applicant the heap holds exactly the teeth strictly above them. Popping the maximum evicts a vacant slot only when a real bracket is smaller, which is why brackets above $1$ --- possible for an applicant who lacks access and is served better than $s$ --- never displace a slot. The cost is $O(|\Gamma| + \sum_s n_s \log q_s)$, dominated by $O(|\Gamma| \log \max_s q_s)$, which on $\text{Block}_2$ is a few hundred thousand operations against an LP re-solve; separation is free relative to the oracle it feeds. \kccomment{eh could be more clear}

The oracle is then a cutting plane in its own right, nested inside the zoning cutting plane: load \eqref{saa:assignment}--\eqref{saa:access}, \eqref{saa:domain} and \eqref{saa:rothblum}, solve, run the descent above, add every comb with $\mathrm{viol}(i,s) > \epsilon$, and repeat until none is found. Termination is finite because the comb family is finite and no row is added twice, and on exit Proposition~2 gives a value equal to deferred-acceptance welfare by Corollary~1. Two implementation notes. Rows should be added in batches --- all violated bases at once --- since the marginal cost of a row is small next to a re-solve. And the emitted row must record its teeth as \emph{indices}, not as the numeric bracket values that selected them: the row is an inequality in $(D,a)$ valid at every zoning, and re-deriving it from $\bar D$ at the next master iterate would silently produce a different, weaker row.

\subsection{Zoning as a Priority Bonus}
\label{app:ap_bonus}

Exclusive access is the sharpest form a zone can take, but it is not the form the district operates. \kccomment{this is imprecise}In the status quo a family receives a priority \emph{boost} at its attendance-area program and may still apply anywhere; the zone changes who wins a contested seat, not who may enter the contest. This subsection writes the polytope for that policy. The construction is the natural one: every applicant is split into a \emph{zoned} and an \emph{unzoned} copy, the priority structure is written once for each copy, and the zoning decides which copy is real.

\noindent\textbf{Data.} Let $\rho_{i,s} \in \mathbb{Z}_{\ge 0}$ be the base priority score of $i$ at $s$, namely the sum of the configured tier weights --- CTIP, sibling, pre-K, round and the rest --- with the zone term removed; let $\beta \in \mathbb{Z}_{\ge 1}$ be the zone weight, the single policy dial of this subsection; and let $\omega_{i,s}^\psi \in (0,1)$ be the lottery draw under seed $\psi$, with $\omega_{i,s}^\psi \neq \omega_{j,s}^\psi$ whenever $i \neq j$. Single and multiple tie-breaking differ only in whether $\omega_{i,s}^\psi$ depends on $s$, and nothing below uses either. Program $s$ orders applicants by decreasing realized score $\rho_{i,s} + \beta a_{i,s} + \omega_{i,s}^\psi$. Integrality of $\rho$ and $\beta$ is used only to keep realized scores separated and may be replaced by the requirement that $\max_{i,s} \omega^\psi_{i,s}$ be smaller than the least positive gap between distinct values of $\rho_{i,s} + \beta a_{i,s}$.

\noindent\textbf{The bonus indicator.} $a_{i,s} = 1$ when $l(s)$ and $v(i)$ are co-zoned, which includes $l(s) = v(i)$. This is the same quantity as in the access model and the same master variable $a_{v(i), l(s)}$ of \eqref{saa:master_access}, so the zoning side of the problem is untouched: only the recourse changes. Two families of pairs are constant rather than variable and must be fixed as data. The outside option carries no bonus, $a_{i,s_0} \equiv 0$. For a citywide program we take residence in the zone containing it to earn no bonus either, $a_{i,s} \equiv 0$. The access model instead fixes $a_{i,s} = 1$ there, but only to make such programs universally reachable; under a bonus the requirement is merely that the value be constant, since a uniform shift of one program's scores leaves its order unchanged, so either constant serves. Letting $a_{i,s}$ be the co-zoning indicator at citywide programs as well is a different \emph{policy}, not a different convention, and is covered by everything below without modification.

Each applicant carries two \emph{keys} at each program,
\begin{equation}\label{eq:ap_keys}
\kappa^{\mathrm{in}}_{i,s} = \rho_{i,s} + \omega_{i,s}^\psi + \beta,
\qquad
\kappa^{\mathrm{out}}_{i,s} = \rho_{i,s} + \omega_{i,s}^\psi ,
\end{equation}
of which the zoning selects one: the realized score of $i$ at $s$ is $\kappa^{\mathrm{in}}_{i,s}$ if $a_{i,s} = 1$ and $\kappa^{\mathrm{out}}_{i,s}$ otherwise. Both keys are data; the selection is the decision.

\noindent\textbf{Merged key order.} Fix $s$ and let $n_s = |\{i : (i,s) \in \Gamma\}|$. The $2 n_s$ keys \eqref{eq:ap_keys} are pairwise distinct, since equality of any two forces two lottery draws to differ by an integer and hence to coincide. Sort them once into strictly decreasing order and write $\mathrm{pos}_s(i,z) \in \{1, \dots, 2 n_s\}$ for the position of $\kappa^{z}_{i,s}$ and $(i_s(m), z_s(m))$ for the applicant and tier occupying position $m$. Every zoning realizes exactly one key per applicant, so the realized order $\succ_s^{\psi,a}$ is the subsequence of the merged order at the realized positions and, writing $z(\cdot)$ for realized tiers,
\begin{equation}\label{eq:ap_pos_order}
k \succ_s^{\psi,a} i
\quad\Longleftrightarrow\quad
\mathrm{pos}_s\big(k, z(k)\big) < \mathrm{pos}_s\big(i, z(i)\big).
\end{equation}
Programs whose $a_{\cdot,s}$ is constant --- citywide and the outside option --- carry $n_s$ live keys rather than $2 n_s$. Because no pair is ever withdrawn under a bonus, $n_s$ is data: the vacant slots of Section~\ref{app:ap_variable} are needed at exactly the programs with $n_s < q_s$, and there are $\max\{0, q_s - n_s\}$ of them, independent of the zoning.

Split the assignment variables to match:
\begin{equation}\label{eq:ap_split}
D_{i,s} = D^{\mathrm{in}}_{i,s} + D^{\mathrm{out}}_{i,s},
\qquad
0 \le D^{\mathrm{in}}_{i,s} \le a_{i,s},
\qquad
0 \le D^{\mathrm{out}}_{i,s} \le 1 - a_{i,s} .
\end{equation}

\noindent\textbf{Lemma 2 (Tier split exactness).} \textit{At every binary $a$ and every point satisfying \eqref{eq:ap_split}, $D^{\mathrm{in}}_{i,s} = a_{i,s} D_{i,s}$ and $D^{\mathrm{out}}_{i,s} = (1 - a_{i,s}) D_{i,s}$, with no integrality required of $D$.}

\textit{Proof.} If $a_{i,s} = 1$ the second bound forces $D^{\mathrm{out}}_{i,s} = 0$; if $a_{i,s} = 0$ the first forces $D^{\mathrm{in}}_{i,s} = 0$. $\blacksquare$

Lemma~2 is the whole point of the split. Under a bonus the shaft of a comb is a sum over an index set that \emph{depends on the decision} --- who is above $i$ at $s$ is a function of $a$ --- and such a sum is bilinear in $(a, D)$. The split linearizes it exactly, because the products $a_{i,s} D_{i,s}$ that the bilinear form needs are now variables in their own right. Define, for a base pair $(i,s) \in \Gamma$ and a tier $z \in \{\mathrm{in}, \mathrm{out}\}$, the \emph{tier shaft}
\begin{equation}\label{eq:ap_tier_shaft}
\Sigma^{z}_{s,i}
= \sum_{(i',s) \in \Gamma}
\Big(
\mathds{1}\big[\kappa^{\mathrm{in}}_{i',s} \ge \kappa^{z}_{i,s}\big]\, D^{\mathrm{in}}_{i',s}
\;+\;
\mathds{1}\big[\kappa^{\mathrm{out}}_{i',s} \ge \kappa^{z}_{i,s}\big]\, D^{\mathrm{out}}_{i',s}
\Big),
\end{equation}
a fixed linear form whose indicators are data: each is evaluated once from \eqref{eq:ap_keys}. Write $\succ_s^{\psi,a}$ for the order the realized scores induce and $S^{\psi,a}(s,i)$ for the shaft \eqref{eq:ap_shaft_tooth} taken under it. By Lemma~2, if $i$'s realized tier is $z$ then $\Sigma^{z}_{s,i} = D(S^{\psi,a}(s,i))$, and the two forms $\Sigma^{\mathrm{in}}_{s,i}$, $\Sigma^{\mathrm{out}}_{s,i}$ together cover both possibilities. The identity is exact and not a relaxation: $\Sigma^{z}_{s,i}$ sums \emph{both} copies of every applicant with a key above position $\mathrm{pos}_s(i,z)$, and Lemma~2 zeroes the copy the zoning did not realize. In particular $\Sigma^{\mathrm{out}}_{s,i}$ contains $D^{\mathrm{in}}_{i,s}$, because $\kappa^{\mathrm{in}}_{i,s} > \kappa^{\mathrm{out}}_{i,s}$; that term vanishes at every zoning where the row's right-hand side is active.

Do not instantiate \eqref{eq:ap_tier_shaft} as written. Each form has $O(n_s)$ nonzeros and there are $2 n_s$ of them per program, which costs $O(\sum_s n_s^2)$ nonzeros. Because the forms are nested prefixes of one order, introduce instead $2 n_s + 1$ prefix variables per program,
\begin{equation}\label{eq:ap_prefix}
E_{s,0} = 0,
\qquad
E_{s,m} = E_{s,m-1} + D^{z_s(m)}_{i_s(m),\, s}
\qquad m = 1, \dots, 2 n_s,
\end{equation}
so that $E_{s,m}$ is the mass at $s$ held by the first $m$ positions of the merged order, and read the shafts off it:
\begin{equation}\label{eq:ap_shaft_prefix}
\Sigma^{z}_{s,i} = E_{s,\, \mathrm{pos}_s(i,z)},
\qquad\text{and the strict shaft is}\qquad
E_{s,\, \mathrm{pos}_s(i,z) - 1} .
\end{equation}
This is $2|\Gamma|$ extra columns and $2|\Gamma|$ extra rows in total, and it leaves every comb row with $O(q_s)$ nonzeros. Do the same on the applicant side, where the bonus changes nothing: let
\begin{equation}\label{eq:ap_tooth_prefix}
t_{i,s} = \sum_{s' \succ_i^\psi s} D_{i,s'},
\qquad D_{i,s'} = D^{\mathrm{in}}_{i,s'} + D^{\mathrm{out}}_{i,s'},
\end{equation}
be the $|\Gamma|$ tooth prefix variables of \eqref{eq:ap_shaft_tooth}. Writing $s^{(1)} \succ_i^\psi \cdots \succ_i^\psi s^{(R_i)}$ for $i$'s list, they are defined by the recursion $t_{i,s^{(1)}} = 0$ and $t_{i, s^{(r)}} = t_{i, s^{(r-1)}} + D_{i, s^{(r-1)}}$, so that $T^\psi(i,s)$ carries mass $t_{i,s} + D_{i,s}$.

That leaves the eligibility of the teeth, which is where the two keys reappear. A candidate tooth base $k$ can stand above $i$ at some zoning only if its zoned key does, so the admissible set is
\begin{equation}\label{eq:ap_bonus_teeth}
\mathcal{K}^{z}(s,i) = \big\{ k \neq i : (k,s) \in \Gamma, \ \mathrm{pos}_s(k, \mathrm{in}) < \mathrm{pos}_s(i,z) \big\} \cup \{\varnothing\},
\end{equation}
and the credit records whether $k$ stands above $i$ unconditionally or only by virtue of its own bonus --- one integer comparison per tooth:
\begin{equation}\label{eq:ap_bonus_credit}
\chi^{z}_{i,k,s} =
\begin{cases}
0 & \text{if } \mathrm{pos}_s(k, \mathrm{out}) < \mathrm{pos}_s(i,z) \quad \text{($k$ ahead in either tier),}\\
1 - a_{k,s} & \text{if } \mathrm{pos}_s(k, \mathrm{out}) > \mathrm{pos}_s(i,z) \quad \text{($k$ ahead only when zoned),}\\
1 & \text{if } k = \varnothing .
\end{cases}
\end{equation}
Collecting everything, the priority-bonus recourse program for sample $\psi$ is
\begin{align}
\max \quad & \sum_{(i,s) \in \Gamma} u_{i,s} \big( D^{\mathrm{in}}_{i,s} + D^{\mathrm{out}}_{i,s} \big) \label{pb:obj} \\
\textrm{s.t.} \quad
& \sum_{s : (i,s) \in \Gamma} \big( D^{\mathrm{in}}_{i,s} + D^{\mathrm{out}}_{i,s} \big) = 1 && \forall i \in \mathcal{I} \label{pb:assign} \\
& \sum_{i : (i,s) \in \Gamma} \big( D^{\mathrm{in}}_{i,s} + D^{\mathrm{out}}_{i,s} \big) \le q_s && \forall s \in \mathcal{S} \label{pb:cap} \\
& D^{\mathrm{in}}_{i,s} \le a_{i,s}, \qquad D^{\mathrm{out}}_{i,s} \le 1 - a_{i,s} && \forall (i,s) \in \Gamma \label{pb:tier} \\
& E_{s,0} = 0, \qquad E_{s,m} = E_{s,m-1} + D^{z_s(m)}_{i_s(m),\, s} && \forall s, \ m \le 2 n_s \label{pb:prefix} \\
& t_{i,s^{(1)}} = 0, \qquad t_{i,s^{(r)}} = t_{i,s^{(r-1)}} + D_{i,s^{(r-1)}} && \forall i, \ 2 \le r \le R_i \label{pb:tooth} \\
& E_{s,\, \mathrm{pos}_s(i,z)} \;+\; t_{i,s} \;+\; \sum_{k \in K} \big[\, t_{k,s} + \chi^{z}_{i,k,s} \,\big]
\ \ge\ q_s\, \alpha^{z}_{i,s} && \forall (i,s), z, K \label{eq:ap_bonus_comb} \\
& D^{\mathrm{in}}_{i,s}, \ D^{\mathrm{out}}_{i,s} \ge 0 && \forall (i,s) \in \Gamma \label{pb:domain}
\end{align}
where \eqref{eq:ap_bonus_comb} ranges over $(i,s) \in \Gamma$, both tiers $z$, and every collection $K$ of $q_s - 1$ elements of $\mathcal{K}^{z}(s,i)$ admissible in the sense of Section~\ref{app:ap_combs}, and where
\begin{equation}\label{eq:ap_alpha}
\alpha^{\mathrm{in}}_{i,s} = a_{i,s},
\qquad
\alpha^{\mathrm{out}}_{i,s} = 1 - a_{i,s}
\end{equation}
switch each row on precisely at the zonings where its tier is the realized one. Four points of practice. The access constraint \eqref{saa:access} is gone: under a bonus no pair is ever withdrawn, and the tier bounds \eqref{pb:tier} carry all of the zoning's influence. Those bounds are written as rows rather than as variable bounds for the same reason \eqref{saa:access} is in \texttt{SaaOracle}: their multipliers are what the master cut of Section~\ref{sec:saa} reads. Every left-hand side is written in the prefix variables \eqref{eq:ap_prefix} and \eqref{eq:ap_tooth_prefix}, so a comb row has $q_s + 1$ nonzeros plus the credits, whatever $n_s$ is. And Corollary~1 transfers verbatim, since at a fixed zoning \eqref{pb:obj}--\eqref{pb:domain} is an ordinary stable admissions market: solved comb-complete, its value is the welfare of deferred acceptance run under the bonus-induced priorities.

As a warm start, the aggregated surrogate of \eqref{saa:rothblum} transfers by substituting the strict shaft of \eqref{eq:ap_shaft_prefix},
\begin{equation}\label{pb:rothblum}
E_{s,\, \mathrm{pos}_s(i,z) - 1}
\;+\; q_s \big( t_{i,s} + D_{i,s} \big)
\;\ge\; q_s\, \alpha^{z}_{i,s}
\qquad \forall (i,s) \in \Gamma, \ z \in \{\mathrm{in}, \mathrm{out}\},
\end{equation}
which is $2|\Gamma|$ rows and is valid because at the active tier it is exactly \eqref{saa:rothblum} for the realized market, the unrealized copies inside the prefix being zero by Lemma~2. Counting columns, $2|\Gamma|$ assignment variables, $2|\Gamma|$ shaft prefixes and $|\Gamma|$ tooth prefixes give $5|\Gamma| = 61{,}850$ on $\text{Block}_2$; counting rows, $|\mathcal{I}| + |\mathcal{S}| + 5|\Gamma| \approx 6.6 \times 10^4$ before a single comb is separated. Lemma~2 fixes one of each variable pair to zero at any binary $a$, so an oracle solved at a fixed candidate loses half the assignment columns in presolve.

\noindent\textbf{Proposition 11 (Priority-bonus polytope).} \textit{Fix $\psi$ and $\beta \ge 1$. The system \eqref{pb:obj}--\eqref{pb:domain} is valid at every binary $a$ simultaneously, its integral points are exactly the stable matchings under the bonus-induced order, and at each fixed binary $a$ its solution set is their convex hull.}

\textit{Proof.} Fix a binary $a$ and let $\succ_s^{\psi,a}$ be the induced strict order. Exactly one tier $z(i)$ has $\alpha^{z}_{i,s} = 1$; rows of the other tier have right-hand side zero and non-negative left-hand side. For the active tier, Lemma~2 gives $\Sigma^{z(i)}_{s,i} = D(S^{\psi,a}(s,i))$, and \eqref{eq:ap_bonus_credit} gives $\chi^{z(i)}_{i,k,s} = 0$ exactly when $k \succ_s^{\psi,a} i$ and $\chi^{z(i)}_{i,k,s} = 1$ otherwise, since a tooth with $\kappa^{\mathrm{out}}_{k,s} < \kappa^{z}_{i,s} \le \kappa^{\mathrm{in}}_{k,s}$ outranks $i$ if and only if $a_{k,s} = 1$. So the active rows are precisely the inequalities \eqref{saa:comb} for the ordinary stable admissions market $(\Gamma, \succ_s^{\psi,a})$ in which every pair is acceptable, and Proposition~10 applies to that market verbatim, with credited and vacant teeth playing the role they played there. $\blacksquare$

Three limits of $\beta$ are worth reading off, because they place the two formulations of this appendix on one axis.
\begin{enumerate}
    \item $\beta = 0$. The two keys coincide, so $\Sigma^{\mathrm{in}} = \Sigma^{\mathrm{out}}$ and every credit vanishes, the second branch of \eqref{eq:ap_bonus_credit} becoming vacuous. The two tiered rows at a base then have identical left-hand sides and right-hand sides $q_s a_{i,s}$ and $q_s(1 - a_{i,s})$, so at every binary $a$ they jointly assert only that the left-hand side is at least $q_s$ --- the zoning-free inequality. The zoning has left the formulation entirely, as it should: with no bonus and no access restriction, welfare does not depend on $x$.
    \item $\beta > \max_{i,s} \rho_{i,s}$, a dominant bonus. Every zoned applicant outranks every unzoned one, so $\succ_s^{\psi,a}$ is the base order restricted to the zone, concatenated with the base order restricted to its complement, and \eqref{eq:ap_tier_shaft} collapses to $\Sigma^{\mathrm{in}}_{s,i} = \sum_{i' \succeq^{0}_s i} D^{\mathrm{in}}_{i',s}$ and $\Sigma^{\mathrm{out}}_{s,i} = \sum_{i'} D^{\mathrm{in}}_{i',s} + \sum_{i' \succeq^{0}_s i} D^{\mathrm{out}}_{i',s}$, where $\succeq^{0}_s$ is the base order: an unzoned applicant is behind the whole zoned block.
    \item Hard access. Add $D^{\mathrm{out}}_{i,s} = 0$ and $\alpha^{\mathrm{out}}_{i,s} = 0$ for every non-citywide $s$, which declares the unzoned copy \emph{unacceptable} rather than merely low-priority. Then \eqref{eq:ap_split} reduces to $D_{i,s} \le a_{i,s}$, which is \eqref{saa:access}; with a dominant bonus every tooth satisfies $\kappa^{\mathrm{out}}_{k,s} < \kappa^{\mathrm{in}}_{i,s}$, so \eqref{eq:ap_bonus_credit} reduces to $\chi_{k,s} = 1 - a_{k,s}$, which is \eqref{eq:ap_credit}; and \eqref{eq:ap_bonus_comb} becomes \eqref{saa:comb}. The access polytope of Section~\ref{app:ap_combs} is the face of the priority-bonus polytope on which the unzoned copies vanish.
\end{enumerate}

\subsection{Comb Separation Under a Priority Bonus}
\label{app:ap_bonus_sep}

At a candidate zoning the separation problem is not merely similar to the one under access, it is the same problem. Once $a$ is fixed, only the tier with $\alpha^{z}_{i,s} = 1$ need be separated, since the other has right-hand side zero, and every credit collapses: for a base $(i,s)$ and a tooth $k$ realized above $i$, either $k$ is realized unzoned, in which case $\mathrm{pos}_s(k,\mathrm{out}) < \mathrm{pos}_s(i,z)$ and \eqref{eq:ap_bonus_credit} gives $0$, or $k$ is realized zoned, in which case both branches give $0$ because $a_{k,s} = 1$. Conversely an admissible tooth that is \emph{not} realized above $i$ has $a_{k,s} = 0$ and credit $1$, so its bracket is at least $1$ and a vacant slot is never worse. The most violated comb may therefore be sought over the realized competitors alone, with base-independent brackets
\begin{equation}\label{eq:ap_bonus_bracket}
\bar b_{k,s} = \bar t_{k,s} \quad\text{(the tooth prefix \eqref{eq:ap_tooth_prefix} evaluated at $\bar D$),}
\end{equation}
and the descent of Section~\ref{app:ap_sep} applies verbatim: walk the merged order at $s$, skipping the $n_s$ unrealized positions; carry $E_{s,\cdot}$ as the running shaft; carry a max-heap of size $q_s - 1$ over \eqref{eq:ap_bonus_bracket}, seeded with $q_s - 1$ vacant slots. The cost is again $O(|\Gamma| \log \max_s q_s)$.

What changes is the \emph{lifting}. The descent sees one order, but the row it emits has to hold at every zoning, and the two are reconciled by recording the row's combinatorics rather than the order that produced it. Three items must be lifted, and each corresponds to a way the row would otherwise be tight at $x^r$ and invalid elsewhere.
\begin{itemize}
    \item \textbf{The base tier.} Record which of $a_{i,s}$ or $1 - a_{i,s}$ is the right-hand side multiplier. A row emitted with the constant $q_s$ would remain in force at zonings that move $i$ to the other tier, where its shaft is the wrong prefix.
    \item \textbf{The shaft.} Record it as the single prefix variable $E_{s,\, \mathrm{pos}_s(i,z)}$ of \eqref{eq:ap_shaft_prefix}. Expanding the realized prefix into coefficients on the aggregate $D_{i',s}$ would freeze the current order into the row; the prefix formulation removes the temptation, since the shaft is one column reference and the merged order lives in \eqref{pb:prefix} where it belongs.
    \item \textbf{The tooth credits.} For each selected tooth, compare $\mathrm{pos}_s(k,\mathrm{out})$ against $\mathrm{pos}_s(i,z)$ and store the constant $0$ or the affine term $1 - a_{k,s}$ accordingly. This is the one place where the collapse noted above is a trap rather than a convenience: at the candidate the two are numerically indistinguishable, both evaluating to $0$, so a row that stores the evaluated credit looks correct and is not. Recorded as $0$, a tooth that outranks $i$ only when zoned yields a row violated by legitimate stable matchings at any zoning that unzones $k$ --- the failure mode of the two-applicant example in Section~\ref{app:ap_variable}.
\end{itemize}
The practical consequence is that a base pair contributes two row families and that at any one candidate only one of them is separated, so the per-iteration row count is what it was under hard access; both families accumulate over iterations as the master moves applicants between tiers.

\subsection{Consequences for the Master Cuts}
\label{app:ap_consequences}

Two consequences of this appendix bear on the formulations in the body.

First, the sign of the cut coefficients. The coefficient of $a_{i,s}$ in the dual cut of Section~\ref{sec:saa} collects the multiplier of \eqref{saa:access}, which is non-negative, together with $q_s$ times the multipliers of the stability rows based at $(i,s)$ --- and, once combs are separated, $-1$ times the multiplier of every comb row in which $i$ appears as a \emph{credited tooth}, since $a_{i,s}$ enters those rows through $\chi_{i,s} = 1 - a_{i,s}$ on the left. With only the aggregated rows \eqref{saa:rothblum} loaded, $a_{i,s}$ occurs in exactly one row and on one side, and the vertex-pair coefficient $\Gamma^{(r,\psi)}_{v,u}$ inherits a single sign. With combs, or with the tiered rows \eqref{eq:ap_bonus_comb}, it does not. The sparse instantiation of the master must therefore key on $\Gamma^{(r,\psi)}_{v,u} \neq 0$, as $\mathcal{A}^*$ in \eqref{saa:master_access} does, and the reading of $\beta^{(r,\psi)}_{i,s}$ as a marginal \emph{value} of access is specific to the aggregated relaxation.

Second, and more consequentially, a priority bonus does not admit the congestion-priced bound developed in Appendix~\ref{app:priced_access}. That construction starts from $T(x)$, the transportation value obtained by keeping the access constraints \eqref{saa:access} and deleting stability. Under a bonus there is no access constraint to keep. Every zoning has $\Gamma(a) = \Gamma$, so $T(x)$ collapses to the constant $\overline{W}^{\mathrm{TP}}$, the bound $\overline{W}^{\mathrm{PA}}(x; \pi)$ stops depending on $x$, and its maximizer is arbitrary. The reason is structural rather than technical. Under hard access the zoning acts on welfare through the \emph{feasible set}, which a transportation relaxation still sees; under a priority bonus it acts on welfare exclusively through \emph{who wins a contested seat}, which lives entirely in the stability constraints. Any relaxation that discards stability discards the whole dependence on the decision. The same argument disposes of the monotonicity that Proposition~6 relies on: granting $i$ the bonus at $s$ weakly helps $i$ and weakly harms every applicant that $i$ thereby overtakes, so welfare is not monotone in $a$ and the threshold family \eqref{eq:pa_threshold_cut} has no counterpart. Of the three budget-set formulations, the comb-based recourse is the only one whose dependence on the zoning is of the right kind to price a priority bonus.

\section{Welfare Bounds for the Budget-Set Masters}
\label{app:welfare_bounds}

This appendix develops the upper bounds that the SAA master of Section~\ref{sec:saa} can put in the epigraph domain \eqref{saa:master_bounds}. They come in three layers, each pricing something the previous one ignored. Section~\ref{app:transport_bound} gives zoning-independent constants that price capacity. Section~\ref{app:zoned_transport} keeps a constant but prices the confinement too, by maximizing the zone-restricted transportation value over a linear relaxation of the zoning model itself. Section~\ref{app:priced_access} abandons the constant altogether for a \emph{function} of the zoning; it is the congestion-priced bound the \texttt{priced\_access} strategy optimizes.

\subsection{The A Priori Transportation Bound}
\label{app:transport_bound}

The constant $\text{Welfare}_{\max}$ in \eqref{saa:master_bounds} is merely the domain of the epigraph variables $\eta_\psi$, but it carries more weight than that role suggests. The optimality gap reported by the master is measured against its dual bound, and that dual bound is the smaller of $\text{Welfare}_{\max}$ and whatever the accumulated cuts \eqref{saa:master_cut} certify. Until the cuts descend below the constant, the constant \emph{is} the reported gap.

The natural choice, obtained by granting every applicant their top-ranked program,
\begin{equation}\label{eq:welfare_fc}
\overline{W}^{\mathrm{FC}} = \sum_{i \in \mathcal{I}} u_{i, 1},
\end{equation}
ignores capacity entirely and is therefore loose by the full extent of oversubscription. We replace it with a bound that respects capacity while remaining independent of the zoning.

\noindent\textbf{Proposition 4 (Transportation Bound).} \textit{Let}
\begin{align}
\overline{W}^{\mathrm{TP}} = \max_{y} \quad & \sum_{i \in \mathcal{I}} \sum_{s \in \mathcal{S}} u_{i,s}\, y_{i,s} \label{eq:welfare_tp_obj} \\
\textrm{s.t.} \quad & \sum_{s \in \mathcal{S}} y_{i,s} \leq 1 && \forall i \in \mathcal{I} \label{eq:welfare_tp_assign} \\
& \sum_{i \in \mathcal{I}} y_{i,s} \leq q_s && \forall s \in \mathcal{S} \label{eq:welfare_tp_cap} \\
& y_{i,s} \geq 0 && \forall i, s. \label{eq:welfare_tp_domain}
\end{align}
\textit{Then for every zoning $x \in \Omega(G, C, k, \delta)$ and every tie-breaking seed $\psi \in \Psi$, the stable-matching welfare $W_\psi(x)$ satisfies}
\begin{equation}\label{eq:welfare_chain}
W_\psi(x) \;\leq\; \overline{W}^{\mathrm{TP}} \;\leq\; \overline{W}^{\mathrm{FC}}.
\end{equation}

\textit{Proof.} Let $D$ be the stable matching induced by $(x, \psi)$, and let $s_0 \in \mathcal{S}$ denote the outside option, with $u_{i, s_0} = 0$ and capacity $q_{s_0} = |\mathcal{I}|$. Take $y_{i,s} = D_{i,s}$ for $s \neq s_0$. Constraint \eqref{eq:welfare_tp_assign} follows from \eqref{saa:assignment} after discarding the mass on $s_0$, constraint \eqref{eq:welfare_tp_cap} is \eqref{saa:capacity}, and \eqref{eq:welfare_tp_domain} is \eqref{saa:domain}; since $u_{i,s_0} = 0$, the objective is unchanged. Hence $y$ is feasible for \eqref{eq:welfare_tp_obj}--\eqref{eq:welfare_tp_domain} and $W_\psi(x) \leq \overline{W}^{\mathrm{TP}}$. Equivalently, \eqref{eq:welfare_tp_obj}--\eqref{eq:welfare_tp_domain} is the sample program \eqref{saa:primal_obj}--\eqref{saa:domain} with the access constraints \eqref{saa:access} and the stability constraints \eqref{saa:comb} and \eqref{saa:rothblum} deleted, so it is a relaxation of every sample, uniformly in $(x, \psi)$. Deleting \eqref{eq:welfare_tp_cap} in turn decouples the students and leaves $\sum_i \max_s u_{i,s} = \overline{W}^{\mathrm{FC}}$, giving the second inequality, strictly whenever some program is oversubscribed at its top-choice demand. $\blacksquare$

Because \eqref{eq:welfare_chain} holds pointwise in $\psi$, it also bounds the sample average, and so $\text{Welfare}_{\max} = \overline{W}^{\mathrm{TP}}$ is a valid domain for $\eta$.

Three properties make this substitution attractive. First, \eqref{eq:welfare_tp_obj}--\eqref{eq:welfare_tp_domain} is a transportation problem, so its polytope is integral and the LP optimum is attained by an integral assignment; no rounding argument is needed. Second, it carries one variable per ranked student-program pair --- $1.24 \times 10^4$ on our instances --- and solves in milliseconds, which is negligible beside a single master solve, so the tightening is free. Third, the substitution is weakly dominant: since the reported bound is $\min\{\text{Welfare}_{\max}, \text{cut bound}\}$ and $\overline{W}^{\mathrm{TP}} \leq \overline{W}^{\mathrm{FC}}$, replacing the constant can never worsen the reported bound.

On the $\text{Block}_2$ instances ($|\mathcal{I}| = 3953$ applicants, $|\mathcal{S}| = 123$ programs, $k = 6$ zones), $\overline{W}^{\mathrm{FC}} = 19{,}311.29$ against $\overline{W}^{\mathrm{TP}} = 18{,}699.31$ for one centroid configuration and $18{,}968.60$ against $18{,}370.41$ for a second, reductions of $611.98$ and $598.19$. The improvement actually realized in the reported bound lies between zero and that reduction, according to how far the cuts have already descended: we observe $573.19$ and $67.76$ respectively, the second case being one in which the cuts had already certified a value close to $\overline{W}^{\mathrm{TP}}$.

We emphasize what this does not address. Neither constant prices the zoning restriction itself: $\overline{W}^{\mathrm{TP}}$ still permits every applicant a globally best available seat, whereas a feasible zoning confines each applicant to the programs co-zoned with their residence, roughly a $1/k$ share of the district. That unpriced restriction, rather than any weakness of the sample programs, accounts for the bulk of the residual gap, and closing it requires a relaxation that reasons about zone membership. Two such relaxations follow. Section~\ref{app:zoned_transport} keeps a zoning-independent constant but computes it over a linear relaxation of the zoning model, so the constant sees the confinement; Section~\ref{app:priced_access} keeps the zoning integral and replaces the constant with a function of it.

\subsection{Zoned Transportation Bounds}
\label{app:zoned_transport}

Proposition 4 obtains its constant by deleting the access constraints \eqref{saa:access} outright, which is exactly why it cannot see the confinement. This subsection keeps those constraints and pays for the zoning model instead, at the price of relaxing that model to a linear program: geographic units may sit fractionally in several zones at once. The optimum is still a \emph{constant}, so nothing in the master of \eqref{saa:master_obj}--\eqref{saa:master_domain} changes --- but it is a constant computed over the zoning polytope rather than in ignorance of it.

\noindent\textbf{Zone-Restricted Transportation Value:}
Reinstating the access constraints \eqref{saa:access} in the transportation problem of Proposition 4 while still deleting the stability constraints \eqref{saa:comb} and \eqref{saa:rothblum} gives, for each zoning $x$,
\begin{align}
T(x) = \max_{y} \quad & \sum_{i \in \mathcal{I}} \sum_{s \in \mathcal{S}} u_{i,s}\, y_{i,s} \label{eq:pa_restricted_obj} \\
\textrm{s.t.} \quad & \sum_{s \in \mathcal{S}} y_{i,s} \leq 1 && \forall i \in \mathcal{I} \label{eq:pa_restricted_assign} \\
& \sum_{i \in \mathcal{I}} y_{i,s} \leq q_s && \forall s \in \mathcal{S} \label{eq:pa_restricted_cap} \\
& 0 \leq y_{i,s} \leq a_{i,s}(x) && \forall i \in \mathcal{I}, s \in \mathcal{S}. \label{eq:pa_restricted_access}
\end{align}

\noindent\textbf{Proposition 14 (Restricted Transportation Bound).} \textit{For every zoning $x \in \Omega(G, C, k, \delta)$ and every seed $\psi \in \Psi$, $W_\psi(x) \leq T(x) \leq \overline{W}^{\mathrm{TP}}$.}

\textit{Proof.} The stable matching induced by $(x, \psi)$ satisfies \eqref{saa:access}, so the feasible point constructed in the proof of Proposition 4 also satisfies \eqref{eq:pa_restricted_access}; and $a_{i,s}(x) \leq 1$ makes \eqref{eq:pa_restricted_access} a restriction of \eqref{eq:welfare_tp_domain}. $\blacksquare$

Both subsections that follow start from $T(x)$. This one bounds $\max_x T(x)$ by relaxing the partition; Section~\ref{app:priced_access} bounds $T(x)$ pointwise by relaxing capacity.

\noindent\textbf{The Relaxed Zoning Polytope:}
Write $Z(v) \subseteq [k]$ for the labels vertex $v$ may take, as induced by \texttt{max\_distance}, the centroid anchors and their graph-hop neighborhoods, and any fixed assignment. Relaxing the integrality of the base model \eqref{cp:obj}--\eqref{cp:b_domain} while retaining every row gives the polytope $\Omega^{\mathrm{LP}}(G, C, k, \delta)$ in the variables $x_{v,z}$ and $b_{i,j}$:
\begin{align}
& \sum_{z \in Z(v)} x_{v,z} = 1, \qquad 0 \leq x_{v,z} \leq 1, \qquad x_{v,z} = 0 \ \ \forall z \notin Z(v) && \forall v \in V \label{eq:zt_assign} \\
& x_{c(z), z} = 1 && \forall z \in [k] \label{eq:zt_anchor} \\
& \left\lfloor \tfrac{|S|}{k} \right\rfloor \leq \sum_{v \in V} s(v)\, x_{v,z} \leq \left\lceil \tfrac{|S|}{k} \right\rceil && \forall z \in [k] \label{eq:zt_schools} \\
& \sum_{v \in V} \big( \gamma_v - \lambda^-\, n_v \big) x_{v,z} \geq 0, \qquad \sum_{v \in V} \big( \gamma_v - \lambda^+\, n_v \big) x_{v,z} \leq 0 && \forall z \in [k] \label{eq:zt_ratio} \\
& b_{i,j} \geq x_{i,z} - x_{j,z}, \quad b_{i,j} \geq x_{j,z} - x_{i,z}, \quad \sum_{(i,j) \in E} b_{i,j} \leq \lfloor \delta_\partial |E| \rfloor && \forall (i,j) \in E, z \in [k] \label{eq:zt_boundary}
\end{align}
Row \eqref{eq:zt_ratio} is the zero-right-hand-side form of \eqref{cp:capacity}, \eqref{cp:frl} and the racial rows of Appendix~\ref{app:diversity_constraints}: one pair per balanced quantity, with $\gamma_v$ the per-vertex value of that quantity ($c_v$, $f_v$, or $r_{v,e}$) and $\lambda^\mp$ its lower and upper ratio bound. Each side is retained exactly when the corresponding tolerance is active and dropped when it is disabled, so a run with the capacity and racial tolerances off carries only the FRL pair. Row \eqref{eq:zt_boundary} is the cut-edge budget the budget-set masters impose in place of minimizing the boundary. Contiguity is written in one of two ways, which is the only difference between the two bounds:
\begin{description}
    \item[(CN)] the closer-neighbor rows of \eqref{cp:contiguity} relaxed to $x_{v,z} \leq \sum_{u \in \text{cn}(v, c(z))} x_{u,z}$ for $v \neq c(z)$, with $x_{v,z} = 0$ where the support set is empty;
    \item[(F)] a rooted single-commodity flow, one commodity per zone with the root fixed at $c(z)$. Replace each undirected edge by both arcs, let $m = |V|$, and let $\varphi^z_{uv} \geq 0$ be the flow on arc $(u,v)$ for label $z$:
\end{description}
\begin{align}
& \sum_{u \in N(v)} \big( \varphi^z_{uv} - \varphi^z_{vu} \big) = x_{v,z}
&& \forall z \in [k],\ v \neq c(z) \label{eq:zt_flow_balance} \\
& \sum_{u \in N(c(z))} \big( \varphi^z_{c(z)u} - \varphi^z_{uc(z)} \big)
= \sum_{v \neq c(z)} x_{v,z}
&& \forall z \in [k] \label{eq:zt_flow_supply} \\
& 0 \leq \varphi^z_{uv} \leq m\, x_{u,z}, \qquad \varphi^z_{uv} \leq m\, x_{v,z}
&& \forall z \in [k],\ (u,v) \label{eq:zt_flow_capacity}
\end{align}
At integral $x$ this enforces exactly that $G[Z_z]$ is connected and contains $c(z)$: a selected component not containing the root would have positive total demand \eqref{eq:zt_flow_balance} and no incoming flow, while a connected selected set can route one unit to every non-root member along a rooted spanning tree. Because the root is fixed rather than chosen, no root-selection variables are needed. The flow description is \emph{weaker} than (CN) at integral $x$ --- it permits connected zones that (CN) rejects for lack of monotone support toward the anchor --- which is why $\widehat{W}^{\mathrm{ZT}}_{\mathrm{flow}} \geq \widehat{W}^{\mathrm{ZT}}_{\mathrm{cn}}$ is the direction observed empirically.

\noindent\textbf{Reified Access:}
The co-zoning indicators are linearized on the relaxed memberships exactly as the masters linearize \eqref{saa:master_access}, using an auxiliary $h_{v,u,z}$ for the conjunction:
\begin{equation}\label{eq:zt_access}
0 \leq h_{v,u,z} \leq x_{v,z}, \qquad h_{v,u,z} \leq x_{u,z}, \qquad 0 \leq a_{v,u} \leq \sum_{z \in Z(v) \cap Z(u)} h_{v,u,z} .
\end{equation}
At binary $x$ this forces $a_{v,u} = 1$ precisely when $v$ and $u$ are co-zoned, and permits it, which is all that is needed below.

\noindent\textbf{The Bound:}
For $\bullet \in \{\mathrm{cn}, \mathrm{flow}\}$ define
\begin{align}
\widehat{W}^{\mathrm{ZT}}_{\bullet} = \max \quad & \sum_{i \in \mathcal{I}} \sum_{s \in \mathcal{S}} u_{i,s}\, y_{i,s} \label{eq:zt_obj} \\
\textrm{s.t.} \quad & x \in \Omega^{\mathrm{LP}}(G, C, k, \delta) \ \text{with contiguity } \bullet, \ \text{and } \eqref{eq:zt_access} \label{eq:zt_omega} \\
& \sum_{s \in \mathcal{S}} y_{i,s} \leq 1 && \forall i \in \mathcal{I} \label{eq:zt_assign_row} \\
& \sum_{i \in \mathcal{I}} y_{i,s} \leq q_s && \forall s \in \mathcal{S} \label{eq:zt_cap} \\
& 0 \leq y_{i,s} \leq a_{v(i), l(s)} && \forall i, \ s \notin F_i \label{eq:zt_access_row} \\
& 0 \leq y_{i,s} \leq 1 && \forall i, \ s \in F_i \label{eq:zt_free}
\end{align}
where $F_i \subseteq \mathcal{S}$ collects the programs whose access no zoning can revoke --- citywide programs, and programs with $l(s) = v(i)$ --- for which $a_{i,s} \equiv 1$ and no access variable exists. Pairs with $Z(v(i)) \cap Z(l(s)) = \emptyset$ carry $y_{i,s} = 0$, since no zoning can ever co-zone them.

\noindent\textbf{Proposition 13 (Zoned Transportation Bound).} \textit{For $\bullet \in \{\mathrm{cn}, \mathrm{flow}\}$, every zoning $x \in \Omega(G, C, k, \delta)$ and every tie-breaking seed $\psi \in \Psi$,}
\begin{equation}\label{eq:zt_chain}
W_\psi(x) \;\leq\; T(x) \;\leq\; \widehat{W}^{\mathrm{ZT}}_{\bullet} \;\leq\; \overline{W}^{\mathrm{TP}} ,
\qquad\text{hence}\qquad
\max_{x \in \Omega} \ \frac{1}{|\Psi|} \sum_{\psi \in \Psi} W_\psi(x) \;\leq\; \widehat{W}^{\mathrm{ZT}}_{\bullet} ,
\end{equation}
\textit{so either value is a valid $\text{Welfare}_{\max}$ for \eqref{saa:master_bounds}.}

\textit{Proof.} The first inequality is Proposition 14. For the third, every feasible point of \eqref{eq:zt_obj}--\eqref{eq:zt_free} has $y \geq 0$ and satisfies \eqref{eq:zt_assign_row} and \eqref{eq:zt_cap}, which are \eqref{eq:welfare_tp_assign} and \eqref{eq:welfare_tp_cap}; so its $y$ is feasible for \eqref{eq:welfare_tp_obj}--\eqref{eq:welfare_tp_domain} with the same objective, and the transportation optimum is at least as large. Conversely, deleting the zoning block \eqref{eq:zt_omega} and the access rows \eqref{eq:zt_access_row} leaves exactly \eqref{eq:welfare_tp_obj}--\eqref{eq:welfare_tp_domain}, so nothing but the access rows separates the two.

For the second inequality, fix $x \in \Omega$ and exhibit a feasible point of \eqref{eq:zt_obj}--\eqref{eq:zt_free} whose objective is $T(x)$. Rows \eqref{eq:zt_assign}--\eqref{eq:zt_boundary} and (CN) are the rows of the base model, which $x$ satisfies integrally, with $b_{i,j}$ set to its induced value. For (F): closer-neighbor feasibility induces a strictly distance-decreasing path from every vertex of $Z_z$ to $c(z)$ inside $Z_z$, so $G[Z_z]$ is connected and contains $c(z)$; routing one unit from $c(z)$ to each other vertex of $Z_z$ along a spanning tree of $G[Z_z]$ gives arc flows in $[0, |V|]$, zero on every arc with an endpoint outside $Z_z$, so \eqref{eq:zt_flow_balance}--\eqref{eq:zt_flow_capacity} hold. Taking $h_{v,u,z} = x_{v,z} x_{u,z}$ satisfies \eqref{eq:zt_access} and yields $a_{v,u} = a_{v,u}(x)$, so any $y$ optimal for \eqref{eq:pa_restricted_obj}--\eqref{eq:pa_restricted_access} is feasible here and $T(x) \leq \widehat{W}^{\mathrm{ZT}}_{\bullet}$. The final claim holds because \eqref{eq:zt_chain} is pointwise in $(x, \psi)$, so it bounds the sample average and then its maximum over $\Omega$. $\blacksquare$

Three remarks on the construction. First, the two variants are both valid but are \emph{not} nested: (CN) relaxes rows the master itself carries and is therefore the natural companion to it, while (F) is an exact description of connectedness at binary $x$ and rules out no $(v,z)$ pair a priori, so it contains every (CN)-feasible integral zoning --- but at fractional $x$ neither set of rows implies the other. We report both rather than choose. Second, the bound is a constant, so unlike the priced bound of Section~\ref{app:priced_access} it prices confinement only in the aggregate: it says what the best \emph{fractional} zoning could achieve, not what the master's current candidate achieves. Third, the value depends only on the feasible set $\Omega$, the market, and which contiguity description is used --- never on how many threads solved it --- which is why the implementation caches it content-addressed by exactly those inputs with the worker count excluded from the key. The LP must be solved to optimality for \eqref{eq:zt_chain} to mean anything: a time-limited LP returns a primal-feasible point, which bounds the value from below.

\noindent\textbf{Remark (Where the tightening comes from).}
It does not come from the contiguity rows. Suppose every unanchored vertex is a candidate for all $k$ labels and the school-count and ratio bounds admit the district averages. Then $\bar{x}$, placing $\bar{x}_{v,z} = 1/k$ at every unanchored $v$, is feasible for both variants: each zone receives a $1/k$ share of every quantity, so \eqref{eq:zt_schools} and \eqref{eq:zt_ratio} hold at the district average; each row of (CN) reads $1/k \leq |\text{cn}(v, c(z))|/k$, satisfied whenever the support is nonempty; (F) routes $1/k$ to each vertex along any spanning tree; and every unanchored vertex carries the same membership vector, so $b_{i,j} = 0$ on every edge between two of them and \eqref{eq:zt_boundary} is slack. At $\bar{x}$ we may take $h_{v,u,z} = 1/k$ for all $z$ and obtain $a_{v,u} = 1$ for \emph{every} pair with two unanchored endpoints: the same pair can be co-zoned a little in every zone at once. On such an instance $\widehat{W}^{\mathrm{ZT}}_\bullet$ can differ from $\overline{W}^{\mathrm{TP}}$ only through the anchors. The bound therefore earns whatever it earns from what pins zone membership --- the anchors and their neighborhoods, the candidate sets $Z(v)$ induced by \texttt{max\_distance}, and vertices whose closer-neighbor support is a single vertex --- and an instance run with \texttt{max\_distance} $= \infty$ should be expected to return $\overline{W}^{\mathrm{TP}}$ almost exactly. This is the same mechanism that made the triangle and cardinality inequalities on $a$ worthless: the all-access point is inside the relaxation, and no symmetric valid inequality in $a$-space removes it.

Since the SFUSD instances are run with a finite \texttt{max\_distance}, the bound does earn something. On $\text{Block}_2$ with $k = 6$, $\texttt{max\_distance} = 3.1$ miles, $\delta_\partial = 0.18$, $\delta_F = 0.15$ and the racial and capacity tolerances disabled:

\begin{center}
\begin{tabular}{lrrrr}
\hline
centroid set & $\overline{W}^{\mathrm{FC}}$ & $\overline{W}^{\mathrm{TP}}$ & $\widehat{W}^{\mathrm{ZT}}_{\mathrm{cn}}$ & $\widehat{W}^{\mathrm{ZT}}_{\mathrm{flow}}$ \\
\hline
6-zone-3 & 18{,}948.93 & 18{,}355.58 & \textbf{17{,}114.93} & 17{,}187.56 \\
6-zone-9 & 18{,}673.33 & 18{,}079.53 & \textbf{16{,}881.51} & 16{,}945.91 \\
\hline
\end{tabular}
\end{center}

The zoned bounds cut $1{,}240.65$ and $1{,}198.02$ off $\overline{W}^{\mathrm{TP}}$ under (CN), and $1{,}168.02$ and $1{,}133.63$ under (F) --- roughly a third of the standing gap, with (CN) tighter than (F) on both instances as the non-nesting remark allows but does not guarantee. The cost is one LP: $2.1 \times 10^4$ columns and $3.2 \times 10^4$ rows under (CN), $2.8 \times 10^4$ and $4.6 \times 10^4$ under (F), solved single-threaded in $1.8$ to $2.8$ seconds under (CN) and $2.3$ to $2.4$ under (F), against a master solve of $300\,\mathrm{s}$; model construction accounts for under $0.3\,\mathrm{s}$ of that and a cache hit for $14\,\mathrm{ms}$. Threads do not help: sweeping one to eight, no count above one was reliably faster and several were clearly slower, since the solver races LP algorithms concurrently and on a model this size they mostly contend --- (F) reaches $4.4\,\mathrm{s}$ at eight threads. The implementation therefore gives this LP a single thread regardless of how many the master gets, less to save the two seconds than to leave the cores to the master. Since (CN) dominates on both tightness and time, it is the default, and (F) is reported for the comparison rather than for use. Validity was checked directly at two integral zonings of the 6-zone-3 instance, where $T(x)$ came back at $14{,}493.94$ and $13{,}331.82$, far below both bounds. Replacing the constant in a two-iteration probe at 6-zone-3 moved the reported absolute gap from $4{,}262.54$ to $2{,}887.29$, which is the constant reduction plus a small incumbent difference; a single probe says nothing about whether changing the epigraph domain helps the search, and we make no claim there.

What this does not do is escape the convexification gap. The bound still sits far above anything any method attains --- on the order of $2{,}000$ units, judging by the budget welfares of Section~\ref{sec:saa}, though those were measured on a differently parameterized instance of the same family and so are indicative rather than a matched comparison --- and the reason is visible in the remark: a fractional membership buys a pro-rata slice of a popular school that no integral zoning could hand to the same applicants, and \eqref{eq:zt_cap} limits only the total. That is a property of relaxing the partition, not of this particular objective, and it is why the values above land in the same range as the root relaxations of the fully joint formulations we have measured on this instance family. The contribution is not a tighter class of bound; it is that a value of that quality is available \emph{a priori}, in seconds, as a constant the master can be handed, instead of being something the master has to spend its root solve discovering.

\noindent\textbf{Relation to the Priced Bound:}
$\widehat{W}^{\mathrm{ZT}}_\bullet$ and the priced bound $\overline{W}^{\mathrm{PA}}(x^\ast; \pi)$ certified by Theorem 2 are both valid upper bounds on $\max_{x \in \Omega} \bar{W}_\Psi(x)$, and neither dominates the other: the former keeps capacity exactly and relaxes the partition, the latter keeps the partition integral and relaxes capacity into prices. The implementation therefore reports the smaller of the two, which is valid because both are. The zoned LP also supplies prices: the duals of \eqref{eq:zt_cap} are non-negative, so Proposition 5 applies to them verbatim, and they price a seat by its marginal value \emph{after} confinement rather than before it. The dominance clause of Proposition 5, that $\overline{W}^{\mathrm{PA}}(x; \pi^\star) \leq \overline{W}^{\mathrm{TP}}$, is specific to the duals of the unrestricted problem and is not inherited by these; validity is, and validity is what the bound needs.

\subsection{Congestion-Priced Access}
\label{app:priced_access}

Section~\ref{app:zoned_transport} priced the confinement while keeping $\text{Welfare}_{\max}$ a constant, by relaxing the partition. This subsection takes the other route suggested at the end of Section~\ref{app:transport_bound}: keep the zoning integral and replace the \emph{constant} with a \emph{function} of it. The result is still a valid upper bound on stable-matching welfare, still separable across applicants, and --- unlike the dual cuts \eqref{saa:master_cut} --- admits a finite closed-form outer description requiring no recourse LP at all.

Recall from Section~\ref{app:zoned_transport} the zone-restricted transportation value $T(x)$ of \eqref{eq:pa_restricted_obj}--\eqref{eq:pa_restricted_access}, and that $W_\psi(x) \leq T(x) \leq \overline{W}^{\mathrm{TP}}$ for every seed. $T(x)$ is the tightest bound obtainable without reasoning about stability, but it is the optimal value of an LP parameterized by $x$; used directly it would reintroduce exactly the per-candidate recourse solve that makes \eqref{saa:master_cut} expensive, and it would again be summarized by one dual hyperplane per seed over all $|\mathcal{I}| \times |\mathcal{S}|$ access coordinates. Lagrangian relaxation of the capacity constraints removes both objections at once.

\noindent\textbf{Priced Access Bound:}
Fix a price vector $\pi \in \mathbb{R}^{|\mathcal{S}|}_{\geq 0}$ and define the \emph{priced value} of program $s$ to applicant $i$ as
\begin{equation}\label{eq:pa_priced_value}
w_{i,s} = u_{i,s} - \pi_s .
\end{equation}
With $F_i$ as in Section~\ref{app:zoned_transport}, write
\begin{equation}\label{eq:pa_student_term}
g_i(a_{i, \cdot}) = \max\Big\{0, \ \max_{s \in \mathcal{S}} w_{i,s}\, a_{i,s}\Big\}, \qquad
\phi_i = \max\Big\{0, \ \max_{s \in F_i} w_{i,s}\Big\},
\end{equation}
so that $g_i$ is the best priced value an applicant can reach under $a_{i,\cdot}$, the outside option being available at value $0$, and $\phi_i \leq g_i(a_{i,\cdot})$ is the part of it that is guaranteed. The priced-access bound is
\begin{equation}\label{eq:pa_bound}
\overline{W}^{\mathrm{PA}}(x; \pi) = \sum_{s \in \mathcal{S}} q_s \pi_s + \sum_{i \in \mathcal{I}} g_i\big(a_{i, \cdot}(x)\big).
\end{equation}

\noindent\textbf{Proposition 5 (Validity and Dominance).} \textit{For every price vector $\pi \geq 0$, every zoning $x \in \Omega(G, C, k, \delta)$ and every tie-breaking seed $\psi \in \Psi$,}
\begin{equation}\label{eq:pa_chain}
W_\psi(x) \;\leq\; T(x) \;\leq\; \overline{W}^{\mathrm{PA}}(x; \pi),
\qquad\text{and}\qquad
\min_{\pi \geq 0} \overline{W}^{\mathrm{PA}}(x; \pi) = T(x).
\end{equation}
\textit{Moreover, if $\pi^\star$ denotes an optimal dual solution to the capacity constraints \eqref{eq:welfare_tp_cap} of the unrestricted problem, then $\overline{W}^{\mathrm{PA}}(x; \pi^\star) \leq \overline{W}^{\mathrm{TP}}$ for every $x$.}

\textit{Proof.} The first inequality is Proposition 14. For the second, relax \eqref{eq:pa_restricted_cap} into the objective with multipliers $\pi \geq 0$; weak duality gives, for any such $\pi$,
$$T(x) \;\leq\; \sum_{s} q_s \pi_s + \max\Big\{ \textstyle\sum_{i,s} w_{i,s} y_{i,s} \;:\; \sum_s y_{i,s} \leq 1, \ 0 \leq y_{i,s} \leq a_{i,s}(x) \Big\}.$$
The remaining program has no constraint coupling distinct applicants, so it separates into $|\mathcal{I}|$ independent problems. Each is a maximization of a linear function over $\{y_{i, \cdot} \geq 0 : \sum_s y_{i,s} \leq 1, \ y_{i,s} \leq a_{i,s}\}$, whose optimum places unit mass on a single accessible program of largest priced value when that value is positive, and no mass at all otherwise; its value is therefore exactly $g_i(a_{i, \cdot}(x))$, establishing \eqref{eq:pa_chain}. Since \eqref{eq:pa_restricted_obj}--\eqref{eq:pa_restricted_access} is a feasible bounded linear program, strong duality applies and the Lagrangian dual attains $T(x)$, giving the stated minimum. Finally, each $g_i$ is nondecreasing in $a_{i, \cdot}$ and $a_{i,s}(x) \leq 1$, so $\overline{W}^{\mathrm{PA}}(x; \pi^\star) \leq \overline{W}^{\mathrm{PA}}(\mathbf{1}; \pi^\star)$; the latter is the Lagrangian dual objective of the unrestricted transportation problem evaluated at its own optimal multipliers, which by strong duality equals $\overline{W}^{\mathrm{TP}}$. $\blacksquare$

Three consequences are worth separating. Validity holds for \emph{every} non-negative $\pi$, so the prices are a free modelling choice rather than a quantity that must be computed correctly; taking $\pi = 0$ recovers the zone-restricted first-choice bound, in which every applicant receives their favourite co-zoned program irrespective of capacity. Dominance means the substitution $\text{Welfare}_{\max} \leftarrow \overline{W}^{\mathrm{PA}}(x; \pi^\star)$ can never weaken the bound of \eqref{saa:master_bounds}, exactly as $\overline{W}^{\mathrm{TP}}$ could not weaken $\overline{W}^{\mathrm{FC}}$. And the price of using a single fixed $\pi$ instead of re-optimizing per candidate is precisely $\overline{W}^{\mathrm{PA}}(x; \pi) - T(x)$, which by the last part of \eqref{eq:pa_chain} vanishes at the restricted problem's own duals; we quantify it below.

\noindent\textbf{Proposition 6 (Exact Threshold Description).} \textit{For each applicant $i$ let}
$$\Theta_i = \{\phi_i\} \cup \{\, w_{i,s} \;:\; s \in \mathcal{S} \setminus F_i, \ w_{i,s} > \phi_i \,\}.$$
\textit{Then for every $\tau \geq \phi_i$ and every $a_{i, \cdot} \in \{0,1\}^{|\mathcal{S}|}$,}
\begin{equation}\label{eq:pa_threshold_cut}
g_i(a_{i, \cdot}) \;\leq\; \tau + \sum_{\substack{s \in \mathcal{S} \setminus F_i \\ w_{i,s} > \tau}} (w_{i,s} - \tau)\, a_{i,s},
\end{equation}
\textit{and the family is exact on the binary cube:}
\begin{equation}\label{eq:pa_threshold_exact}
g_i(a_{i, \cdot}) \;=\; \min_{\tau \in \Theta_i} \Bigg[ \tau + \sum_{\substack{s \in \mathcal{S} \setminus F_i \\ w_{i,s} > \tau}} (w_{i,s} - \tau)\, a_{i,s} \Bigg].
\end{equation}

\textit{Proof.} Write $m = g_i(a_{i, \cdot})$ and let $\tau \geq \phi_i$. If $m \leq \tau$ the right-hand side of \eqref{eq:pa_threshold_cut} is at least $\tau \geq m$. If $m > \tau$ then $m$ is attained by some program $s^\ast$ with $a_{i, s^\ast} = 1$ and $w_{i, s^\ast} = m$; since $m > \tau \geq \phi_i$, the definition of $\phi_i$ forces $s^\ast \notin F_i$, so $s^\ast$ contributes the term $(m - \tau) a_{i, s^\ast} = m - \tau$ and the right-hand side is at least $m$. This proves \eqref{eq:pa_threshold_cut}, hence that the minimum in \eqref{eq:pa_threshold_exact} is at least $m$.

For the reverse, choose $\tau = m$ if $m > \phi_i$ and $\tau = \phi_i$ otherwise; in both cases $\tau \in \Theta_i$, because $m > \phi_i$ implies $m = w_{i, s^\ast}$ for some $s^\ast \in \mathcal{S} \setminus F_i$ as above. Every accessible program satisfies $w_{i,s} \leq m = \tau$ by maximality of $m$, so no surviving term has $a_{i,s} = 1$, the sum vanishes, and the bracket equals $\tau = m$. $\blacksquare$

Proposition 6 is the structural payoff. The function $g_i$ is piecewise linear, concave and monotone in $a_{i, \cdot}$, and \eqref{eq:pa_threshold_exact} exhibits it as the minimum of at most one more affine function than the number of programs $i$ ranks with coefficients available in closed form from the sorted priced values --- no LP, no dual solution, and no dependence on the seed $\psi$. Each inequality \eqref{eq:pa_threshold_cut} touches only applicant $i$'s own access row, which is what makes the family exact rather than merely valid: the cut set describes $g_i$ on the whole cube, whereas a dual cut \eqref{saa:master_cut} is one hyperplane, tangent at a single point of a $|\mathcal{I}| \times |\mathcal{S}|$-dimensional domain, and the multicut master of Section~\ref{sec:saa} buys $|\Psi|$ of them per iteration rather than a description. The restriction $\tau \geq \phi_i$ is not cosmetic. A threshold below $\phi_i$ would leave the guaranteed value $\phi_i$ uncovered by any term of \eqref{eq:pa_threshold_cut}, since the programs realizing it lie in $F_i$ and carry no access variable, and the inequality would then cut away welfare the applicant retains under every zoning.

\noindent\textbf{Vertex Aggregation:}
Because access is geographic, $a_{i,s} \equiv a_{v(i), l(s)}$, so applicants sharing a home vertex share an access row and their inequalities may be summed with no loss --- the same exact aggregation invoked for \eqref{saa:master_cut}, but applied to a family that is already per-applicant rather than to a dual hyperplane. Introduce per-vertex epigraph variables $\theta_v \in \mathbb{R}$ and, for any assignment of thresholds $\tau_i \geq \phi_i$ to the applicants of $\mathcal{I}_v$, define
\begin{equation}\label{eq:pa_vertex_coeff}
c_v(\tau) = \sum_{i \in \mathcal{I}_v} \tau_i,
\qquad
\Lambda_{v, u}(\tau) = \sum_{i \in \mathcal{I}_v} \ \sum_{\substack{s \in \mathcal{S}_u \setminus F_i \\ w_{i,s} > \tau_i}} (w_{i,s} - \tau_i).
\end{equation}
Summing \eqref{eq:pa_threshold_cut} over $\mathcal{I}_v$ gives the vertex cut
\begin{equation}\label{eq:pa_vertex_cut}
\theta_v \;\leq\; c_v(\tau) + \sum_{u \in S} \Lambda_{v, u}(\tau)\, a_{v, u} \qquad \forall v \in V,
\end{equation}
valid for every zoning and every admissible $\tau$. A vertex with no resident applicants has $\mathcal{I}_v = \emptyset$, hence $c_v = 0$ and $\Lambda_{v, \cdot} = 0$; together with $\theta_v \geq 0$ this pins $\theta_v = 0$ exactly, which matters because the master instantiates one epigraph variable per vertex and an uncut variable would otherwise be free up to its domain.

Note the exact structural correspondence with the choice-set cut \eqref{eq:benders_cut}: $c_v(\tau)$ plays the role of the incumbent utility $U_v(S_v^t)$ and $\Lambda_{v,u}(\tau)$ the role of the marginal gains $\Delta^+_{t,s,v}$, while the losses $\Delta^-_{s,v}$ are represented not by a separate term but by the choice of a lower threshold. Consequently the same CP-SAT and MIP master solvers serve all three formulations unchanged.

\noindent\textbf{Master Problem and Separation:}
The priced-access master problem is
\begin{align}
\max_{x, a, \theta} \quad & \sum_{v \in V} \theta_v \label{pa:master_obj} \\
\textrm{s.t.} \quad & x \in \Omega(G, C, k, \delta) \label{pa:master_omega} \\
& \theta_v \leq c_v(\tau^{(r)}) + \sum_{u \in S} \Lambda_{v, u}(\tau^{(r)})\, a_{v, u} && \forall v \in V, \ r = 1, \dots, R \label{pa:master_cut} \\
& a_{v, u} \iff \bigvee_{z \in [k]} (x_{v, z} \land x_{u, z}) && \forall (v, u) \in \mathcal{A}^{\mathrm{PA}} \label{pa:master_access} \\
& 0 \leq \theta_v \leq \textstyle\sum_{i \in \mathcal{I}_v} \max\{0, \max_s w_{i,s}\} && \forall v \in V, \label{pa:master_bounds}
\end{align}
where $\mathcal{A}^{\mathrm{PA}} = \{(v, u) \in V \times S \mid \exists r, \Lambda_{v, u}(\tau^{(r)}) \neq 0\}$ is instantiated on demand exactly as $\mathcal{A}^*$ is in \eqref{saa:master_access}. The constant $\sum_s q_s \pi_s$ is added to the objective when its value is read as a welfare bound; being independent of $x$, it does not affect the maximizer. Given a candidate zoning $x^r$, separation requires no optimization: for each applicant, sort the accessible priced values and emit the thresholds
\begin{equation}\label{eq:pa_levels}
\tau_i^{(r, \ell)} = \max\Big\{\phi_i, \ \text{the } (\ell+1)\text{-th largest } w_{i,s} \text{ over } \{s : a_{i,s}(x^r) = 1\}\Big\}
\end{equation}
for $\ell = 0, \dots, L-1$, taking the outside value $0$ once an applicant's accessible options are exhausted. Level $\ell = 0$ is tight at $x^r$ by the exactness argument of Proposition 6; levels $\ell \geq 1$ price what the applicant loses when their $\ell$ best accessible programs are withdrawn, which is the direction a family generated only at the incumbent cannot see. All $L$ levels are aggregated by \eqref{eq:pa_vertex_coeff} into $L$ cuts per vertex per iteration, at a cost of one pass over the preference lists.

\noindent\textbf{Theorem 2 (Finite Convergence to the Bound).} \textit{Suppose each master problem \eqref{pa:master_obj}--\eqref{pa:master_bounds} is solved to optimality and duplicate cuts are discarded. Then the priced-access cutting-plane algorithm terminates finitely at a zoning $x^\ast$ maximizing $\overline{W}^{\mathrm{PA}}(\cdot\,; \pi)$ over $\Omega(G, C, k, \delta)$, and}
\begin{equation}\label{eq:pa_certificate}
\max_{x \in \Omega} \ \frac{1}{|\Psi|} \sum_{\psi \in \Psi} W_\psi(x) \;\leq\; \overline{W}^{\mathrm{PA}}(x^\ast; \pi).
\end{equation}

\textit{Proof.} By Proposition 6 every cut \eqref{eq:pa_vertex_cut} is valid for all $x \in \Omega$, so the master maximizes an outer approximation of $\sum_v \sum_{i \in \mathcal{I}_v} g_i$ and its optimum weakly exceeds $\max_x \overline{W}^{\mathrm{PA}}(x; \pi) - \sum_s q_s \pi_s$ at every iteration. The threshold set $\Theta_i$ is finite and \eqref{eq:pa_levels} draws from it, so only finitely many distinct cuts \eqref{eq:pa_vertex_cut} exist; discarding duplicates therefore admits finitely many iterations that add a cut. At an iteration adding none, the level-$0$ cuts at $x^r$ are already present, and by exactness they force $\theta_v = \sum_{i \in \mathcal{I}_v} g_i(a_{i, \cdot}(x^r))$; the master objective thus equals the true value of $\sum_i g_i$ at $x^r$ while dominating it everywhere on $\Omega$, so $x^r$ is a maximizer. Inequality \eqref{eq:pa_certificate} follows because Proposition 5 holds pointwise in $\psi$, hence for the sample average, and $x^\ast$ maximizes the dominating function. $\blacksquare$

Two limitations deserve emphasis, since \eqref{eq:pa_certificate} is a statement about the bound and not about the matching. First, $\overline{W}^{\mathrm{PA}}$ strictly exceeds sample-average stable-matching welfare, so its maximizer need not maximize welfare; the algorithm returns a certified optimum of a surrogate, and the usefulness of that optimum rests on the surrogate's bias being close to constant across the zonings the master visits. Second, Theorem 2 presumes exact master solves, whereas a time-limited CP-SAT solve returns an incumbent; under a per-iteration limit the method retains validity of \eqref{eq:pa_certificate} for the cuts generated but forfeits the optimality claim, and terminates on the iteration budget rather than on the absence of separation.

\noindent\textbf{Decomposition of the Bias:}
Proposition 5 splits the residual optimism into two independently interpretable pieces,
\begin{equation}\label{eq:pa_bias}
\underbrace{\overline{W}^{\mathrm{PA}}(x; \pi^\star) - T(x)}_{\text{fixed rather than re-optimized prices}}
\;+\;
\underbrace{T(x) - \tfrac{1}{|\Psi|}\textstyle\sum_\psi W_\psi(x)}_{\text{stability discarded}} .
\end{equation}
On the $\text{Block}_2$ instance with $k = 6$ zones and $|\Psi| = 15$ seeds, evaluated at the final zonings of five different strategies and with $\pi^\star$ taken from the single unrestricted transportation LP of Proposition 4, which gives $\overline{W}^{\mathrm{TP}} = 18{,}699.31$ for this centroid configuration. Writing $\bar{W}_\Psi(x) = \frac{1}{|\Psi|}\sum_{\psi} W_\psi(x)$ for sample-average stable-matching welfare:

\begin{center}
\begin{tabular}{lrrrrr}
\hline
zoning from & $\overline{W}^{\mathrm{PA}}$ & $T(x)$ & $\bar{W}_\Psi(x)$ & prices & stability \\
\hline
SAA               & 15{,}084.29 & 14{,}997.76 & 14{,}822.77 &  86.53 & 175.00 \\
short bursts      & 14{,}735.40 & 14{,}647.44 & 14{,}486.66 &  87.96 & 160.78 \\
choice set        & 14{,}635.02 & 14{,}486.79 & 14{,}324.86 & 148.23 & 161.94 \\
MID decomposition & 14{,}220.26 & 14{,}098.75 & 13{,}884.10 & 121.50 & 214.65 \\
MID               & 14{,}203.24 & 14{,}011.49 & 13{,}849.74 & 191.75 & 161.75 \\
\hline
\end{tabular}
\end{center}

The last two columns are the two terms of \eqref{eq:pa_bias}. The zoning-independent constant $\overline{W}^{\mathrm{TP}}$ overstates sample-average welfare by $3{,}877$ to $4{,}850$ depending on the zoning evaluated; $\overline{W}^{\mathrm{PA}}$ overstates it by $249$ to $354$, of which $87$ to $192$ is attributable to holding $\pi$ fixed and $161$ to $215$ to discarding stability. Two features of the table matter more than its absolute magnitudes. The bound is valid at every one of the five zonings and at every individual seed, not merely on average, as Proposition 5 requires. And the overstatement is nearly constant across zonings produced by five structurally unrelated methods and spanning almost a thousand units of welfare, which is the empirical condition under which maximizing $\overline{W}^{\mathrm{PA}}$ approximates maximizing welfare --- a regularity we observe rather than a property we can prove, and the reason the limitations noted after Theorem 2 are stated as limitations rather than dismissed.

\noindent\textbf{Role of the Prices:}
Setting $\pi = 0$ isolates what congestion pricing contributes. The resulting surrogate --- each applicant's favourite co-zoned program, capacity ignored --- is cheaper to optimize and in our runs converges faster, terminating on the absence of separation with a model-versus-realized gap of $0.01\%$; but it converges to a materially worse zoning, because an unpriced surrogate is indifferent between a program with spare seats and an oversubscribed one and is therefore rewarded for concentrating demand on popular schools. Charging $\pi^\star$ recovered $1{,}357$ units of realized welfare relative to $\pi = 0$ on one instance at a matched iteration allowance --- a single-instance observation, but a difference an order of magnitude larger than the bias figures tabulated above. This is the clearest illustration of the gap between Theorem 2 and the objective one actually cares about: the algorithm certifies optimality for whichever surrogate it is given, and the surrogate must be the right one.

\section{Dantzig--Wolfe Decomposition over Whole Zones}
\label{app:dw_mid}
Every formulation so far has described a zoning by its assignment variables and
handed the whole thing to one solver. This appendix describes the alternative:
take an entire zone as the atom of decision, let the master choose one zone per
label so that the zones tile the district, and generate the zones themselves on
demand from the master's shadow prices. What makes the exchange worthwhile is
that a single zone's welfare is a well-posed object --- the welfare of a
self-contained submarket --- so a column carries an exactly evaluated value
rather than a surrogate, and the decomposition inherits whichever welfare
definition we care to price.

Three things have to hold for that to work, and they organize the appendix.
Section~\ref{app:dw_families} shows that the base feasible set
$\Omega(G, C, k, \delta)$ of Section~2 factors label by label, so that ``legal
zone for label $z$'' is a well-defined family $\mathcal F_z$ and the master's
rows are exactly what is left over. Section~\ref{app:dw_value} shows that zone
welfare is additive over a partition, under each of the three definitions we
price, and that the sampled stable-matching definition is computable by running
deferred acceptance rather than by solving the access polytope of
Appendix~\ref{app:access_polytope}. Section~\ref{app:dw_pricing} writes the
pricing problem for one anchored zone and Section~\ref{app:dw_bounds} gives the
bound and explains why it survives an arbitrary dual point. That section also
reports where the method actually stalls, which is neither the cost of pricing
nor the choice of dual point but the \emph{completability} of the columns
generated; Section~\ref{app:dw_redraw} answers it with a second, bound-free
column source that re-partitions two zones at once and shows why adding one is
sound while substituting it for pricing would not be.
Section~\ref{app:dw_convergence} closes the integrality gap by branching.

The decomposition requires that no program offer citywide access. The
restriction is substantive rather than cosmetic: a citywide program's seats are
contested by applicants in every zone through one shared capacity row, so zone
welfare would not be additive and a column would have no well-defined value. We
therefore remove citywide schools and their programs from the market and from
the school and capacity counts. Applicants remain at their residential
vertices; their lists retain the order, utilities and priority tiers of the
remaining programs, and an applicant with no remaining acceptable program
receives the outside option of utility zero.

\subsection{Admissible Zones for One Label}
\label{app:dw_families}

Write $\mathcal C(v) \subseteq [k]$ for the labels vertex $v$ may take, as
determined by $\texttt{max\_distance}$, explicit candidate restrictions and
fixed assignments; $\text{cn}(v, c(z))$ for the graph neighbours of $v$ whose
geometry is strictly closer to label $z$'s school point than $v$'s own, the
relation already used in \eqref{cp:contiguity}; and $B_r(c)$ for the set of
vertices within $r$ graph hops of $c$, with $r$ the configured centroid
neighbourhood radius. Define the \emph{support}
\begin{equation}\label{eq:dw_support}
\text{cn}^{*}(v, z) = \big\{ u \in \text{cn}(v, c(z)) : z \in \mathcal C(u),
\ u = c(z) \text{ or } \text{cn}^{*}(u, z) \neq \emptyset \big\},
\end{equation}
the greatest fixed point of the displayed recursion --- a neighbour counts as
support only if it is itself a candidate and can itself continue toward the
anchor. A vertex with $\text{cn}^{*}(v,z) = \emptyset$ can never carry label
$z$; the base model reaches the same conclusion by propagation from
\eqref{cp:contiguity}, so nothing is added or removed by taking the fixed point
explicitly.

For each label $z \in [k]$ let $\mathcal F_z$ be the family of $A \subseteq V$
satisfying
\begin{align}
& c_z \in A \quad\text{and}\quad B_r(c_z) \subseteq A, \label{eq:dw_anchor} \\
& A \cap B_r(c_{z'}) = \emptyset \qquad \forall z' \neq z, \label{eq:dw_exclude} \\
& z \in \mathcal C(v) \qquad \forall v \in A, \label{eq:dw_candidacy} \\
& A \cap \text{cn}^{*}(v, z) \neq \emptyset \qquad \forall v \in A \setminus \{c_z\}, \label{eq:dw_support_row} \\
& \underline{s} \leq \textstyle\sum_{v \in A} s(v) \leq \overline{s},
\qquad \text{and \eqref{cp:capacity}, \eqref{cp:frl} and their per-ethnicity
analogues, evaluated on } A. \label{eq:dw_local}
\end{align}
Condition \eqref{eq:dw_exclude} subsumes the exclusion of every other centroid,
since $c_{z'} \in B_r(c_{z'})$. Conditions \eqref{eq:dw_anchor} and
\eqref{eq:dw_local} are exactly \eqref{cp:centroid} and
\eqref{cp:schools}--\eqref{cp:frl} read at one label, and
\eqref{eq:dw_support_row} is \eqref{cp:contiguity} read at one label. Only
individual admissibility is required: $V \setminus A$ need not itself be an
admissible zone.

\noindent\textbf{Lemma 3 (Monotone support).} \textit{Let $A \in \mathcal F_z$.
Then every $v \in A$ is joined to $c_z$ by a path inside $A$ along which the
distance to label $z$'s school point strictly decreases; in particular $G[A]$ is
connected. The converse fails: there are connected sets containing $c_z$ and
satisfying \eqref{eq:dw_anchor}--\eqref{eq:dw_local} that violate
\eqref{eq:dw_support_row}.}

\textit{Proof.} Let $d(v)$ denote the distance from $v$'s geometry to label
$z$'s school point. If $v \in A \setminus \{c_z\}$ then
\eqref{eq:dw_support_row} supplies $u \in A$ adjacent to $v$ with $d(u) < d(v)$.
Iterating produces a strictly decreasing sequence of distances along vertices of
$A$; since $A$ is finite the sequence terminates, and it can only terminate at a
vertex carrying no obligation, i.e.\ at $c_z$. Every vertex of $A$ therefore
reaches $c_z$ inside $A$, so $G[A]$ is connected.

For the converse, take the four-cycle $v_0 - v_1 - v_2 - v_3 - v_0$ with
$c_z = v_0$ and geometry giving $d(v_0) = 0$, $d(v_1) = 2$, $d(v_2) = 1$,
$d(v_3) = 3$ --- nothing prevents $v_2$ from being geometrically nearer the
school than the vertex one has to walk through to reach it. Then
$A = \{v_0, v_1, v_2\}$ is connected and contains the anchor, but $v_2$'s only
neighbour in $A$ is $v_1$, which is farther from the school, so
\eqref{eq:dw_support_row} fails at $v_2$ and $A \notin \mathcal F_z$.
$\blacksquare$

The strictness in Lemma 3 is a modelling decision, not an oversight. It is the
contiguity description the base solvers of Section~2 already use, and adopting
it here is what makes the decomposition describe the same feasible set as the
rest of the paper. It is also what makes pricing tractable: connectedness costs
one clause per candidate vertex, where an unanchored rooted-flow description
costs $2|E|$ arc variables with big-$M$ capacities and lets one label's
subproblem range over every connected subset of the graph.

\noindent\textbf{Proposition 15 (Label separability).} \textit{A tuple
$(A_0, \dots, A_{k-1})$ with $A_z \in \mathcal F_z$ for every $z$ and
$\bigsqcup_z A_z = V$ is precisely the zone tuple of a point
$x \in \Omega(G, C, k, \delta)$, and conversely. If a boundary budget
$\overline b$ is imposed, add $\sum_z b(A_z) \le \overline b$ on the left and
the corresponding row on the right.}

\textit{Proof.} Inspect the base model
\eqref{cp:obj}--\eqref{cp:b_domain}. Every row involves the variables of a
single label --- \eqref{cp:centroid}, \eqref{cp:contiguity},
\eqref{cp:schools}, \eqref{cp:capacity}, \eqref{cp:frl} and the per-ethnicity
rows all do --- except \eqref{cp:assignment}, which states that the labels
partition $V$, and the boundary rows \eqref{cp:cut}, whose only aggregate use
is the optional cap. Fix $x \in \Omega$ and set
$A_z = \{v : x_{v,z} = 1\}$. Then \eqref{cp:assignment} gives
$\bigsqcup_z A_z = V$, and the label-$z$ rows evaluated at $x$ are literally
\eqref{eq:dw_anchor}--\eqref{eq:dw_local}, so $A_z \in \mathcal F_z$.
Conversely, given such a tuple, define $x_{v,z} = \mathds{1}\{v \in A_z\}$;
disjointness and coverage give \eqref{cp:assignment} and each family membership
gives that label's rows. $\blacksquare$

\subsection{The Value of a Zone}
\label{app:dw_value}

Let $\mathcal G_A = \{g \in \mathcal G : v(g) \in A\}$ and
$\mathcal S_A = \{s \in \mathcal S : l(s) \in A\}$. We price three definitions
of $W(A)$.

\noindent\textbf{(i) Finite-grid MID welfare.}
Fix an integer lottery scale $L \geq 1$ and restrict each cutoff to
$P_s = p_s / L$ with $p_s$ a non-negative integer, so every assignment
probability is a multiple of $1/L$. Writing
$\widehat R_{g,r} = L R_{g,r}$ and $\widehat d_{g,r} = L d_{g,r}$, the mass
recurrence \eqref{eq:mid_recurrence} becomes
\begin{align}
\widehat{t}_{s,\rho} &= \min\{L,\max\{p_s-L\rho,0\}\}, \label{eq:dw_threshold} \\
\widehat{e}_{g,r} &= a_{g,s(g,r)}\widehat{t}_{s(g,r),\rho(g,r)}
                      +(1-a_{g,s(g,r)})L, \label{eq:dw_effective} \\
\widehat{R}_{g,0} &= L, \qquad
\widehat{R}_{g,r} = \min\{\widehat{R}_{g,r-1},\widehat{e}_{g,r}\}, \label{eq:dw_recurrence} \\
\widehat{d}_{g,r} &= \widehat{R}_{g,r-1}-\widehat{R}_{g,r}. \label{eq:dw_mass}
\end{align}
For a fixed $A$ set $a_{g,s} = 1$ exactly when $g \in \mathcal G_A$ and
$s \in \mathcal S_A$; running \eqref{eq:dw_threshold}--\eqref{eq:dw_mass} over
all types and ranks then gives zero assigned mass to every omitted applicant.
The zone value is
\begin{equation}\label{eq:dw_zone_value}
\begin{aligned}
W^{\mathrm{MID}}_L(A) = \max_{p,\widehat{R},\widehat{d}}\quad
& \frac{1}{L}\sum_{g \in \mathcal{G}} n_g
       \sum_{r=1}^{R_g} u(g,r)\widehat{d}_{g,r} \\
\text{s.t.}\quad
& \eqref{eq:dw_threshold}\text{--}\eqref{eq:dw_mass}, \\
& \sum_{g \in \mathcal{G}}\sum_{r:s(g,r)=s}
        n_g\widehat{d}_{g,r} \leq Lq_s
        \quad \forall s \in \mathcal{S}, \\
& p_s \in \{0,\ldots,\overline{p}_s\}
        \quad \forall s \in \mathcal{S},
\end{aligned}
\end{equation}
where $\overline p_s = L(1 + \max_{g,r : s(g,r) = s} \rho(g,r))$, an empty
maximum read as zero. This bound permits rejecting every applicant, so
\eqref{eq:dw_zone_value} is always feasible, and all masses are integers in
$[0, L]$.

\noindent\textbf{Proposition 7 (Least-Cutoff Zone Value).} \textit{For every
$A \in \mathcal F_z$, the value $W^{\mathrm{MID}}_L(A)$ is the welfare at the
componentwise least capacity-feasible cutoff vector.}

\textit{Proof.} Start with $p = 0$. Whenever a program is overloaded, raise only
its cutoff to the smallest integer value clearing its capacity, holding the
others fixed. Demand is non-increasing in a program's own cutoff and
$\overline p_s$ clears it, so each update strictly increases a bounded integer
coordinate and the procedure terminates at a capacity-feasible $p^\ast$.

Let $\widetilde p$ be any capacity-feasible vector. Initially
$p \le \widetilde p$. If this holds before an update of program $s$, then
evaluating $s$'s demand at $p_s = \widetilde p_s$ with all other cutoffs at
their current, weakly smaller values yields demand no greater than at
$\widetilde p$: lower cutoffs at earlier choices only reduce the mass remaining
to demand $s$, and later choices do not affect it. So the smallest clearing
update satisfies $p_s \le \widetilde p_s$, and induction gives
$p^\ast \le \widetilde p$.

Set $u(g,R_g+1) = 0$. By \eqref{eq:dw_recurrence} lower cutoffs weakly decrease
every remaining mass, and for a non-empty list
\begin{equation}\label{eq:dw_welfare_telescoping}
\frac{n_g}{L}\sum_{r=1}^{R_g}u(g,r)\widehat{d}_{g,r}
= n_g u(g,1)-\frac{n_g}{L}\sum_{r=1}^{R_g}
\big(u(g,r)-u(g,r+1)\big)\widehat{R}_{g,r},
\end{equation}
whose remaining-mass coefficients are all non-positive because utilities are
rank-ordered. Hence $p^\ast$ weakly maximizes welfare among capacity-feasible
cutoff vectors; empty lists contribute zero. $\blacksquare$

\noindent\textbf{(ii) Sampled stable-matching welfare.}
Fix a single tie-breaking seed $\psi$, drawn under the same single- or
multiple-tie-breaking rule as Section~\ref{sec:saa} and after citywide programs
have been removed, so that the realized strict orders $\succ_s^\psi$ refer to
the programs that remain. Let $\Gamma(a_A)$ be the acceptable pairs of the zone
submarket, that is those $(i,s) \in \Gamma$ with $v(i) \in A$ and
$l(s) \in A$, and define
\begin{equation}\label{eq:dw_sampled_value}
W^{\psi}(A) = \max\Big\{ \textstyle\sum_{(i,s)} u_{i,s} D_{i,s}
: D \text{ a stable matching of } \Gamma(a_A) \Big\}.
\end{equation}
This is the object Appendix~\ref{app:access_polytope} describes: the maximum of
a linear utility over the stable admissions polytope
\eqref{saa:primal_obj}--\eqref{saa:domain} of the submarket.

\noindent\textbf{Proposition 16 (The sampled zone value is a deferred-acceptance
run).} \textit{For every $A$, $W^{\psi}(A)$ equals the utilitarian welfare of
the applicant-proposing deferred-acceptance outcome on $\Gamma(a_A)$.}

\textit{Proof.} This is Corollary~1 applied to the submarket. The feasible set
of \eqref{saa:primal_obj}--\eqref{saa:domain} at the binary access pattern
$a_A$ is, by Proposition~2, the convex hull of the stable matchings of
$\Gamma(a_A)$, so the maximum is attained at one of them. Deferred acceptance
returns the applicant-optimal stable matching, every applicant weakly preferring
it to every other, and $u_{i,\cdot}$ is consistent with $\succ_i^\psi$, so the
sum is weakly largest there. $\blacksquare$

Proposition 16 is the practical content of the whole access-polytope
development for this appendix, and it cuts in a direction worth naming. The
$\texttt{saa}$ master of Section~\ref{sec:saa} cannot run the mechanism,
because it is searching over zonings and the access pattern is a variable, so
it loads the aggregated stability row \eqref{saa:rothblum} and pays for the
slack: on $\text{Block}_2$ that surrogate over-reports realized
deferred-acceptance welfare by $194.8$ to $251.8$ units. Here the same quantity
is needed only at a \emph{fixed} zone, where the mechanism can simply be run.
The column's value is therefore exactly realized welfare, at the cost of one
deferred-acceptance pass, and none of the comb machinery is needed to compute
it. The inequalities come back in Section~\ref{app:dw_pricing}, where
membership is again a variable.

\noindent\textbf{(iii) Boundary cost.}
With $w_{u,v} \geq 0$ the shared boundary length, define the half-perimeter
\begin{equation}\label{eq:dw_half_perimeter}
b(A)=\frac{1}{2}\sum_{\substack{\{u,v\}\in E\\ |\{u,v\}\cap A|=1}} w_{u,v},
\qquad W^{\partial}(A) = -b(A).
\end{equation}
With $w \equiv 1$ this is half the number of edges leaving $A$. This is not a
welfare; it reproduces the base compactness objective \eqref{cp:obj} and is
included because it makes the decomposition checkable against an enumerated
optimum with no market at all.

\noindent\textbf{Proposition 17 (Additivity of zone value).} \textit{Let
$(A_0, \dots, A_{k-1})$ be an admissible partition and let
$W \in \{W^{\mathrm{MID}}_L, W^{\psi}, W^{\partial}\}$. Then the district value
under $W$'s own definition equals $\sum_{z} W(A_z)$.}

\textit{Proof.} Because no program is citywide, $\mathcal S_{A_z}$ partitions
$\mathcal S$ and $\mathcal G_{A_z}$ partitions $\mathcal G$, and an applicant
has access only to programs in their own zone. For $W^{\mathrm{MID}}_L$, both
the recurrence \eqref{eq:dw_threshold}--\eqref{eq:dw_mass} and the capacity rows
separate across zones: independent optimizers of \eqref{eq:dw_zone_value}
combine into a feasible district allocation, and every district allocation
restricts to a feasible allocation in each zone. For $W^{\psi}$, the acceptable
pairs decompose as $\Gamma(a(x)) = \bigsqcup_z \Gamma(a_{A_z})$ with disjoint
applicant and program sets, so a pair of $\Gamma(a_{A_z})$ blocks a matching
$D$ iff it blocks $D$ restricted to zone $z$; stability, and hence the
optimization \eqref{eq:dw_sampled_value}, therefore separates. For
$W^{\partial}$, each cut edge lies in exactly two zone perimeters, so
$\sum_z b(A_z)$ is the district boundary cost. $\blacksquare$

\subsection{The Zone-Column Master}
\label{app:dw_master}
Let $\overline b \geq 0$ be an optional district boundary budget and introduce a
binary $\lambda_{z,A}$ selecting zone $A$ for label $z$. The master is
\begin{align}
\max_{\lambda}\quad
& \sum_{z\in[k]}\sum_{A\in\mathcal{F}_z}W(A)\lambda_{z,A}
\label{eq:dw_master_obj} \\
\text{s.t.}\quad
& \sum_{z\in[k]}\sum_{\substack{A\in\mathcal{F}_z\\v\in A}}\lambda_{z,A}=1
&& \forall v\in V, \label{eq:dw_master_cover} \\
& \sum_{A\in\mathcal{F}_z}\lambda_{z,A}=1
&& \forall z\in[k], \label{eq:dw_master_label} \\
& \sum_{z\in[k]}\sum_{A\in\mathcal{F}_z}b(A)\lambda_{z,A}\leq\overline{b},
\label{eq:dw_master_boundary} \\
& \lambda_{z,A}\in\{0,1\}
&& \forall z\in[k],\ A\in\mathcal{F}_z. \label{eq:dw_master_binary}
\end{align}
If compactness is unrestricted, delete \eqref{eq:dw_master_boundary}.

\noindent\textbf{Proposition 8 (Exact Decomposition).} \textit{For each of the
three definitions of $W$, the integer master
\eqref{eq:dw_master_obj}--\eqref{eq:dw_master_binary} has the same optimal value
as maximizing the corresponding district objective over
$\Omega(G, C, k, \delta)$.}

\textit{Proof.} Rows \eqref{eq:dw_master_cover} and \eqref{eq:dw_master_label}
say that exactly one member of each family is selected and that the selected
sets cover each vertex exactly once, i.e.\ that they form a partition with
$A_z \in \mathcal F_z$. By Proposition 15 these are precisely the admissible
zonings, and \eqref{eq:dw_master_boundary} is the boundary cap because each cut
edge appears in two perimeters. By Proposition 17 the objective
\eqref{eq:dw_master_obj} is the district objective at that zoning. Both
feasibility and objective are therefore preserved in both directions.
$\blacksquare$

\subsection{Pricing One Anchored Zone}
\label{app:dw_pricing}
Maintain a restricted pool $\mathcal F'_z \subseteq \mathcal F_z$ and relax
\eqref{eq:dw_master_binary} to $\lambda \geq 0$. Let $\alpha_v$ and $\gamma_z$
be the unrestricted duals of \eqref{eq:dw_master_cover} and
\eqref{eq:dw_master_label} and $\mu \geq 0$ the multiplier of
\eqref{eq:dw_master_boundary}; the complete dual requires
\begin{equation}\label{eq:dw_dual_row}
\sum_{v\in A}\alpha_v+\gamma_z+\mu b(A)\geq W(A)
\qquad \forall z\in[k],\ A\in\mathcal{F}_z,
\end{equation}
with objective $\sum_v \alpha_v + \sum_z \gamma_z + \mu \overline b$. Take
$\mu = 0$ and drop the budget term when there is no boundary constraint.
Pricing label $z$ maximizes the \emph{reduced welfare}
\begin{equation}\label{eq:dw_pricing}
f_z(A) \;=\; W(A)-\sum_{v\in A}\alpha_v-\gamma_z-\mu b(A),
\qquad A \in \mathcal F_z .
\end{equation}
A strictly positive optimum identifies an improving column; non-positive optima
for all labels certify that the restricted LP solves the full LP relaxation.

To solve \eqref{eq:dw_pricing} globally, introduce binary memberships
$y_v = \mathds{1}\{v \in A\}$ for the vertices with $z \in \mathcal C(v)$ and
$\text{cn}^{*}(v,z) \neq \emptyset$, and write
\eqref{eq:dw_anchor}--\eqref{eq:dw_local} in them. Conditions
\eqref{eq:dw_anchor}--\eqref{eq:dw_candidacy} fix variables. Condition
\eqref{eq:dw_support_row} is the clause family
\begin{equation}\label{eq:dw_support_clause}
y_v \;\leq\; \sum_{u \in \text{cn}^{*}(v,z)} y_u
\qquad \forall v \neq c_z ,
\end{equation}
one row per candidate vertex, and \eqref{eq:dw_local} is the same inequality
with each sum over $A$ replaced by its membership-weighted sum. By Lemma 3 no
connectivity variables are required: \eqref{eq:dw_support_clause} already
implies connectedness, and it is the description the base model uses. For the
boundary term, $\beta_{uv} = |y_u - y_v|$ is imposed by
\begin{equation}\label{eq:dw_edge_linearization}
\begin{aligned}
\beta_{uv}&\geq y_u-y_v, &\quad \beta_{uv}&\geq y_v-y_u,\\
\beta_{uv}&\leq y_u+y_v, &\quad \beta_{uv}&\leq 2-y_u-y_v,
\end{aligned}
\end{equation}
and is created only when $\mu$ can be nonzero or $W = W^{\partial}$. An edge
with exactly one candidate endpoint needs no variable: it is cut precisely when
that endpoint joins.

\noindent\textbf{Access, with exactly one zone.}
The zoning enters a welfare block only through the access indicators, and with
a single zone
\begin{equation}\label{eq:dw_pricing_access}
a_{i,s} \;=\; y_{v(i)}\, y_{l(s)}, \qquad
a_{i,s}\leq y_{v(i)}, \quad a_{i,s}\leq y_{l(s)}, \quad
a_{i,s}\geq y_{v(i)}+y_{l(s)}-1 ,
\end{equation}
one variable and three rows per distinct vertex pair; if $v(i) = l(s)$ then
$a_{i,s} = y_{v(i)}$. This is the structural saving relative to a
whole-district master, which needs one conjunction variable per label plus the
row $\sum_z a^{z}_{i,s} \le 1$ to reify the disjunction over labels
\eqref{saa:master_access}. It can be pushed further.

\noindent\textbf{Proposition 18 (Single-zone access elimination).}
\textit{Suppose the access variables of a welfare block appear only (a) as upper
bounds $D \leq a_{i,s}$ on non-negative quantities and (b) on the right-hand
side of rows $\Lambda \geq \theta\, a_{i,s}$ with $\theta \geq 0$ and
$\Lambda \geq 0$. Replacing (a) by the pair $D \leq y_{v(i)}$,
$D \leq y_{l(s)}$ and (b) by
$\Lambda \geq \theta\,(y_{v(i)} + y_{l(s)} - 1)$ leaves the feasible set
unchanged at every integral $y$, and eliminates every access variable.}

\textit{Proof.} Fix integral $y$. If $y_{v(i)} = y_{l(s)} = 1$ then
$a_{i,s} = 1$, and both substitutions reproduce the original row exactly, since
$\min\{y_{v(i)}, y_{l(s)}\} = 1$ and $y_{v(i)} + y_{l(s)} - 1 = 1$. Otherwise
$a_{i,s} = 0$; the original (a) forces $D = 0$, which the pair does too because
$\min\{y_{v(i)}, y_{l(s)}\} = 0$, and the original (b) reads $\Lambda \geq 0$,
which is implied since $y_{v(i)} + y_{l(s)} - 1 \leq 0 \leq \Lambda$.
$\blacksquare$

The sampled-matching block below uses access in exactly those two ways ---
\eqref{eq:dw_f34} is (a), and \eqref{eq:dw_f6} rearranges to
$\pi_{i,s} - z_{i,s} + 1 \geq a_{i,s}$ with a non-negative left-hand side while
\eqref{eq:dw_s2} is (b) outright --- so all of its access variables can be
eliminated. The implementation nonetheless keeps them, because with a
constraint-programming backend the reification \eqref{eq:dw_pricing_access}
propagates in both directions where the substituted form does not, and the
number of distinct vertex pairs is small. Proposition 18 is what says those
variables are a convenience rather than a requirement.

The MID block is different, and the difference is instructive.
\eqref{eq:dw_effective} does not bound the effective rejected mass; it
\emph{selects} between two values, $\widehat t$ and $L$, and the selecting row
$\widehat e_{g,r} \geq L(1 - a_{g,s})$ becomes over-restrictive under the
substitution --- at $y_{v(g)} = y_{l(s)} = 0$ it would read
$\widehat e_{g,r} \geq 2L$. So the conjunction has to be represented there, and
\eqref{eq:dw_pricing_access} is not optional for the finite-grid block.

\noindent\textbf{The MID block.} With $a_{g,s}$ as above, impose
\eqref{eq:dw_threshold}--\eqref{eq:dw_mass} and, for each program,
\begin{equation}\label{eq:dw_mid_capacity}
\sum_{g}\sum_{r : s(g,r)=s} n_g \widehat d_{g,r} \;\leq\; L\, q_s\, y_{l(s)} .
\end{equation}
The right-hand side is the single-zone strengthening of the capacity row in
\eqref{eq:dw_zone_value}: a program outside the zone seats nobody, so it is
valid, and it is far stronger than $\le L q_s$ in the relaxation. The nonlinear
operators have exact bounded linearizations --- for $a, b \in [0,M]$,
$h = \min(a,b)$ is
\begin{equation}\label{eq:dw_min_linearization}
0\leq h,\quad h\leq a,\quad h\leq b,\quad
h\geq a-M\chi,\quad h\geq b-M(1-\chi)
\end{equation}
with $\chi$ binary, and $h = \max(a, 0)$ for $a \in [-M, M]$ is
\begin{equation}\label{eq:dw_max_linearization}
h\geq a,\quad h\geq0,\quad
h\leq a+M(1-\chi),\quad h\leq M\chi ,
\end{equation}
using $M = \overline p_s$ for the threshold and $M = L$ for the recurrence ---
and \eqref{eq:dw_effective} is equivalent to
\begin{equation}\label{eq:dw_effective_linearization}
\begin{aligned}
0\leq\widehat e_{g,r}&\leq L, &
\widehat e_{g,r}&\geq\widehat t_{s(g,r),\rho(g,r)},\\
\widehat e_{g,r}&\geq L(1-a_{g,s(g,r)}), &
\widehat e_{g,r}&\leq\widehat t_{s(g,r),\rho(g,r)}+L(1-a_{g,s(g,r)}).
\end{aligned}
\end{equation}
A constraint-programming backend needs none of \eqref{eq:dw_min_linearization}
or \eqref{eq:dw_max_linearization}: $\min$ and $\max$ are primitive
constraints there and \eqref{eq:dw_effective_linearization} is a conditional
equality. This is the reason the implementation prices with CP-SAT rather than a
mixed-integer solver --- the recurrence \eqref{eq:dw_recurrence} is the
expensive half of the model, roughly $24{,}000$ auxiliary binaries on
$\text{Block}_2$ under \eqref{eq:dw_min_linearization}, and it is also the
weakest part of that relaxation. One valid inequality repairs much of the
weakness and is worth stating: since
$\min(a,b) = a + b - \max(a,b) \geq a + b - L$ for $a, b \in [0, L]$,
\begin{equation}\label{eq:dw_min_valid}
\widehat R_{g,r} \;\geq\; \widehat R_{g,r-1} + \widehat e_{g,r} - L ,
\end{equation}
equivalently $\widehat d_{g,r} \le L - \widehat e_{g,r}$: a type cannot take
more mass at rank $r$ than its acceptance chance there. It is valid because
\eqref{eq:dw_threshold} clamps the thresholds to the lottery interval, and it
changes no integral point.

\noindent\textbf{The sampled-matching block.} Here the access polytope is
needed, because $A$ is a variable and the mechanism cannot be run. Let
$y_{i,s} \in \{0,1\}$ seat applicant $i$ at program $s$, let
$z_{i,s} \in \{0,1\}$ record that $i$ clears $s$'s realized cutoff, let
$\pi_{i,s} = \sum_{s' \succeq_i^\psi s} y_{i,s'}$ be the weakly-preferred
prefix, and let $\Gamma(s)$ be the drawn order restricted to the retained
applicants. The block is
\begin{align}
& \pi_{i,s} \in \{0,1\}, \qquad
\pi_{i,s} = \pi_{i,\mathrm{prev}(i,s)} + y_{i,s},
\quad \pi_{i,\mathrm{prev}(i, \text{first})} = 0 \label{eq:dw_f1} \\
& \sum_{i} y_{i,s} \leq q_s\, y_{l(s)} \label{eq:dw_f2} \\
& y_{i,s} \leq a_{i,s}, \qquad y_{i,s} \leq z_{i,s} \label{eq:dw_f34} \\
& z_{i,s} \leq z_{i',s} \qquad \text{for consecutive } (i', i) \text{ in } \Gamma(s) \label{eq:dw_f5} \\
& \pi_{i,s} \geq a_{i,s} + z_{i,s} - 1 \label{eq:dw_f6} \\
& \sum_{i'} y_{i',s} \geq q_s\,(1 - z_{i,s}) \label{eq:dw_s1} \\
& q_s \pi_{i,s} + \sum_{i' \succ_s^\psi i} y_{i',s} \geq q_s\, a_{i,s} \label{eq:dw_s2}
\end{align}
with the objective $\sum_{i,s} u_{i,s} y_{i,s}$. Row \eqref{eq:dw_f5} makes
$\{i : z_{i,s} = 1\}$ a prefix of $\Gamma(s)$, and a prefix of the realized
priority order \emph{is} a cutoff, so no cutoff variable is needed and nothing
has to live on a grid --- the contrast with \eqref{eq:dw_threshold} is the whole
difference between the two blocks. Declaring $\pi$ binary in \eqref{eq:dw_f1}
states ``at most one seat per applicant'' as a domain rather than a row.
Row \eqref{eq:dw_f6} is the no-blocking condition of Lemma 1 and
\eqref{eq:dw_s2} is the aggregated stability row \eqref{saa:rothblum}, which
implies it and tightens the relaxation; \eqref{eq:dw_s1} is non-wastefulness,
unnecessary for exactness but worth roughly two orders of magnitude in solve
time through presolve. Row \eqref{eq:dw_f2} is again the single-zone
strengthening. The inner sum of \eqref{eq:dw_s2} is carried by one running
prefix variable per pair, so the family is $O(|\Gamma(s)|)$ rows rather than
$O(|\Gamma(s)|^2)$.

\noindent\textbf{Proposition 19 (Exactness of the sampled-matching block).}
\textit{Fix integral $y_v$ with $A = \{v : y_v = 1\} \in \mathcal F_z$. Then the
integral points of \eqref{eq:dw_f1}--\eqref{eq:dw_s2}, projected onto the
applicants of $A$ and the programs of $A$, are exactly the stable matchings of
$\Gamma(a_A)$, and the maximum of $\sum u_{i,s} y_{i,s}$ over them is
$W^{\psi}(A)$. Retained applicants may be restricted to those with
$v(i) \in A'$ for any $A' \supseteq A$ without changing this.}

\textit{Proof.} At integral $y_v$ the reification
\eqref{eq:dw_pricing_access} pins $a_{i,s} = \mathds{1}\{(i,s) \in \Gamma(a_A)\}$.
Rows \eqref{eq:dw_f1}, \eqref{eq:dw_f2} and the first half of
\eqref{eq:dw_f34} then make $y$ a matching of $\Gamma(a_A)$ within capacity.

\emph{Feasible implies stable.} Suppose $(i,s) \in \Gamma(a_A)$ blocked $y$.
By Lemma 1, $D(T^\psi(i,s)) = 0$, i.e.\ $\pi_{i,s} = 0$, and
$D(S^\psi(s,i)) \leq q_s - 1$. Since $\pi_{i,s} = 0$ we have $y_{i,s} = 0$, so
$\sum_{i' \succ_s^\psi i} y_{i',s} = D(S^\psi(s,i)) \leq q_s - 1$, and
\eqref{eq:dw_s2} at $a_{i,s} = 1$ demands $q_s$ --- a contradiction. (Without
\eqref{eq:dw_s2} the same conclusion follows from \eqref{eq:dw_f34}--\eqref{eq:dw_s1}:
$\pi_{i,s} = 0$ forces $z_{i,s} = 0$ through \eqref{eq:dw_f6}, then
\eqref{eq:dw_s1} makes $s$ full, and the prefix structure \eqref{eq:dw_f5}
together with \eqref{eq:dw_f34} puts all $q_s$ admits weakly above $i$,
contradicting $D(S^\psi(s,i)) \leq q_s - 1$.)

\emph{Stable implies feasible.} Let $D$ be a stable matching of $\Gamma(a_A)$
and set $z_{i,s} = 1$ exactly when $i$ is weakly above $s$'s realized cutoff,
meaning $s$'s worst admit if $s$ is full and the bottom of $\Gamma(s)$
otherwise. Then \eqref{eq:dw_f5} holds by construction and \eqref{eq:dw_f34}
because an admit is weakly above the worst admit. For \eqref{eq:dw_f6}, suppose
$a_{i,s} = z_{i,s} = 1$ and $\pi_{i,s} = 0$: then $i$ has access, receives
nothing weakly preferred to $s$, and either $s$ has a free seat or $s$ admitted
someone no better than $i$ --- so $(i,s)$ blocks, contradiction. For
\eqref{eq:dw_s1}, $z_{i,s} = 0$ places $i$ strictly below the cutoff, which by
construction happens only when $s$ is full. For \eqref{eq:dw_s2}, if
$a_{i,s} = 1$ and $\pi_{i,s} = 0$ then non-blocking forces $s$ full with every
admit strictly preferred to $i$, so the inner sum is $q_s$.

The claimed maximum is then \eqref{eq:dw_sampled_value}. For the last sentence,
an applicant with $v(i) \notin A$ has $a_{i,s} = 0$ for every $s$, so
\eqref{eq:dw_f34} forces $y_{i,\cdot} = 0$ and \eqref{eq:dw_f6} and
\eqref{eq:dw_s2} are vacuous; their $z_{i,s}$ is unconstrained except through
\eqref{eq:dw_f5}, where an absent variable and a free variable impose the same
thing --- nothing. $\blacksquare$

The joint pricing problem is therefore \eqref{eq:dw_pricing} maximized over the
memberships together with the chosen welfare block:
\begin{equation}\label{eq:dw_joint_pricing_obj}
\max \quad W\big(\{v : y_v = 1\}\big)
-\sum_v\alpha_v y_v-\gamma_z
-\frac{\mu}{2}\sum_{\{u,v\}\in E}w_{u,v}\beta_{uv} ,
\end{equation}
exact by Proposition 7 for $W^{\mathrm{MID}}_L$ and by Proposition 19 for
$W^{\psi}$. It requires no simultaneous partition of the complement, and no
variable of any other label.

\noindent\textbf{Integer scaling.}
A constraint-programming backend needs integer objective coefficients, so
\eqref{eq:dw_joint_pricing_obj} is multiplied by a scale $K$ and rounded. The
rounding must not destroy the bound of Section~\ref{app:dw_bounds}, and the
direction is what secures it.

\noindent\textbf{Proposition 20 (Directional rounding gives a valid bound with a
computable allowance).} \textit{Let $\Upsilon \geq 1$, let $K = L\Upsilon$ for
$W^{\mathrm{MID}}_L$ and $K = \Upsilon$ for $W^{\psi}$, and form the integer
objective $\Theta$ from \eqref{eq:dw_joint_pricing_obj} by replacing each
utility coefficient $u$ by $\lceil \Upsilon u \rceil$, each dual $\alpha_v$ by
$\lfloor K\alpha_v \rfloor$, and each boundary coefficient
$\mu w_{u,v}/2$ by $\lfloor K \mu w_{u,v}/2 \rfloor$, dropping the constant
$\gamma_z$. Then for every $A \in \mathcal F_z$}
\begin{equation}\label{eq:dw_rounding}
f_z(A) \;\leq\; \frac{\Theta(A)}{K} - \gamma_z \;\leq\; f_z(A) + \varepsilon_z,
\qquad
\varepsilon_z \;=\; \frac{m_z + |V_z| + |E_z|}{K},
\end{equation}
\textit{where $V_z$ and $E_z$ are the membership and boundary variables of label
$z$ and $m_z = L|\mathcal G_{z}|$ for $W^{\mathrm{MID}}_L$,
$m_z = |\mathcal I_{z}|$ for $W^{\psi}$, and $m_z = 0$ for $W^{\partial}$.
Consequently any upper bound on $\max_A \Theta(A)/K - \gamma_z$ is an upper
bound on $\max_A f_z(A)$, and if that bound does not exceed $\varepsilon_z$ then
no $A \in \mathcal F_z$ has $f_z(A) > \varepsilon_z$.}

\textit{Proof.} Each of the three substitutions moves the objective weakly up:
$\lceil \Upsilon u \rceil \geq \Upsilon u$ multiplies a non-negative mass, and
$\lfloor K \alpha_v \rfloor \le K\alpha_v$ and
$\lfloor K\mu w/2 \rfloor \le K\mu w/2$ are \emph{subtracted}. This gives the
left inequality. For the right, each substitution moves it up by less than one
scaled unit per term it multiplies. The dual terms contribute at most $|V_z|$
and the boundary terms at most $|E_z|$, one per variable. For the utility
terms, the masses of a single MID type sum to at most $L$ by
\eqref{eq:dw_recurrence}, so that type contributes under $L$; an applicant
holds at most one seat by \eqref{eq:dw_f1}, so that applicant contributes under
$1$; and $W^{\partial}$'s coefficients $K w_{u,v}/2$ are integers already.
Dividing by $K$ gives \eqref{eq:dw_rounding}. The two consequences are
immediate. $\blacksquare$

The allowance $\varepsilon_z$ is reported, not subtracted: subtracting it would
break the left inequality of \eqref{eq:dw_rounding} and with it the bound's
validity, whereas adding $\sum_z \varepsilon_z$ to the incumbent comparison
turns the search's guarantee into a stated absolute tolerance. On
$\text{Block}_2$ at $\Upsilon = 10^3$ and $L = 20$ it is a fraction of a unit
per label against a welfare near $15{,}000$, and it shrinks like $1/\Upsilon$.
Nothing else about the rounding is trusted: a zone returned by pricing is
re-tested against \eqref{eq:dw_anchor}--\eqref{eq:dw_local}, re-scored by the
exact evaluator of Section~\ref{app:dw_value}, and its reduced cost recomputed
by the master, so a rounding artefact can waste a round but cannot admit a
column that does not improve.

\subsection{Bounds and the Choice of Dual Point}
\label{app:dw_bounds}

\noindent\textbf{Proposition 9 (Pricing-Corrected Upper Bound).} \textit{Let
$(\alpha, \gamma, \mu)$ be \emph{any} vector with $\mu \geq 0$, and suppose
$\overline r_z$ upper-bounds $\max_{A \in \mathcal F_z} f_z(A)$ at that vector
for each label, where $f_z$ is \eqref{eq:dw_pricing}. Then}
\begin{equation}\label{eq:dw_corrected_bound}
\mathrm{UB}=\sum_v\alpha_v+\sum_z\gamma_z+\mu\overline b
              +\sum_z\max\{0,\overline r_z\}
\end{equation}
\textit{is an upper bound on the complete LP relaxation of
\eqref{eq:dw_master_obj}--\eqref{eq:dw_master_boundary}, and hence on every
admissible partition at that branch.}

\textit{Proof.} Set $\gamma'_z = \gamma_z + \max\{0, \overline r_z\}$. For every
$A \in \mathcal F_z$ the reduced-cost bound gives
$W(A) - \sum_{v \in A}\alpha_v - \gamma_z - \mu b(A) \leq \overline r_z \leq
\max\{0, \overline r_z\}$, i.e.\
$\sum_{v\in A}\alpha_v+\gamma'_z+\mu b(A)\geq W(A)$, which is
\eqref{eq:dw_dual_row}. With $\mu \geq 0$ the triple
$(\alpha, \gamma', \mu)$ is therefore feasible for the complete dual, and its
objective is \eqref{eq:dw_corrected_bound}. Weak duality proves the claim.
$\blacksquare$

The statement deliberately does not require $(\alpha, \gamma, \mu)$ to be an
optimal dual solution of the restricted master, or a dual solution at all
beyond $\mu \geq 0$. That freedom is what the implementation spends, and the
reason is worth recording because it was the binding constraint on this method.

The master is a set-partitioning LP and is therefore massively degenerate. An
earlier build solved it with dual simplex, pinned so that a given column pool
always produced the same duals, and the consequence was measurable: on
$\text{BlockGroup}_0$ with a $404$-column pool, three to six columns of
strictly positive reduced cost entered on every round and the LP value stayed at
$11{,}142.80$ for eight consecutive rounds. A positive reduced cost that does
not move the objective is a degenerate pivot --- the entering column's step
length is zero. Only $132$ of $579$ cover rows carried a nonzero dual, so
\eqref{eq:dw_pricing} was close to maximizing raw welfare, and six labels priced
independently under an unanchored family returned six zones of $340$ to $392$
vertices out of $579$, no six of which can tile anything.

Two standard remedies apply, and Proposition 9 is what makes both free.
\emph{Interior duals}: solve the master by a barrier method and stop at the
barrier iterate rather than crossing over to a vertex, which spreads the prices
across the cover rows instead of concentrating them on a basis.
\emph{Dual smoothing}: price at
$\widetilde\alpha^{(t)} = a\,\alpha^{(t)} + (1-a)\,\widetilde\alpha^{(t-1)}$ for
$a \in (0,1]$, the Wentges rule, and likewise for $\gamma$ and $\mu$; a convex
combination of non-negative $\mu$'s is non-negative, so Proposition 9 still
applies, and the bound is simply evaluated at the smoothed point. A third
remedy is structural and is Section~\ref{app:dw_families}'s: an anchored family
cannot return two thirds of the district, so a priced zone is at least the right
size to be part of a tiling.

\noindent\textbf{Where the method actually stalls.}
None of the three is sufficient, and the reason is worth recording because it is
not the dual degeneracy those remedies were chosen for. On a district-sized
instance ($|V| = 576$, six anchors, $4{,}032$ applicants, $|\Gamma| = 12{,}096$)
seeded with one admissible partition, all four combinations of
$\{\text{vertex},\ \text{interior}\} \times \{a = 1,\ a = 0.5\}$ admitted
$400$ to $424$ columns of strictly positive reduced cost in ten rounds and left
the restricted LP at the seed value $13{,}161.9302$ to four decimal places.
Deleting the $k$ seed columns from the resulting $413$-column pool and re-solving
the integer master returns \textsc{infeasible}: the pool contains exactly one
tiling.

The obstruction is a dimension count rather than a choice of basis. The
restricted LP carries $|V| + k$ equality rows --- $582$ here --- against the
number of columns generated, so while the pool is small relative to $|V|$ the
system is over-determined and its solution set is essentially the tilings it
happens to contain. Pricing each label independently maximizes
\eqref{eq:dw_pricing} with no reference to the other labels, so $k$ priced zones
are under no pressure to be mutually complementary, and a pool can grow by
hundreds of individually improving zones without acquiring a second tiling.
What a pool generated this way delivers is recombination of \emph{seeded}
partitions through a set-partitioning master, which is a useful primal heuristic
and not column generation. Note that this is a statement about the pool and not
about the cost of pricing, which anchoring and Section~\ref{app:dw_pricing}'s
constraint-programming formulation reduced by roughly two orders of magnitude.

The diagnosis names its own repair: what is needed is a source of columns that
are mutually \emph{completable} by construction, and pricing cannot be it,
because the per-label question is exactly the question Proposition 9 requires
and it makes no reference to the other labels. Section~\ref{app:dw_redraw}
supplies such a source --- optimal re-partitions of two zones of an incumbent
partition, every feasible solution of which is a complete tiling --- and
Proposition 25 there records why adding it preserves both Proposition 9 and
Theorem 3 while replacing pricing with it would destroy both.

Two further repairs suggest themselves, neither implemented here. The first is
an \emph{elastic} master: replace \eqref{eq:dw_master_cover} by
$\sum \lambda + d_v - e_v = 1$ with $d, e \geq 0$ penalized at a large $M$,
which is Phase I extended to Phase II. That makes the LP full-dimensional, so a
column can enter with a positive step length and the duals become informative;
it remains a relaxation of the exact master, so Proposition 9 survives, and a
zero-slack optimum is still a partition. The second is pricing that knows its
columns must tile: a cardinality or student-mass window per label, read off the
incumbent, imposed on the column-generating pass while the bound-producing pass
still ranges over all of $\mathcal F_z$. Until one of them is in place the
\emph{bound} reported by this appendix's algorithm does not improve on the
a-priori constants of Appendix~\ref{app:welfare_bounds}, whatever the
incumbent does.

One caveat on smoothing is load-bearing. If pricing at the smoothed point finds
no column improving at the \emph{true} duals --- a \emph{mis-priced} round --- it
must be repeated at $\alpha^{(t)}$ before any conclusion is drawn, and the
smoothing centre reset there. Otherwise a smoothed point reports a convergence
the restricted LP does not have, and \eqref{eq:dw_corrected_bound} evaluated at
it is needlessly weak. Note also that Proposition 9 remains valid in Phase I
below: a deficit artificial has coefficient $+1$ in its own row, so raising a
convexity dual preserves its dual inequality, and a convex combination of two
Phase-I dual points still satisfies those inequalities --- which is why the
smoothing centre is reset when the phase changes.

\noindent\textbf{Phase I.} Seeding is not required for correctness. If the
restricted pool cannot cover the graph, add non-negative deficit-only
artificials to \eqref{eq:dw_master_cover} and \eqref{eq:dw_master_label} and
maximize minus their sum. An empty pool is Phase-I feasible, objective zero is
equivalent to feasibility of the original master, and Phase-I pricing uses zero
column cost with the same geometry and branch restrictions. Artificials never
appear in a returned partition, and a strictly negative certified Phase-I value
of \eqref{eq:dw_corrected_bound} proves the branch infeasible. Time-limited or
failed pricing is never read as proof that no improving column exists.

\noindent\textbf{Seeding.} Two sources supply initial columns, and they do
different jobs. A single feasibility solve of the base model
$\Omega(G, C, k, \delta)$ produces an admissible partition, hence $k$ columns
that make the master feasible immediately. ReCom then samples from that
partition and each sampled zone is \emph{rejection-filtered} against
$\mathcal F_z$. The filter is per zone rather than per partition, which is what
makes it worth running: a recombination step rewrites two zones and leaves the
rest of an admissible partition alone, and ReCom's own cut candidates already
respect $\mathcal C(\cdot)$, hence centroid anchoring, so what its zones
typically fail is \eqref{eq:dw_support_row}. The filter is a set-membership test
and runs before the welfare evaluator, so a rejected sample costs almost
nothing. Neither source needs coverage of the feasible partition space or a
mixing assumption; both are finite and bounded independently of the exact
search.

\subsection{Redrawing Two Zones at Once}
\label{app:dw_redraw}
The diagnosis at the end of Section~\ref{app:dw_bounds} identifies the pool,
not the prices, as the obstruction, and it also says what a repair has to do:
produce columns that are \emph{mutually completable} rather than individually
attractive. This section gives one, and states precisely why it can be run
alongside the pricing of Section~\ref{app:dw_pricing} without disturbing
Proposition 9 or Theorem 3 --- and why it must not be run instead of it.

Fix an admissible partition $\mathcal A = (A_0, \dots, A_{k-1})$, so
$A_z \in \mathcal F_z$ for every $z$ and $\bigsqcup_z A_z = V$. Fix two labels
$a \neq b$, write $U = A_a \cup A_b$ for their joint \emph{territory}, and let
\begin{equation}\label{eq:dw_pair_splits}
\mathcal S_{ab}(\mathcal A) \;=\;
\big\{ (B_a, B_b) \;:\; B_a \in \mathcal F_a,\ B_b \in \mathcal F_b,\
B_a \sqcup B_b = U \big\} .
\end{equation}
The incumbent pair $(A_a, A_b)$ belongs to $\mathcal S_{ab}(\mathcal A)$, so the
set is never empty. Write $\delta(A)$ for the edges with exactly one endpoint in
$A$, $E_U(B_a, B_b)$ for the edges of $G[U]$ with one endpoint in each part, and
$w(\cdot)$ for total edge weight, so that $b(A) = \tfrac12 w(\delta(A))$.

\noindent\textbf{Proposition 21 (Every pair split is a tiling).} \textit{Let
$(B_a, B_b) \in \mathcal S_{ab}(\mathcal A)$ and let $\mathcal A'$ be obtained
from $\mathcal A$ by replacing $A_a$ by $B_a$ and $A_b$ by $B_b$. Then
$\mathcal A'$ is an admissible partition of $V$, its objective exceeds
$\mathcal A$'s by exactly
$[W(B_a) + W(B_b)] - [W(A_a) + W(A_b)]$, and}
\begin{equation}\label{eq:dw_pair_boundary}
b(B_a) + b(B_b) \;=\; \tfrac12\, w(\delta(U)) \;+\; w\big(E_U(B_a, B_b)\big) .
\end{equation}
\textit{Consequently the district cap \eqref{eq:dw_master_boundary} holds at
$\mathcal A'$ if and only if}
\begin{equation}\label{eq:dw_pair_cap}
b(B_a) + b(B_b) \;\leq\; \overline b - \sum_{z \notin \{a, b\}} b(A_z) ,
\end{equation}
\textit{an inequality in the pair alone.}

\textit{Proof.} $B_a \sqcup B_b = U = A_a \sqcup A_b$, so $\mathcal A'$ covers
$V$ exactly once, and $B_a \in \mathcal F_a$, $B_b \in \mathcal F_b$ by
\eqref{eq:dw_pair_splits} while the remaining $A_z$ are unchanged; Proposition
15 then makes $\mathcal A'$ admissible. By Proposition 17 the objective is
$\sum_z W(\cdot)$ over the parts, and the $k - 2$ untouched parts contribute the
same value to both, which gives the displayed difference. For
\eqref{eq:dw_pair_boundary}, an edge of $\delta(B_a)$ either has its other
endpoint in $B_b$, hence lies in $E_U(B_a, B_b)$, or leaves $U$; the same holds
for $\delta(B_b)$; each edge of $E_U(B_a, B_b)$ is therefore counted once in
each of the two perimeters, while each edge of $\delta(U)$ is counted once in
total, because its endpoint in $U$ lies in exactly one part. Halving gives
\eqref{eq:dw_pair_boundary}. Finally $\sum_z b$ at $\mathcal A'$ equals
$\sum_{z \notin \{a,b\}} b(A_z) + b(B_a) + b(B_b)$, so
\eqref{eq:dw_master_boundary} is \eqref{eq:dw_pair_cap}. $\blacksquare$

Proposition 21 is the whole point. \emph{Every} member of
$\mathcal S_{ab}(\mathcal A)$ --- not only the optimum, and not only the
improving ones --- yields two columns that provably complete the $k-2$ columns
already in the pool into a tiling. A pool seeded with one partition and swept
over $m$ pair splits therefore admits at least $m+1$ tilings, which is exactly
the dimension the restricted master was missing. It also shows that the boundary
budget need not be priced here: the untouched perimeters and the part of the
pair's perimeter that faces them are constants, so the cap is an exact linear
row in the pair's own cut indicators.

\noindent\textbf{Proposition 22 (The cover prices cancel).} \textit{Fix any
$(\alpha, \gamma, \mu)$ with $\mu \geq 0$ and let $f_z$ be
\eqref{eq:dw_pricing}. For every $(B_a, B_b) \in \mathcal S_{ab}(\mathcal A)$,}
\begin{equation}\label{eq:dw_pair_reduced}
f_a(B_a) + f_b(B_b) \;=\;
\Big[ W(B_a) + W(B_b) - \mu\, w\big(E_U(B_a,B_b)\big) \Big]
\;-\; \kappa ,
\end{equation}
\textit{where $\kappa = \sum_{v \in U}\alpha_v + \gamma_a + \gamma_b
+ \tfrac{\mu}{2} w(\delta(U))$ does not depend on $(B_a, B_b)$. In particular,
if $\mu = 0$ then maximizing the pair's summed welfare over
$\mathcal S_{ab}(\mathcal A)$ maximizes its summed reduced cost, and a split
whose summed welfare strictly exceeds the incumbent pair's supplies a column of
strictly positive reduced cost.}

\textit{Proof.} Summing \eqref{eq:dw_pricing} over the two labels,
$\sum_{v \in B_a}\alpha_v + \sum_{v \in B_b}\alpha_v = \sum_{v \in U}\alpha_v$
because $B_a \sqcup B_b = U$, and $\gamma_a + \gamma_b$ is constant; the
$\mu$ term is $-\mu[b(B_a) + b(B_b)]$, which \eqref{eq:dw_pair_boundary}
rewrites as $-\tfrac{\mu}{2} w(\delta(U)) - \mu\, w(E_U(B_a,B_b))$. Collecting
the constants gives \eqref{eq:dw_pair_reduced}. For the last claim, if
$f_a(B_a) + f_b(B_b) > f_a(A_a) + f_b(A_b)$ and the incumbent columns are in the
pool with non-positive reduced cost, at least one of the two new columns has
$f > 0$. $\blacksquare$

So the redraw needs no dual solution, no restricted master and no LP: with no
boundary row it maximizes the summed reduced cost at \emph{every} dual point
simultaneously, and with one it solves the primal-exact problem, maximizing
summed welfare subject to \eqref{eq:dw_pair_cap}, which is tighter than pricing
$\mu$ and keeps every returned split admissible for
\eqref{eq:dw_master_boundary}. It is therefore both an exact local search on the
map and a reduced-cost-maximizing column generator, and it can be run before the
first master solve.

\noindent\textbf{Proposition 23 (Territory invariance).} \textit{If
$\mathcal A'$ arises from $\mathcal A$ by an $(a,b)$-redraw then
$\mathcal S_{ab}(\mathcal A') = \mathcal S_{ab}(\mathcal A)$. Hence a pair
solved to optimality remains optimal until some other pair alters $A_a$ or
$A_b$; and if no pair admits an improvement, $\mathcal A$ is a local maximum of
the objective over the two-zone neighbourhood
$\bigcup_{a \neq b}\mathcal S_{ab}(\mathcal A)$.}

\textit{Proof.} $\mathcal S_{ab}(\mathcal A)$ depends on $\mathcal A$ only
through $U = A_a \cup A_b$, and every element of it has $B_a \cup B_b = U$, so
the territory --- and with it the whole set --- is unchanged. $\blacksquare$

Proposition 23 is what makes a sweep cheap. Memoizing solved pairs on
$(a, b, U)$ needs no invalidation rule: an entry expires precisely when another
pair moves a vertex across $A_a$ or $A_b$, because that changes $U$.

\noindent\textbf{Proposition 24 (Non-adjacent pairs are inert).} \textit{If no
edge of $G$ joins $A_a$ to $A_b$ then
$\mathcal S_{ab}(\mathcal A) = \{(A_a, A_b)\}$.}

\textit{Proof.} Let $(B_a, B_b) \in \mathcal S_{ab}(\mathcal A)$ and suppose
$A_a \cap B_b \neq \emptyset$. Let $d(\cdot)$ denote distance to label $b$'s
school point and take $v \in A_a \cap B_b$ minimizing $d$. Since $c_b \in A_b$
and $A_a \cap A_b = \emptyset$ we have $v \neq c_b$, so
\eqref{eq:dw_support_row} supplies $u \in B_b$ adjacent to $v$ with
$d(u) < d(v)$. As $B_b \subseteq U$ and no edge joins $A_a$ to $A_b$, the
neighbour $u$ lies in $A_a$, hence in $A_a \cap B_b$ with $d(u) < d(v)$,
contradicting minimality. So $A_a \subseteq B_a$, and symmetrically
$A_b \subseteq B_b$; since both pairs partition $U$, equality holds.
$\blacksquare$

\noindent\textbf{The pair model.} Introduce
$x_{v, \zeta} = \mathds{1}\{v \in B_\zeta\}$ for $v \in U$ and
$\zeta \in \{a, b\}$ with $\zeta \in \mathcal C(v)$, and impose
\begin{equation}\label{eq:dw_pair_exactly_one}
\sum_{\zeta \in \{a,b\} \cap\, \mathcal C(v)} x_{v,\zeta} \;=\; 1
\qquad \forall v \in U ,
\end{equation}
together with, for each $\zeta$ separately,
\eqref{eq:dw_anchor}--\eqref{eq:dw_local} restricted to $U$ --- that is
\eqref{eq:dw_support_clause} with $\text{cn}^{*}(v, \zeta)$ intersected with
$U$, since $B_\zeta \subseteq U$ --- the cut indicators
\eqref{eq:dw_edge_linearization} when they are needed, the cap
\eqref{eq:dw_pair_cap}, and one welfare block per label, taken verbatim from
Section~\ref{app:dw_pricing}. The objective is
$\max\, W(B_a) + W(B_b)$. Four remarks on structure:

\begin{enumerate}
\item Row \eqref{eq:dw_pair_exactly_one} is the \emph{only} coupling between the
two labels, and it is what turns two zone problems into a re-partition. Where
both labels are candidates it has two terms, so $x_{v,b} = 1 - x_{v,a}$ and one
variable per vertex is eliminated at presolve: the solver sees the same
one-Boolean-per-vertex zoning model as \eqref{eq:dw_support_clause} does at a
single label. Where only one label is a candidate it fixes the vertex.
\item By Proposition 21 the model is always feasible --- the incumbent split
satisfies it --- so it is supplied as a warm start and an interrupted solve
still returns a tiling.
\item The welfare blocks are per label in both settings, so their structure is
unchanged. In particular the sampled-matching block
\eqref{eq:dw_f1}--\eqref{eq:dw_s2} remains Boolean apart from the per-program
seat counts and the running \eqref{eq:dw_s2} prefixes, and Proposition 18's
elimination still applies to it label by label; the finite-grid block retains
its small-domain cutoffs, thresholds and mass recurrence and, as before, needs
the access conjunction.
\item The one structural cost of two labels is that access
\eqref{eq:dw_pricing_access} is reified once per label rather than once, which
is the midpoint between the single-zone form and the $k$-joint form of
\eqref{saa:master_access}. Against it, $U$ is two zones' worth of vertices
rather than a label's whole candidate set.
\end{enumerate}

\noindent\textbf{Proposition 25 (The redraw is additive, and cannot replace
pricing).} \textit{(i) Enlarging the restricted pool by any columns whose
memberships lie in their own families leaves \eqref{eq:dw_corrected_bound} valid
and every step of Theorem 3 intact. (ii) Let
$\mathcal R_a = \{B_a : (B_a, B_b) \in \mathcal S_{ab}(\mathcal A)\}$. Then
$\mathcal R_a \subseteq \mathcal F_a$, and the inclusion is proper whenever some
member of $\mathcal F_a$ meets $V \setminus U$. Hence
$\max_{A \in \mathcal R_a} f_a(A)$ is in general a strict \emph{lower} bound on
$\max_{A \in \mathcal F_a} f_a(A)$ and does not satisfy the hypothesis of
Proposition 9.}

\textit{Proof.} (i) Proposition 9 is a statement about the \emph{complete} dual
of \eqref{eq:dw_master_obj}--\eqref{eq:dw_master_boundary} and its hypothesis
mentions no pool, so it is insensitive to which columns have been generated;
enlarging the pool enlarges the restricted primal's feasible set, so its LP
value can only move toward the complete LP value. Step 1 of Theorem 3 requires
only that a round which adds columns adds \emph{new} members of a finite set,
which deduplication and membership in $\mathcal F_z$ guarantee irrespective of
the source, and steps 2--4 read only the branch structure and the certified
bound. (ii) $\mathcal R_a \subseteq \mathcal F_a$ is immediate from
\eqref{eq:dw_pair_splits}. Every $A \in \mathcal R_a$ satisfies $A \subseteq U$,
so any $A \in \mathcal F_a$ with $A \not\subseteq U$ witnesses properness, and
the maximum over a subset cannot exceed the maximum over the whole.
$\blacksquare$

Part (ii) is the reason the redraw is an \emph{addition} to the algorithm rather
than a replacement for pricing. A node is closed only when every label's
pricing problem is solved globally over all of $\mathcal F_z$ and returns
nothing improving; a redraw contributes columns and incumbents and never a node
bound, so a round in which only the redraw fires re-solves the master and
returns to pricing. Were the redraw ever substituted for pricing, both
Proposition 9's bound and Theorem 3's guarantee would be lost, and what remained
would be a good heuristic with no certificate.

\noindent\textbf{Measured on $\text{Block}_2$.} Take the district instance with
$\texttt{max\_distance} = 3.1$, $k = 6$ centroids from the \texttt{6-zone-9}
configuration, $\texttt{program\_population} = \texttt{All}$ and citywide
programs removed as this appendix requires: $|V| = 501$ geographic units,
$1{,}340$ edges and $4{,}154$ applicants, giving $2{,}309$ compressed MID types
over $105$ programs and $4{,}154$ individual applicants over $100$ programs.
Anchoring leaves candidate sets of $147, 147, 253, 309, 248$ and $170$ units,
and $10$ of the $15$ label pairs are graph-adjacent at the seed, so
Proposition 24 discards a third of the neighbourhood before any model is built.
The seed is a single \texttt{cp\_bool} feasibility solve --- the
\emph{feasibility} hint, which optimizes nothing --- and its zones are
correspondingly lopsided, at $62, 65, 67, 81, 104$ and $122$ units. ReCom
sampling is switched off so that the comparison isolates the construction of
this section. Six workers, $600\,\mathrm{s}$.

\begin{center}
\begin{tabular}{lrrrrr}
\hline
& hint & no redraw & redraw & $\overline W^{\mathrm{TP}}$ & first-choice \\
\hline
$W^{\text{MID}}$ & $10{,}307.27$ & $10{,}307.27$ & $11{,}822.46$ & $15{,}956.50$ & $16{,}830.95$ \\
$W^{\text{SM}}$  & $10{,}344.55$ & $10{,}344.55$ & $12{,}105.20$ & $15{,}180.98$ & $15{,}983.65$ \\
\hline
\end{tabular}
\end{center}

The redraw is worth $+14.70\%$ under $W^{\text{MID}}$ and $+17.02\%$ under
$W^{\text{SM}}$ over the seeded partition. Without it the improvement is
$+0.000\%$, and that figure is exact rather than rounded: both objectives ran
$21$ column-generation rounds and generated $542$ to $555$ columns of strictly
positive reduced cost inside the budget, and the restricted LP sat at the hint
value to four decimal places on every single round, with the resulting pool
still admitting exactly one tiling. The diagnosis of
Section~\ref{app:dw_bounds} is therefore not an artefact of the synthetic
instance it was found on; it is what the method does on the district. With the
redraw enabled the same pool admits more than twenty-five tilings and the
restricted LP climbs through a sequence of distinct values ---
$10{,}307.27 \to 11{,}579.37 \to 11{,}624.64 \to 11{,}820.40 \to 11{,}822.46$
under $W^{\text{MID}}$ --- rather than stepping once. Extending the budget to
$1{,}800\,\mathrm{s}$ reaches $12{,}050.60$, or $+16.91\%$, over $107$ pair
solves and $28$ improvements, so the construction is budget-limited rather than
neighbourhood-limited at ten minutes. That partition was re-validated
independently of the search: it covers $V$ exactly once, each of its six zones
satisfies $\mathcal F_z$ on its own, and a fresh evaluation of the welfare
oracle on each zone reproduces the reported value to $10^{-6}$. Its zone sizes
are $65, 65, 77, 90, 94$ and $110$, against the hint's $62$ to $122$: the gain
is the rebalancing of an unbalanced feasibility solve, which is exactly what
program capacity rewards.

\noindent\textbf{Which side of the gap converges.} The two sides behave
completely differently and the distinction is the honest summary of this
construction. The \emph{primal} side converges, and on the district it is still
converging when the budget expires. The \emph{bound} side does not move at all:
\eqref{eq:dw_corrected_bound} evaluates to the first-choice constant
$\overline W^{\mathrm{FC}}$ of Appendix~\ref{app:welfare_bounds} in every arm,
with the redraw and without it, at every budget tried. The mechanism is visible
in the arithmetic. By LP duality the first two terms of
\eqref{eq:dw_corrected_bound} are the restricted master's own objective, so for
the bound to have stayed at $16{,}830.95$ while the restricted LP sat at
$12{,}050.60$, the pricing residuals must be summing to at least $4{,}780$ ---
close to $800$ per label. The entire headroom to the transportation bound is
$15{,}956.50 - 12{,}050.60 = 3{,}906$. Six labels are each separately reporting
more available improvement than provably exists in total, and no six of the
zones they return can occupy one tiling. That is not a defect in the pricing
problem, which answers the question Proposition 9 asks it and answers it to
proved optimality; it is the price of asking a per-label question. It is also
why the reported gap of roughly a quarter is mostly slack in the certificate:
measured against $\overline W^{\mathrm{TP}}$ the same incumbents are within
$24.5\%$ and $20.3\%$, and against the zoned transportation bound
$\widehat W^{\mathrm{ZT}}_{\mathrm{cn}}$ of Appendix~\ref{app:zoned_transport},
which on this instance at this $\texttt{max\_distance}$ cuts a further $1{,}134$
to $1{,}241$ units off $\overline W^{\mathrm{TP}}$, closer still.

This is Proposition 25(ii) seen from the other direction. The redraw supplies
the tilings the pool was missing, which is a statement about the primal, and by
construction it cannot drive the dual objective up, because its optimum ranges
over a subset of one label's family. What remains open is therefore exactly the
dual side, and it admits two quite different attacks: repair the decomposition's
own duals, via the elastic master or completability-targeted pricing described
in Section~\ref{app:dw_bounds}; or replace the constant that
\eqref{eq:dw_corrected_bound} falls back on with one of the zoning-aware bounds
this paper already develops, which requires nothing of the decomposition at
all. On a synthetic grid ($|V| = 324$, $k = 6$) with capacity tightened to
$80\%$ of the applicants needing seats the same comparison gives $+3.40\%$ and
$+0.58\%$ against $+0.000\%$; the district gains are larger because a
feasibility hint leaves far more on the table than the nearest-centroid tiling
that grid was seeded with. All figures are single runs and vary between them,
since a constraint solver using more than one search worker is
nondeterministic.

\subsection{Branching and Finite Convergence}
\label{app:dw_convergence}
Exact column generation solves the master LP, whose optimum may still be
fractional. Define the geographic assignment marginals
\begin{equation}\label{eq:dw_marginal}
\overline x_{v,z}=\sum_{\substack{A\in\mathcal F'_z\\v\in A}}\lambda_{z,A}.
\end{equation}
If $0<\overline x_{v,z}<1$, branch into $x_{v,z}=0$ and $x_{v,z}=1$. In the
first child, exclude $v$ from every column of label $z$; in the second, require
$v$ in every column of label $z$ and exclude it from every other label. Apply
each decision both to the stored pool and to all future pricing problems. These
operations define branch-specific subfamilies of $\mathcal F_z$ without adding
any cross-zone constraint to the pricing problem.

\begin{algorithm}[t]
\caption{Anchored Zone-Column Branch-and-Price}
\label{alg:dw_mid}
\begin{algorithmic}[1]
\Require Graph $G$, families $\mathcal F_z$, a zone value $W$, optional budget $\overline b$.
\State Seed the pool with one feasibility solve of $\Omega$, then with rejection-filtered ReCom samples; either may be empty.
\State Sweep two-zone redraws \eqref{eq:dw_pair_splits} from the seeded partition, retaining every split.
\State Retain the best sampled partition, if any; initialize the root search node.
\While{an unexplored node remains}
    \State Select a node; restrict pool and pricing to its branch decisions.
    \State Solve the restricted master; if artificials remain, stay in Phase I.
    \State Form the dual point: interior duals, smoothed against the previous centre.
    \State Price every label globally by \eqref{eq:dw_joint_pricing_obj}, in parallel, against one deadline.
    \If{no column improves at the smoothed point}
        \State Re-price at the raw duals and move the centre there.
    \EndIf
    \State Add every improving column.
    \If{the restricted master did not improve, or pricing added nothing}
        \State Sweep two-zone redraws from the incumbent; add every split; contribute no bound.
    \EndIf
    \State Repeat from the restricted master until neither source adds a column.
    \If{Phase I is certified strictly negative}
        \State Discard the node as infeasible and continue.
    \EndIf
    \State Optionally improve the incumbent using an integer restricted master.
    \If{the certified node bound does not exceed incumbent welfare $+\ \sum_z \varepsilon_z$}
        \State Discard the node.
    \ElsIf{all marginals \eqref{eq:dw_marginal} are integral}
        \State Recover the partition, update the incumbent, and discard the node.
    \Else
        \State Branch on a fractional marginal and retain both children.
    \EndIf
\EndWhile
\State \Return the incumbent, or infeasibility if no partition was found.
\end{algorithmic}
\end{algorithm}

\noindent\textbf{Theorem 3 (Finite Convergence).} \textit{Fix a finite graph,
the families $\mathcal F_z$, a zone value $W$ from
Section~\ref{app:dw_value} --- with a finite lottery scale $L$ in the MID case
and a fixed seed $\psi$ in the sampled case --- and a scale $\Upsilon$ with
allowances $\varepsilon_z$ as in Proposition 20. Suppose all restricted LPs and
pricing problems are solved globally in exact arithmetic and generated columns
are retained and deduplicated by label and membership. Then
Algorithm~\ref{alg:dw_mid} terminates after finitely many steps and returns
either an admissible partition whose value is within $\sum_z \varepsilon_z$ of
the optimum of \eqref{eq:dw_master_obj}--\eqref{eq:dw_master_binary}, or a proof
that no admissible partition exists. At $\Upsilon \to \infty$, or with rational
duals and exact utilities, $\varepsilon_z = 0$ and the returned partition is
globally optimal.}

\textit{Proof.} Each family contains at most $2^{|V|}-1$ memberships. At an
optimal restricted LP solution a stored column has non-positive reduced cost, so
every column added because its reduced cost at the raw duals is strictly
positive is new; a redraw round continues only when it contributes a column that
the pool did not already hold, which by \eqref{eq:dw_pair_splits} and
Proposition 25(i) is again a new member of its own family. Each node's
column-generation loop therefore adds at least one previously unseen member of
$\bigcup_z \mathcal F_z$ per iteration and so performs finitely many
iterations, whichever source supplies the column. When every label's pricing optimum is non-positive at the raw duals,
the restricted dual --- corrected as in Proposition 9 --- is feasible for the
complete branch-specific dual, so the node's LP bound is certified; by
Proposition 20 the certificate is reached once the scaled bound falls to
$\varepsilon_z$, which is why the closing test in
Algorithm~\ref{alg:dw_mid} carries $\sum_z \varepsilon_z$. A node whose
restricted master is still infeasible once its column family is exhausted
admits no admissible partition, every member of $\mathcal F_z$ compatible with
the branch having been generated by then. The pricing problems are bounded:
membership, cutoff, seat and clearing variables all have finite domains and the
boundary variables are binary.

Suppose all marginals \eqref{eq:dw_marginal} are integral. For each label the
non-negative column weights sum to one; a marginal of one forces every
positive-weight column of that label to contain the vertex and a marginal of
zero excludes it from all of them, so their memberships coincide and
deduplication leaves one column of weight one per label. By Proposition 15 that
is an admissible partition, and by Proposition 17 its value is
$\sum_z W(A_z)$.

Otherwise branching on a fractional marginal is disjoint and exhaustive for
integer partitions and fixes a previously unfixed binary assignment, which
pricing then respects, so tree depth is at most $k|V|$ and the tree is finite.
Infeasibility pruning removes no admissible partition and bound pruning is
justified by Proposition 9 up to the stated $\sum_z \varepsilon_z$. An integer
restricted master can improve the incumbent but is never needed to justify a
bound. On exhaustion no admissible partition exceeds the incumbent by more than
$\sum_z \varepsilon_z$; Proposition 8 identifies the incumbent's value with the
district objective. If no incumbent exists, every branch was proved infeasible.
$\blacksquare$

The guarantee is for the fixed graph, the fixed lottery scale or seed, and the
fixed scaling, with no polynomial bound on computation. If stopped early,
incumbent value is a lower bound and its maximum with the valid bounds of
unresolved nodes is a global upper bound; closing that gap to
$\sum_z \varepsilon_z$ certifies optimality at the stated tolerance. Seeding
contributes candidate zones, while pricing and branching establish the
guarantee.

Two limitations are worth stating plainly, because they bound what this appendix
claims. First, the decomposition is exact for the \emph{specified} value ---
finite-grid MID at scale $L$, or sampled stable matching at one seed $\psi$ ---
and not for continuous-lottery MID or for an expectation over seeds; a single
seed fixes the lottery, which is exactly what makes the column value an exactly
realized welfare rather than an average of them, and the variance across seeds
is not addressed here. Second, the implementation uses floating-point LP duals
and integer constraint-programming bounds, not formal exact-arithmetic
certificates, so the theorem's hypothesis of exact arithmetic is an idealization
of what is computed.


\end{document}
